% ============================================================
% C++ 核心指南详解 — 主文档
% 基于 ISO C++ Core Guidelines
% ============================================================

% !TEX program = xelatex
\documentclass[12pt,a4paper]{ctexbook}

% ==================== 页面 ====================
\usepackage[top=2.5cm,bottom=2.5cm,left=2.5cm,right=2.5cm]{geometry}

% ==================== 字体 ====================
\usepackage{fontspec}
\usepackage{iffont}
\IfFontExistsTF{Consolas}
  {\setmonofont{Consolas}}
  {\setmonofont{DejaVu Sans Mono}[Scale=MatchLowercase]}

% ==================== 颜色 ====================
\usepackage{xcolor}
\definecolor{codebg}{HTML}{F5F5F5}
\definecolor{codeframe}{HTML}{E0E0E0}
\definecolor{cppkeyword}{HTML}{1A1AFF}
\definecolor{cppcomment}{HTML}{2E8B2E}
\definecolor{cppstring}{HTML}{B22222}
\definecolor{linenumgray}{HTML}{999999}
\definecolor{algheadbg}{HTML}{1A3C5E}
\definecolor{csharpkeyword}{HTML}{8B008B}
\definecolor{csharptype}{HTML}{006400}

% ==================== listings ====================
\usepackage{listings}

% ---- listing style ----
\lstdefinestyle{cppstyle}{
  language=C++,
  basicstyle=\small\ttfamily,
  keywordstyle=\color{cppkeyword}\bfseries,
  commentstyle=\color{cppcomment}\itshape,
  stringstyle=\color{cppstring},
  numberstyle=\tiny\color{linenumgray},
  numbers=left,
  frame=single,
  rulecolor=\color{codeframe},
  framesep=6pt,
  breaklines=true,
  tabsize=4,
  columns=flexible,
  keepspaces=true,
  backgroundcolor=\color{codebg},
  showstringspaces=false,
  morekeywords={constexpr,consteval,constinit,decltype,auto,
    concept,requires,co_await,co_yield,co_return,
    noexcept,override,final,nullptr,static_assert,
    template,typename,namespace,using,mutable,volatile,
    explicit,operator,friend,inline,virtual,
    thread_local,alignas,alignof,if,consteval,constexpr},
  deletekeywords={int,char,bool,float,double,long,short,unsigned,signed,void,wchar_t},
  emph={size_t,int32_t,uint32_t,int64_t,uint64_t,ssize_t,
    string,vector,map,set,unique_ptr,shared_ptr,optional,variant,any,
    span,string_view,
    ifstream,ofstream,fstream,iostream,
    ostream,istream,stringstream,ostringstream,
    path,filesystem},
  emphstyle=\color{teal!80!black},
  escapeinside={(*@}{@*)},
}
\lstset{style=cppstyle}

% ==================== tcolorbox ====================
\usepackage[most]{tcolorbox}

\newtcolorbox{guideline}[1][]{
  enhanced,breakable,
  colback=blue!5!white,colframe=blue!75!black,
  fonttitle=\bfseries,
  title=要点,#1,
  boxed title style={colback=blue!75!black},
  attach boxed title to top left={yshift=-2mm,xshift=2mm},
  arc=2mm,boxrule=0.8mm,
}
\newtcolorbox{compare}[1][]{
  enhanced,breakable,
  colback=green!3!white,colframe=green!50!black,
  fonttitle=\bfseries,
  title=对比,#1,
  boxed title style={colback=green!50!black},
  attach boxed title to top left={yshift=-2mm,xshift=2mm},
  arc=2mm,boxrule=0.8mm,
}
\newtcolorbox{warning}[1][]{
  enhanced,breakable,
  colback=red!5!white,colframe=red!75!black,
  fonttitle=\bfseries,
  title=常见陷阱,#1,
  boxed title style={colback=red!75!black},
  attach boxed title to top left={yshift=-2mm,xshift=2mm},
  arc=2mm,boxrule=0.8mm,
}
\newtcolorbox{note}[1][]{
  enhanced,breakable,
  colback=orange!5!white,colframe=orange!80!black,
  fonttitle=\bfseries,
  title=注意,#1,
  boxed title style={colback=orange!80!black},
  attach boxed title to top left={yshift=-2mm,xshift=2mm},
  arc=2mm,boxrule=0.8mm,
}
\newtcolorbox{bestpractice}[1][]{
  enhanced,breakable,
  colback=purple!5!white,colframe=purple!60!black,
  fonttitle=\bfseries,
  title=最佳实践,#1,
  boxed title style={colback=purple!60!black},
  attach boxed title to top left={yshift=-2mm,xshift=2mm},
  arc=2mm,boxrule=0.8mm,
}

% ==================== 章节标题 ====================
\usepackage{titlesec}
\usepackage{fancyhdr}
\usepackage{graphicx}
\usepackage{amssymb}
\usepackage{amsmath}
\usepackage{tabularx}
\usepackage{booktabs}
\usepackage{longtable}
\usepackage{array}
\usepackage{multirow}
\usepackage{enumitem}
\usepackage{seqsplit}

\titleformat{\chapter}[display]
  {\normalfont\huge\bfseries\color{algheadbg}}
  {\chaptertitlename\ \thechapter}{20pt}{\Huge}
\titleformat{\section}
  {\normalfont\Large\bfseries\color{blue!70!black}}
  {\thesection}{1em}{}
\titleformat{\subsection}
  {\normalfont\large\bfseries\color{blue!60!black}}
  {\thesubsection}{1em}{}

\pagestyle{fancy}
\fancyhf{}
\fancyhead[LE,RO]{\thepage}
\fancyhead[RE]{\leftmark}
\fancyhead[LO]{\rightmark}

\setlist[itemize]{leftmargin=2em,itemsep=2pt}
\setlist[enumerate]{leftmargin=2em,itemsep=2pt}

% ==================== 超链接 ====================
\usepackage{hyperref}
\hypersetup{
  colorlinks=true,
  linkcolor=blue!70!black,
  urlcolor=blue!60!black,
  bookmarks=true,
  pdfauthor={C++ Core Guidelines Reference},
  pdftitle={C++ 核心指南详解},
}

\usepackage{makeidx}
\makeindex


% ==================== 自定义命令 ====================
\newcommand{\cpp}[1]{C++#1}
\newcommand{\Fn}[1]{\texttt{#1()}}
\newcommand{\Tp}[1]{\texttt{#1}}

\newcommand{\guideline}[1]{\textbf{\textcolor{blue!70!black}{#1}}}

% CoreBook 特有关键要点框
\newtcolorbox{keypoint}{
  enhanced,breakable,
  colback=blue!5!white,colframe=blue!75!black,
  fonttitle=\bfseries,title=要点,
  attach boxed title to top left={yshift=-2mm,xshift=2mm},
  boxed title style={colback=blue!75!black},
  arc=2mm,boxrule=0.8mm,
}

% ====================================================================
\begin{document}

% ------------------------------------------------------------
% 封面
% ------------------------------------------------------------
\begin{titlepage}
\centering

\vspace*{3cm}

{\Huge\bfseries\color{blue!60!black} C\texttt{++} 核心指南详解\par}

\vspace{1.2cm}

{\LARGE\color{blue!40!black} 基于 ISO C\texttt{++} Core Guidelines\par}

\vspace{2cm}

\tikzset{
  guidelineline/.style={
    draw=blue!60!black,
    line width=1.2pt
  }
}
\noindent\makebox[\textwidth]{\color{blue!60!black}\rule{0.6\textwidth}{1.2pt}}

\vspace{2cm}

{\large 面向现代 C\texttt{++} 开发者的权威参考\par}

\vspace{0.6cm}

{\large 涵盖类型安全、所有权、并发、错误处理等核心主题\par}

\vspace{3cm}

{\large\itshape \today\par}

\vfill

{\small\color{gray} 本书基于 Bjarne Stroustrup 与 Herb Sutter 维护的\\
C\texttt{++} Core Guidelines 社区项目编写\par}

\end{titlepage}

% ------------------------------------------------------------
% 前言
% ------------------------------------------------------------
\frontmatter

\chapter{前言}

C\texttt{++} 是一门历史悠久且持续演进的编程语言。自 2011 年 C\texttt{++}11 标准发布以来，语言本身经历了深刻的变革——移动语义、智能指针、Lambda 表达式、概念（Concepts）、协程（Coroutines）等现代特性的引入，使得 C\texttt{++} 的编程范式与最佳实践发生了根本性的变化。

然而，语言的演进并不意味着旧有的编程习惯会自动消失。大量项目仍然使用着过时的模式，许多开发者对新特性的理解也参差不齐。正是在这一背景下，Bjarne Stroustrup 与 Herb Sutter 于 2015 年发起了 C\texttt{++} Core Guidelines 项目，旨在为社区提供一套经过实践检验的、可被工具自动检查的编程规范。

本书是对 C\texttt{++} Core Guidelines 的系统性解读。我们不仅逐条阐释指南中的核心规则，还结合大量实际代码示例，展示如何在真实项目中应用这些规则。本书的目标是：

\begin{itemize}
  \item 帮助读者理解每条规则背后的设计哲学与技术原因；
  \item 通过正反两方面的代码示例，让读者直观感受规则的价值；
  \item 提供从旧式 C\texttt{++} 向现代 C\texttt{++} 渐进迁移的实用建议；
  \item 为团队制定编码规范提供参考框架。
\end{itemize}

\section*{如何使用本书}

本书按照主题分为多个部分，每部分包含若干章节。各部分之间相对独立，读者可以根据自己的需要选择阅读顺序。建议初学者从第一部分开始按序阅读；有经验的开发者可以直接跳到感兴趣的主题。

每条规则的讲解遵循统一的格式：首先给出规则的核心思想，然后通过代码示例展示正确与错误的做法，最后讨论规则的适用场景与例外情况。

\section*{约定}

\begin{itemize}
  \item \textbf{推荐做法}：用绿色边框的代码块表示。
  \item \textbf{错误做法}：用红色边框的代码块表示。
  \item 规则编号采用原指南的编号体系，如 \texttt{R.1}、\texttt{ES.5} 等。
  \item 书中代码默认使用 C\texttt{++}17 标准，除非另有说明。
\end{itemize}

% ------------------------------------------------------------
% 目录
% ------------------------------------------------------------
\tableofcontents

% ------------------------------------------------------------
% 正文
% ------------------------------------------------------------
\mainmatter

% ============================================================
% Part I: 基础与哲学 (Foundations & Philosophy)
% ============================================================

\part{基础与哲学}

% ------------------------------------------------------------
% Chapter 1: 引言
% ------------------------------------------------------------
\chapter{引言}

\section{C++ 核心指南的目标}

C++ 核心指南（C++ Core Guidelines）是由 Bjarne Stroustrup 和 Herb Sutter 于 2015 年发起的一项社区工程，旨在为现代 C++ 编程提供一套经过实践检验的编程规范。这套指南并非要取代 C++ 语言标准，而是在语言标准之上提炼出一组\textbf{高层原则}与\textbf{具体规则}，帮助程序员写出更加安全、高效、可维护的代码。

核心指南的设计目标可以归纳为以下几点：

\begin{enumerate}
  \item \textbf{帮助程序员拥抱现代 C++：}许多程序员对 C++ 的印象停留在 C with Classes 的阶段。核心指南通过具体示例展示 C++11、C++14、C++17 乃至 C++20 的新特性如何在实际项目中发挥作用。
  \item \textbf{提升代码的安全性：}所谓"安全"并非仅指内存安全，而是涵盖类型安全（type safety）、边界安全（bounds safety）和生命周期安全（lifetime safety）三个维度。
  \item \textbf{提供可执行的规范：}指南中的规则力求足够精确，使得静态分析工具能够自动检查大部分规则的遵守情况。
  \item \textbf{保持务实：}指南承认遗留代码的存在，提供了渐进式迁移的建议，而非要求一夜之间推倒重来。
\end{enumerate}

\section{目标读者}

核心指南面向所有使用 C++ 进行开发的程序员，无论你是：
\begin{itemize}
  \item 刚刚学习 C++ 的新手——指南可以帮助你从一开始就养成良好的编程习惯；
  \item 有多年经验的 C++ 老手——指南可以帮助你了解现代 C++ 的最佳实践；
  \item 系统架构师——指南中的接口设计原则和抽象策略对架构决策至关重要；
  \item 技术主管——指南可以作为团队编码规范的基础。
\end{itemize}

\section{安全剖面（Safety Profiles）}

核心指南定义了三个关键的安全剖面，它们分别关注程序安全性的不同方面：

\subsection{类型安全（Type Safety）}

类型安全意味着每个对象都以其被声明的类型来使用。类型安全的代码不会出现"将一个 \texttt{int} 当作指针来解引用"这样的错误。通过强类型接口、\texttt{enum class}、智能指针等手段，可以在编译期消除大量类型混淆的缺陷。

\begin{lstlisting}[language=C++]
// 类型不安全的示例
void process(int* p) {
    *p = 42;  // p 是否有效？p 是否指向一个 int？
}

void caller() {
    int x = 0;
    process(&x);  // OK
    process(nullptr);  // 未定义行为！
    int arr[10];
    process(arr);  // 可以编译，但语义可能不对
}
\end{lstlisting}

\begin{lstlisting}[language=C++]
// 类型安全的示例
void process(std::span<int> data) {
    if (!data.empty()) {
        data[0] = 42;  // 安全：有边界检查
    }
}

void caller() {
    std::vector<int> v = {1, 2, 3};
    process(v);  // OK，类型明确
}
\end{lstlisting}

\subsection{边界安全（Bounds Safety）}

边界安全要求所有对数组和容器的访问都在有效范围内进行。C 风格的数组退化是指针后，编译器无法追踪其大小，这导致越界访问成为 C/C++ 中最常见的安全漏洞之一。核心指南推荐使用 \texttt{std::array}、\texttt{std::vector}、\texttt{std::span} 等带有边界信息的类型。

\begin{lstlisting}[language=C++]
// 边界不安全的示例
void fill(int* arr, int size) {
    for (int i = 0; i <= size; ++i) {  // 错误：off-by-one
        arr[i] = 0;
    }
}
\end{lstlisting}

\begin{lstlisting}[language=C++]
// 边界安全的示例
void fill(std::span<int> arr) {
    for (auto& elem : arr) {  // 安全：范围 for 自动在边界内
        elem = 0;
    }
}
\end{lstlisting}

\subsection{生命周期安全（Lifetime Safety）}

生命周期安全关注的是对象在其被访问时确实存在。悬垂指针（dangling pointer）和使用已释放的内存是 C++ 中最难调试的问题之一。核心指南通过所有权语义（智能指针）、\texttt{not\_null}、\texttt{owner} 等机制来帮助程序员管理对象的生命周期。

\begin{lstlisting}[language=C++]
// 生命周期不安全的示例
int* dangerous() {
    int local = 42;
    return &local;  // 返回局部变量的地址——悬垂指针！
}

void also_dangerous() {
    int* p = new int(42);
    delete p;
    std::cout << *p << std::endl;  // 使用已释放的内存！
}
\end{lstlisting}

\begin{lstlisting}[language=C++]
// 生命周期安全的示例
std::unique_ptr<int> safe_factory() {
    return std::make_unique<int>(42);  // 所有权明确转移
}

void safe_usage() {
    auto ptr = safe_factory();
    std::cout << *ptr << std::endl;  // 安全：ptr 拥有有效的对象
}  // 自动释放
\end{lstlisting}

\section{如何使用本指南}

核心指南的规则按类别组织，每个类别用一个简短的前缀标识：

\begin{table}[h]
\centering
\begin{tabular}{ll}
\hline
\textbf{前缀} & \textbf{含义} \\
\hline
P & 哲学（Philosophy） \\
I & 接口（Interface） \\
F & 函数（Function） \\
C & 类与类层次（Class & Class Hierarchy） \\
Enum & 枚举（Enumeration） \\
Con & 常量与不可变性（Constancy） \\
T & 模板与泛型编程（Templates） \\
R & 资源管理（Resource Management） \\
ES & 表达式与语句（Expressions & Statements） \\
E & 错误处理（Error Handling） \\
Con & 并发（Concurrency） \\
SL & 标准库（Standard Library） \\
\hline
\end{tabular}
\end{table}

每条规则的编号格式为"类别.序号"，例如 \textbf{P.1} 表示哲学类的第一条规则。规则之间经常相互引用，形成一张有机的规则网络。

\section{关键概念与记号}

在阅读核心指南时，以下概念和记号反复出现，需要首先理解：

\subsection{不变量（Invariant）}

不变量是指在对象的生命周期中始终为真的条件。例如，一个表示日期的类，其月份字段应始终在 1 到 12 之间。不变量是类的契约的核心。

\subsection{所有权（Ownership）}

所有权决定了谁负责释放资源。核心指南强调所有权应该是明确的：要么通过值语义（对象自己管理自己的生命周期），要么通过智能指针（\texttt{unique\_ptr} 表示独占所有权，\texttt{shared\_ptr} 表示共享所有权）。

\subsection{前置条件与后置条件}

前置条件（precondition）是调用函数时必须满足的条件；后置条件（postcondition）是函数执行完毕后保证满足的条件。核心指南使用 \texttt{Expects()} 宏表示前置条件，\texttt{Ensures()} 宏表示后置条件。

\subsection{零规则、五规则}

零规则（Rule of Zero）建议类不应该自定义析构函数、拷贝/移动构造函数和拷贝/移动赋值运算符，而应该让编译器自动生成。如果类需要自定义其中任何一个，则通常需要自定义全部五个（五规则，Rule of Five）。

\section*{练习题}

\begin{enumerate}
  \item 解释类型安全、边界安全和生命周期安全之间的区别，并各举一个违反对应安全规则的代码示例。
  \item 核心指南的目标读者是谁？你认为它对于维护大型遗留 C++ 代码库有什么实际价值？
  \item 阅读以下代码，指出其中违反了哪些安全剖面，并给出修复方案：
\begin{lstlisting}[language=C++]
char* get_name() {
    char name[100];
    strcpy(name, "Hello World");
    return name;
}
\end{lstlisting}
  \item 为什么核心指南推荐使用 \texttt{std::span} 而非 \texttt{(pointer, size)} 对来传递数组？
  \item 用自己的话解释"不变量"的概念，并设计一个具有有意义不变量的 \texttt{Date} 类。
\end{enumerate}


% ------------------------------------------------------------
% Chapter 2: 编程哲学
% ------------------------------------------------------------
\chapter{编程哲学}

编程哲学是核心指南中最基础的部分。它确立了所有后续规则的指导思想。理解这些哲学原则，有助于在面对指南未明确覆盖的场景时做出正确的判断。

% === P.1 ===
\section{规则 P.1：直接在代码中表达思想}

\textbf{规则 P.1:} \textit{直接在代码中表达思想。}

\subsection{为什么重要}

代码是设计思想的最精确、最可靠的记录。注释会过时，文档会丢失，但代码始终与程序同步。如果一段逻辑无法直接用语言特性表达，说明可能需要寻找更好的抽象方式，或者需要自定义一个类型来封装这个概念。

\subsection{错误示例}

\begin{lstlisting}[language=C++]
// 用注释来表达"这个 int 表示米为单位的距离"
int d;  // 距离，单位：米
// ... 200 行代码之后 ...
d = 5;  // 这里的 5 是什么意思？需要回去找注释
\end{lstlisting}

\subsection{正确示例}

\begin{lstlisting}[language=C++]
// 用类型来表达
class Distance {
    double value_in_meters_;
public:
    constexpr explicit Distance(double m) : value_in_meters_(m) {}
    double in_meters() const { return value_in_meters_; }
};

Distance d{5.0};  // 清晰明确
\end{lstlisting}

\subsection{要点}
直接在代码中表达思想意味着使用类型系统、标准库算法和语言特性来传达设计意图，而不是依赖注释和约定。当你发现自己在写大段注释来解释"为什么"时，应该考虑是否可以用更好的代码结构来替代。

% === P.2 ===
\section{规则 P.2：使用 ISO 标准 C++ 编写}

\textbf{规则 P.2:} \textit{使用 ISO 标准 C++ 编写。}

\subsection{为什么重要}

依赖特定编译器扩展会使代码不可移植，阻碍不同平台和工具链之间的协作。标准 C++ 经过了精心设计，绝大多数需求都可以通过标准特性来满足。当确实需要平台特定功能时，应该将其封装在明确的接口之后。

\subsection{错误示例}

\begin{lstlisting}[language=C++]
// GCC 扩展：可变长数组（VLA）
void process(int n) {
    int arr[n];  // 非标准！某些编译器不支持
    // ...
}

// MSVC 扩展：__declspec
__declspec(dllexport) void my_function() { }
\end{lstlisting}

\subsection{正确示例}

\begin{lstlisting}[language=C++]
// 使用标准容器替代 VLA
void process(int n) {
    std::vector<int> arr(n);
    // ...
}

// 使用标准属性或条件编译
#ifdef _MSC_VER
  #define EXPORT __declspec(dllexport)
#else
  #define EXPORT __attribute__((visibility("default")))
#endif

EXPORT void my_function() { }
\end{lstlisting}

\subsection{要点}
坚持使用标准 C++ 可以确保代码的可移植性和长期可维护性。如果必须使用平台特定功能，应通过封装将其影响限制在最小范围内，并使用条件编译来保证在其他平台上的可编译性。

% === P.3 ===
\section{规则 P.3：表达意图}

\textbf{规则 P.3:} \textit{表达意图。}

\subsection{为什么重要}

代码的读者（包括未来的你自己）需要快速理解代码的\textbf{目的}，而不仅仅是它的\textbf{行为}。表达意图的代码是自解释的，它让读者无需深入实现细节就能理解高层逻辑。

\subsection{错误示例}

\begin{lstlisting}[language=C++]
// 读者需要仔细分析才能理解意图
for (int i = 0; i < v.size(); ++i) {
    if (v[i] == target) {
        result = i;
        break;
    }
}
\end{lstlisting}

\subsection{正确示例}

\begin{lstlisting}[language=C++]
// 意图一目了然
auto it = std::find(v.begin(), v.end(), target);
if (it != v.end()) {
    result = std::distance(v.begin(), it);
}
\end{lstlisting}

\subsection{要点}
使用标准库算法（如 \texttt{std::find}、\texttt{std::sort}、\texttt{std::accumulate}）不仅使代码更简洁，更重要的是它直接表达了"我在做什么"而非"我怎么做"。当标准库不能满足需求时，应该通过命名良好的函数和类型来传达意图。

% === P.4 ===
\section{规则 P.4：理想情况下程序应该是静态类型安全的}

\textbf{规则 P.4:} \textit{理想情况下，程序应该是静态类型安全的。}

\subsection{为什么重要}

静态类型安全意味着在编译时就能捕获类型错误。这比在运行时才发现错误要好得多——编译时错误零成本，而运行时错误可能导致崩溃、数据损坏甚至安全漏洞。

\subsection{错误示例}

\begin{lstlisting}[language=C++]
// 使用 union 进行类型双关——未定义行为
union Value {
    int i;
    float f;
};

Value v;
v.f = 3.14f;
int x = v.i;  // 未定义行为！读取不活跃的 union 成员
\end{lstlisting}

\subsection{正确示例}

\begin{lstlisting}[language=C++]
// 使用 std::variant 进行类型安全的联合
std::variant<int, float> v;
v = 3.14f;
// int x = std::get<int>(v);  // 运行时异常（类型不匹配）
if (auto pval = std::get_if<float>(&v)) {
    float fval = *pval;  // 类型安全
}
\end{lstlisting}

\subsection{要点}
完全消除所有动态类型检查既不现实也不必要，但应该将不安全的代码限制在尽可能小的范围内，并通过安全的接口将其封装起来。使用 \texttt{std::variant}、\texttt{std::any}、\texttt{enum class} 等现代 C++ 特性来替代不安全的类型转换和联合体。

% === P.5 ===
\section{规则 P.5：优先编译时检查而非运行时检查}

\textbf{规则 P.5:} \textit{优先在编译时进行检查，而非运行时。}

\subsection{为什么重要}

编译时检查是免费的——它不增加运行时开销，不会导致运行时失败，并且能在开发阶段就发现问题。\texttt{constexpr}、\texttt{static\_assert}、类型系统和模板元编程都是实现编译时检查的有力工具。

\subsection{错误示例}

\begin{lstlisting}[language=C++]
// 运行时检查——增加了运行时开销
int factorial(int n) {
    if (n < 0) {
        throw std::invalid_argument("n must be non-negative");
    }
    if (n == 0) return 1;
    return n * factorial(n - 1);
}
\end{lstlisting}

\subsection{正确示例}

\begin{lstlisting}[language=C++]
// 编译时计算——零运行时开销
constexpr int factorial(int n) {
    if (n < 0) {
        // 在 C++26 中可以用 std::unreachable()
        // 在 constexpr 上下文中，负数会导致编译错误
        return -1;  // 哨兵值，或者抛出异常
    }
    if (n == 0) return 1;
    return n * factorial(n - 1);
}

constexpr int fact5 = factorial(5);  // 编译时计算
static_assert(factorial(5) == 120);  // 编译时验证
\end{lstlisting}

\subsection{要点}
尽可能将计算和检查推到编译时。\texttt{constexpr} 函数既可以在编译时求值，也可以在运行时调用，是最灵活的编译时计算工具。\texttt{static\_assert} 则用于在编译时验证条件。

% === P.6 ===
\section{规则 P.6：编译时无法检查的应该在运行时可检查}

\textbf{规则 P.6:} \textit{编译时无法检查的内容，应该使其在运行时可检查。}

\subsection{为什么重要}

有些属性（例如函数的前置条件、输入的合法性）无法在编译时完全验证。但这不意味着我们应该放弃检查——相反，应该让这些属性在运行时容易验证，并在可能的情况下使用工具（如静态分析器、sanitizer）来辅助检查。

\subsection{错误示例}

\begin{lstlisting}[language=C++]
// 完全无法在运行时检查的隐式约定
void process(const int* data) {
    // data 必须指向至少 100 个元素的数组
    // 但没有任何机制来保证这一点
    for (int i = 0; i < 100; ++i) {
        // 使用 data[i]...
    }
}
\end{lstlisting}

\subsection{正确示例}

\begin{lstlisting}[language=C++]
// 运行时可检查
void process(std::span<const int> data) {
    Expects(data.size() >= 100);  // 前置条件，可检查
    for (size_t i = 0; i < 100; ++i) {
        // 使用 data[i]，且 span 提供边界信息
    }
}
\end{lstlisting}

\subsection{要点}
使用带有大小信息的类型（如 \texttt{std::span}、\texttt{std::string\_view}）替代裸指针，使得运行时检查成为可能。使用 \texttt{Expects()} 和 \texttt{Ensures()} 宏来标注前置条件和后置条件，使其对工具和读者都可见。

% === P.7 ===
\section{规则 P.7：尽早捕获运行时错误}

\textbf{规则 P.7:} \textit{尽早捕获运行时错误。}

\subsection{为什么重要}

错误被发现的时机越晚，诊断和修复的成本就越高。一个在输入验证阶段就能发现的错误，不应该等到它导致数据库中出现脏数据后才被察觉。

\subsection{错误示例}

\begin{lstlisting}[language=C++]
// 错误在很晚之后才被发现
void process(std::vector<int>& v) {
    // 没有检查 v 是否为空
    int first = v[0];  // 如果 v 为空，未定义行为
    // ... 大量处理 ...
    // 最终在某个深层调用中崩溃
}
\end{lstlisting}

\subsection{正确示例}

\begin{lstlisting}[language=C++]
// 在入口处就检查
void process(std::vector<int>& v) {
    Expects(!v.empty());  // 尽早检查前置条件
    int first = v.at(0);  // 使用 at() 进行边界检查
    // ...
}
\end{lstlisting}

\subsection{要点}
在函数的入口处验证前置条件，在出口处验证后置条件。使用断言、异常和标准库提供的安全接口（如 \texttt{at()} 而非 \texttt{operator[]}）来确保错误在发生点附近被捕获。

% === P.8 ===
\section{规则 P.8：不要泄漏任何资源}

\textbf{规则 P.8:} \textit{不要泄漏任何资源。}

\subsection{为什么重要}

资源泄漏（内存、文件句柄、网络连接、互斥锁等）会导致程序性能逐渐下降，最终可能导致系统崩溃。在长期运行的服务器程序中，即使是微小的泄漏也会造成严重后果。

\subsection{错误示例}

\begin{lstlisting}[language=C++]
void read_data(const char* filename) {
    FILE* f = fopen(filename, "r");
    if (!f) return;
    char* buffer = new char[1024];
    fread(buffer, 1, 1024, f);
    // 如果 fread 抛出异常，f 和 buffer 都泄漏了
    // 即使正常返回，也忘记释放了
}
\end{lstlisting}

\subsection{正确示例}

\begin{lstlisting}[language=C++]
void read_data(const std::string& filename) {
    std::ifstream f(filename);
    if (!f) return;
    std::vector<char> buffer(1024);
    f.read(buffer.data(), 1024);
    // 所有资源自动释放：ifstream 和 vector 的析构函数
}
\end{lstlisting}

\subsection{要点}
使用 RAII（资源获取即初始化）原则管理所有资源。标准库中的智能指针、容器、锁和文件流都是 RAII 的范例。避免直接使用 \texttt{new}/\texttt{delete}、\texttt{malloc}/\texttt{free} 和手动的 \texttt{open}/\texttt{close}。

% === P.9 ===
\section{规则 P.9：不要浪费时间和空间}

\textbf{规则 P.9:} \textit{不要浪费时间和空间。}

\subsection{为什么重要}

这条规则体现了 C++ 的核心设计哲学："你不需要为你不使用的东西付出代价"。在保证安全和可维护的前提下，应该选择高效的实现方式。

\subsection{错误示例}

\begin{lstlisting}[language=C++]
// 不必要的拷贝
std::string get_greeting() {
    std::string result = "Hello, World!";
    return result;  // 某些旧编译器可能无法执行 NRVO
}

void use() {
    std::string s = get_greeting();
    std::string s2 = s;       // 不必要的拷贝
    std::string s3 = s2;      // 又一次不必要的拷贝
}
\end{lstlisting}

\subsection{正确示例}

\begin{lstlisting}[language=C++]
// 避免不必要的拷贝
std::string_view get_greeting() {
    return "Hello, World!";  // 零拷贝
}

void use() {
    std::string s(get_greeting());  // 只在需要可变字符串时构造
    std::string_view sv = s;        // 视图，零拷贝
}
\end{lstlisting}

\subsection{要点}
这条规则与 P.8 看似矛盾（安全 vs. 效率），实际上它们互补：RAII 本身几乎不带来运行时开销，智能指针的开销微乎其微。应该先保证正确和安全，然后在此基础上进行有针对性的优化。不要过早优化，但也不要故意低效。

% === P.10 ===
\section{规则 P.10：优先使用不可变数据}

\textbf{规则 P.10:} \textit{优先使用不可变数据而非可变数据。}

\subsection{为什么重要}

不可变数据（immutable data）天然是线程安全的——既然不能被修改，就不存在数据竞争。不可变数据也更易于推理：你不需要追踪它的修改历史。在函数式编程中，不可变性是核心原则；在 C++ 中，\texttt{const} 和 \texttt{constexpr} 提供了类似的能力。

\subsection{错误示例}

\begin{lstlisting}[language=C++]
// 可变状态导致难以追踪的 bug
int counter = 0;

void increment() { counter++; }
void print() { std::cout << counter << std::endl; }

// 在多线程环境中，counter 的值不可预测
\end{lstlisting}

\subsection{正确示例}

\begin{lstlisting}[language=C++]
// 不可变数据
class Counter {
    int value_;
public:
    constexpr explicit Counter(int v) : value_(v) {}
    constexpr int value() const { return value_; }
    constexpr Counter incremented() const {
        return Counter(value_ + 1);  // 返回新对象，不修改原对象
    }
};

const Counter c{0};
Counter c2 = c.incremented();  // c 仍然是 0，c2 是 1
\end{lstlisting}

\subsection{要点}
默认将变量声明为 \texttt{const}，只在确实需要修改时才去掉 \texttt{const}。优先使用值语义和纯函数来替代可变状态。在并发编程中，不可变数据可以完全消除对锁的需求。

% === P.11 ===
\section{规则 P.11：封装杂乱的构造}

\textbf{规则 P.11:} \textit{封装杂乱的构造，而不是让它们散落在代码中。}

\subsection{为什么重要}

低级、容易出错的代码（如裸指针操作、C 风格 API 调用）应该被封装在紧凑的、经过充分测试的抽象之后。这样，代码库的大部分都是安全、清晰的高层代码，而复杂性被限制在少数几个经过审查的模块中。

\subsection{错误示例}

\begin{lstlisting}[language=C++]
// 杂乱的 C API 调用散落在业务逻辑中
void display_image(const char* path) {
    FILE* f = fopen(path, "rb");
    fseek(f, 0, SEEK_END);
    long size = ftell(f);
    fseek(f, 0, SEEK_SET);
    char* data = (char*)malloc(size);
    fread(data, 1, size, f);
    fclose(f);
    // ... 图像处理 ...
    free(data);
}
\end{lstlisting}

\subsection{正确示例}

\begin{lstlisting}[language=C++]
// 封装在 RAII 类中
class ImageFile {
    std::vector<char> data_;
public:
    explicit ImageFile(const std::string& path) {
        std::ifstream f(path, std::ios::binary);
        f.seekg(0, std::ios::end);
        data_.resize(f.tellg());
        f.seekg(0, std::ios::beg);
        f.read(data_.data(), data_.size());
    }
    std::span<const char> data() const { return data_; }
};

void display_image(const std::string& path) {
    ImageFile img(path);  // 所有 messy 细节被封装
    process(img.data());
}
\end{lstlisting}

\subsection{要点}
将不安全、低层的代码封装在小的、有明确接口的函数或类中。这样，代码库中 95\% 的代码可以是安全、高层的，只有 5\% 的代码需要处理底层细节。这 5\% 的代码需要更仔细的审查和测试。

% === P.12 ===
\section{规则 P.12：适当使用辅助工具}

\textbf{规则 P.12:} \textit{适当使用辅助工具。}

\subsection{为什么重要}

现代 C++ 开发有丰富的工具支持：编译器警告、静态分析器（如 clang-tidy、Coverity）、动态分析器（如 Valgrind、AddressSanitizer）、格式化工具（如 clang-format）等。这些工具可以发现人工审查难以察觉的错误。

\subsection{错误示例}

\begin{lstlisting}[language=C++]
// 不使用任何工具，完全依赖人工审查
// 以下代码中的问题可能逃过人眼：
void process(int* p, int n) {
    for (int i = 0; i <= n; ++i) {  // off-by-one
        p[i] *= 2;
    }
    char buf[10];
    strcpy(buf, "long string");  // 缓冲区溢出
}
\end{lstlisting}

\subsection{正确示例}

\begin{lstlisting}[language=C++]
// 使用工具链来捕获问题
// 编译时：-Wall -Wextra -Werror -Wpedantic
// 静态分析：clang-tidy
// 运行时：-fsanitize=address,undefined
void process(std::span<int> data) {
    for (auto& elem : data) {  // 安全：范围 for
        elem *= 2;
    }
    std::string buf = "long string";  // 安全：自动管理内存
}
\end{lstlisting}

\subsection{要点}
将工具集成到开发流程中。启用所有编译器警告并将它们视为错误。在 CI/CD 流水线中运行静态分析和 sanitizer。这些工具的投入产出比极高。

% === P.13 ===
\section{规则 P.13：适当使用支持库}

\textbf{规则 P.13:} \textit{适当使用支持库。}

\subsection{为什么重要}

C++ 标准库提供了大量经过优化和测试的组件。使用标准库而非自己造轮子，可以减少 bug、提高性能、增强可移植性。同样，经过验证的第三方库（如 Boost、fmt、nlohmann/json）也值得信赖。

\subsection{错误示例}

\begin{lstlisting}[language=C++]
// 自己实现链表——容易出错且低效
template<typename T>
struct Node {
    T data;
    Node* next;
};

// 自己实现字符串分割
std::vector<std::string> split(const std::string& s, char delim) {
    std::vector<std::string> result;
    // 50 行容易出错的代码...
    return result;
}
\end{lstlisting}

\subsection{正确示例}

\begin{lstlisting}[language=C++]
// 使用标准库
std::list<int> lst;  // 标准库链表
std::vector<int> vec;  // 通常更好的选择

// 使用标准库算法
std::vector<std::string> split(const std::string& s, char delim) {
    std::vector<std::string> result;
    std::istringstream iss(s);
    std::string token;
    while (std::getline(iss, token, delim)) {
        result.push_back(token);
    }
    return result;
}
\end{lstlisting}

\subsection{要点}
优先使用标准库的容器、算法和工具。当标准库不能满足需求时，考虑使用经过社区验证的第三方库。只有在充分理由的情况下才自己实现基础设施，并且要确保实现经过充分测试。

\section*{练习题}

\begin{enumerate}
  \item 选择规则 P.1 到 P.3 中的一条，在你熟悉的代码库中找到一个违反该规则的实例，并说明如何修复。
  \item 规则 P.8（不泄漏资源）和规则 P.9（不浪费时间和空间）看似可能冲突。讨论在什么情况下它们可能冲突，以及如何平衡。
  \item 使用 \texttt{constexpr} 重写一个你熟悉的运行时计算函数，使其能在编译时求值。使用 \texttt{static\_assert} 验证结果。
  \item 设计一个 RAII 包装类，封装 POSIX 的 \texttt{socket}/\texttt{close} 接口，确保套接字资源不会泄漏。
  \item 解释为什么"封装杂乱的构造"（P.11）与"直接在代码中表达思想"（P.1）并不矛盾。
\end{enumerate}


% ============================================================
% Part II: 接口与函数
% ============================================================

\part{接口与函数}

% ------------------------------------------------------------
% Chapter 3: 接口设计
% ------------------------------------------------------------
\chapter{接口设计}

接口是代码模块之间的契约。好的接口使得使用者无需了解实现细节就能正确地使用模块；坏的接口则导致误用、耦合和维护噩梦。本章涵盖核心指南中关于接口设计的规则。

\section{使接口显式化}

\subsection{规则 I.1：使接口显式化}

\textbf{规则 I.1:} \textit{使接口显式化。}

\subsection{为什么重要}

函数的签名（参数的类型和数量、返回类型、是否 \texttt{const}、是否 \texttt{noexcept}）是调用者与实现之间的第一道契约。如果接口的信息被隐藏在注释或约定中，就容易被忽视和违反。

\subsection{错误示例}

\begin{lstlisting}[language=C++]
// 接口信息隐藏在注释中
void draw(int x, int y, int w, int h);
// x, y 是左上角坐标
// w, h 是宽度和高度
// 调用者需要查看文档才能知道参数含义
\end{lstlisting}

\subsection{正确示例}

\begin{lstlisting}[language=C++]
// 接口信息显式化
struct Point { int x; int y; };
struct Size  { int w; int h; };

void draw(Point top_left, Size dimensions);
// 参数含义一目了然，且类型系统防止参数顺序错误
\end{lstlisting}

\subsection{要点}
使用强类型来替代注释。如果两个参数的类型相同但含义不同（如 \texttt{int x, int y}），调用者很容易搞混顺序。使用不同的类型（如 \texttt{Point} 和 \texttt{Size}）可以让编译器帮你捕获这类错误。

\subsection{规则 I.2：避免非 const 全局变量}

\textbf{规则 I.2:} \textit{避免非 const 全局变量。}

全局变量破坏了模块化——任何代码都可能在任何时候修改它们，这使得推理程序行为变得极其困难。全局变量还导致线程安全问题，并且使单元测试变得复杂（因为测试之间可能通过全局变量产生隐式依赖）。

\begin{lstlisting}[language=C++]
// 错误：全局可变状态
int g_counter = 0;
std::string g_config_path;

void increment() { g_counter++; }
\end{lstlisting}

\begin{lstlisting}[language=C++]
// 正确：使用局部变量或通过参数传递
class Counter {
    int value_ = 0;
public:
    void increment() { ++value_; }
    int value() const { return value_; }
};

// 如果确实需要全局配置，使用 const 或受控访问
const std::string config_path = "/etc/app.conf";
\end{lstlisting}

\subsection{规则 I.3：避免单例}

\textbf{规则 I.3:} \textit{避免单例。}

单例本质上是伪装成面向对象的全局变量。它带来与全局变量相同的问题：隐式依赖、难以测试、线程安全问题。

\begin{lstlisting}[language=C++]
// 错误：单例
class Database {
public:
    static Database& instance() {
        static Database db;
        return db;
    }
    void query(const std::string& sql);
private:
    Database() = default;
};

// 使用
Database::instance().query("SELECT ...");
\end{lstlisting}

\begin{lstlisting}[language=C++]
// 正确：依赖注入
class Database {
public:
    explicit Database(const std::string& conn_str);
    void query(const std::string& sql);
};

class UserService {
    Database& db_;  // 通过构造函数注入
public:
    explicit UserService(Database& db) : db_(db) {}
    User find_user(int id) { /* 使用 db_ */ }
};
\end{lstlisting}

\section{精确和强类型的接口}

\subsection{规则 I.4：使接口精确和强类型}

\textbf{规则 I.4:} \textit{使接口精确和强类型。}

精确和强类型的接口能在编译时捕获错误。当参数类型精确地反映了其含义和约束时，许多常见的编程错误就不可能发生。

\begin{lstlisting}[language=C++]
// 不精确的接口
void set_speed(double value);  // 单位是什么？km/h? m/s?
void connect(int timeout);     // 毫秒？秒？
\end{lstlisting}

\begin{lstlisting}[language=C++]
// 精确的接口
enum class SpeedUnit { km_per_hour, meters_per_second };
void set_speed(double value, SpeedUnit unit);

// 或者更好的方式
class Speed {
    double meters_per_second_;
public:
    static Speed from_kmh(double kmh) {
        return Speed{kmh / 3.6};
    }
    static Speed from_mps(double mps) {
        return Speed{mps};
    }
    double kmh() const { return meters_per_second_ * 3.6; }
    double mps() const { return meters_per_second_; }
};

using namespace std::chrono_literals;
void connect(5s);  // 使用 chrono _DURATION 类型，无歧义
\end{lstlisting}

\section{前置条件与后置条件}

\subsection{规则 I.5：明确声明前置条件}

\textbf{规则 I.5:} \textit{明确声明前置条件。}

前置条件是调用者有责任满足的条件。使用 \texttt{Expects()} 宏（来自 GSL）来显式声明前置条件，使其对读者和工具都可见。

\begin{lstlisting}[language=C++]
// 未声明前置条件
double sqrt(double x);  // x 必须 >= 0，但接口没有说明
\end{lstlisting}

\begin{lstlisting}[language=C++]
// 显式声明前置条件
double my_sqrt(double x) {
    Expects(x >= 0);
    // ...
}

// 更好的方式：用类型系统消除前置条件
double sqrt_non_negative(double x) {
    // x 的类型是 double，但函数名表达了约束
    // 或者使用 std::optional<double> 来处理负数输入
}
\end{lstlisting}

\subsection{规则 I.6：优先用类型表达前置条件}

\textbf{规则 I.6:} \textit{优先使用类型来表达前置条件，使其不可能被违反。}

如果前置条件可以用类型系统表达，那么违反它就会成为编译错误，而非运行时检查。这比 \texttt{Expects()} 更强。

\begin{lstlisting}[language=C++]
// 使用 Expects 的前置条件——运行时检查
void set_age(int age) {
    Expects(age >= 0 && age < 150);
    // ...
}
\end{lstlisting}

\begin{lstlisting}[language=C++]
// 使用类型系统——编译时保证
class Age {
    int value_;
    Age(int v) : value_(v) {
        if (v < 0 || v >= 150)
            throw std::invalid_argument("invalid age");
    }
public:
    static Age make(int v) { return Age(v); }
    int value() const { return value_; }
};

void set_age(Age age);  // 不可能传入无效值
\end{lstlisting}

\subsection{规则 I.7：明确声明后置条件}

\textbf{规则 I.7:} \textit{明确声明后置条件。}

后置条件是函数保证在返回时成立的条件。对于没有副作用的函数，后置条件由返回值完全描述。对于有副作用的函数，\texttt{Ensures()} 宏可以帮助文档化和验证后置条件。

\begin{lstlisting}[language=C++]
// 后置条件不明确
void sort(std::vector<int>& v);
// 排序后 v 是有序的——但这只是约定
\end{lstlisting}

\begin{lstlisting}[language=C++]
// 显式声明后置条件
void my_sort(std::vector<int>& v) {
    auto copy = v;
    // ... 排序实现 ...
    Ensures(std::is_sorted(v.begin(), v.end()));
    Ensures(std::is_permutation(v.begin(), v.end(), copy.begin()));
}
\end{lstlisting}

\subsection{规则 I.8：尽量使接口具有强异常安全性}

\textbf{规则 I.8:} \textit{尽量使接口具有强异常安全性保证。}

强异常安全性保证意味着：如果函数抛出异常，程序状态就像该函数从未被调用过一样。这极大地简化了错误处理逻辑。

\begin{lstlisting}[language=C++]
// 弱异常安全性——出错后对象处于未知状态
void Account::transfer_to(Account& target, double amount) {
    balance_ -= amount;      // 第一步：扣款
    target.balance_ += amount;  // 第二步：入账——如果这里抛异常？
    // balance_ 已经被扣款但 target 没有入账！
}
\end{lstlisting}

\begin{lstlisting}[language=C++]
// 强异常安全性
void Account::transfer_to(Account& target, double amount) {
    // 先准备所有可能失败的操作
    auto guard = std::make_unique<Guard>(*this, amount);
    target.balance_ += amount;
    balance_ -= amount;
    guard->dismiss();  // 全部成功，取消回滚
}
\end{lstlisting}

\section{模板接口与概念}

\subsection{规则 I.9：当接口是模板时，用概念（Concepts）记录参数}

\textbf{规则 I.9:} \textit{当接口是模板时，使用概念（Concepts）来明确记录模板参数的要求。}

在 C++20 引入概念之前，模板参数的要求只能通过文档和 \texttt{static\_assert} 来表达。概念使得模板参数的要求成为接口的一部分，编译器可以在实例化时检查这些要求，并给出更清晰的错误信息。

\begin{lstlisting}[language=C++]
// 没有概念——错误信息难以理解
template<typename T>
T add(T a, T b) {
    return a + b;
}
// add(std::vector<int>{}, std::vector<int>{});
// 编译器报错可能涉及数十行模板实例化信息
\end{lstlisting}

\begin{lstlisting}[language=C++]
// 使用概念——错误信息清晰
template<std::regular T>
T add(T a, T b) {
    return a + b;
}

// 或者自定义概念
template<typename T>
concept Addable = requires(T a, T b) {
    { a + b } -> std::convertible_to<T>;
};

template<Addable T>
T add(T a, T b) {
    return a + b;
}
\end{lstlisting}

\section{错误处理与所有权}

\subsection{规则 I.10：使用异常表示失败}

\textbf{规则 I.10:} \textit{不要使用返回码来处理失败。使用异常。}

返回码容易被忽略，而且会污染正常逻辑。异常不能被静默忽略，并且可以将错误处理逻辑与正常逻辑分离。

\begin{lstlisting}[language=C++]
// 使用返回码——容易被忽略
int connect(const char* host) {
    if (/* 连接失败 */) return -1;
    if (/* 认证失败 */) return -2;
    return 0;  // 成功
}

int fd = connect("example.com");
// 忘记检查返回值——错误被忽略
\end{lstlisting}

\begin{lstlisting}[language=C++]
// 使用异常——无法被忽略
class ConnectionError : public std::runtime_error {
public:
    using std::runtime_error::runtime_error;
};

Connection connect(const std::string& host) {
    if (/* 连接失败 */)
        throw ConnectionError("Cannot connect to " + host);
    // ...
}

auto conn = connect("example.com");  // 异常会自动传播
\end{lstlisting}

\subsection{规则 I.11：永远不要通过裸指针或裸引用转移所有权}

\textbf{规则 I.11:} \textit{永远不要通过裸指针（raw pointer）或裸引用转移所有权。使用 \texttt{unique\_ptr} 或 \texttt{shared\_ptr}。}

如果函数返回一个裸指针，调用者无法知道是否需要释放它，或者它指向的对象何时失效。

\begin{lstlisting}[language=C++]
// 所有权不明确
Widget* create_widget() {
    return new Widget();  // 调用者需要 delete 吗？
}

Widget* get_widget() {
    static Widget w;
    return &w;  // 这个呢？能 delete 吗？
}
\end{lstlisting}

\begin{lstlisting}[language=C++]
// 所有权明确
std::unique_ptr<Widget> create_widget() {
    return std::make_unique<Widget>();  // 明确：调用者拥有所有权
}

Widget& get_widget() {
    static Widget w;
    return w;  // 明确：调用者不拥有所有权
}
\end{lstlisting}

\subsection{规则 I.12：使用 not\_null 指针表明不应为 null}

\textbf{规则 I.12:} \textit{如果指针不应为 null，使用 \texttt{not\_null} 或引用。}

\begin{lstlisting}[language=C++]
// 不确定是否可以为 null
int get_length(const char* str);
// get_length(nullptr) 会怎样？
\end{lstlisting}

\begin{lstlisting}[language=C++]
// 明确不可为 null
int get_length(gsl::not_null<const char*> str);
// 或者使用引用
int get_length(const char& first_char);
// 或者使用 string_view
int get_length(std::string_view str);
\end{lstlisting}

\subsection{规则 I.13：不要将数组作为单指针传递}

\textbf{规则 I.13:} \textit{不要将数组作为单个指针传递。}

\begin{lstlisting}[language=C++]
// 丢失了大小信息
void print(const int* arr);  // arr 有多少个元素？
\end{lstlisting}

\begin{lstlisting}[language=C++]
// 保留大小信息
void print(std::span<const int> arr);
// 或
void print(const std::vector<int>& arr);
\end{lstlisting}

\section{复杂初始化与参数设计}

\subsection{规则 I.22：避免复杂初始化}

\textbf{规则 I.22:} \textit{对于使用复杂或默认构造的对象，避免隐式初始化。}

当对象的初始化涉及多个步骤或可能失败时，应该使用显式的工厂函数或构造函数来确保对象在创建时就处于有效状态。

\begin{lstlisting}[language=C++]
// 复杂且容易出错的初始化
Socket s;
s.set_protocol(TCP);
s.set_timeout(30);
s.bind("0.0.0.0", 8080);
s.listen();  // 如果忘记调用 listen 就使用 s？
\end{lstlisting}

\begin{lstlisting}[language=C++]
// 工厂函数确保完整初始化
class Socket {
    Socket(/* 参数 */) { /* 完整初始化 */ }
public:
    static Socket create_tcp(int port) {
        Socket s(/* ... */);
        s.bind_and_listen(port);
        return s;
    }
};

auto s = Socket::create_tcp(8080);  // 保证已完全初始化
\end{lstlisting}

\subsection{规则 I.23：保持参数数量较少}

\textbf{规则 I.23:} \textit{保持函数的参数数量较少。}

参数过多是接口设计不良的信号。超过三到四个参数通常意味着需要将这些参数组织成结构体，或者函数承担了过多的职责。

\begin{lstlisting}[language=C++]
// 参数过多
void create_user(const std::string& first,
                 const std::string& last,
                 const std::string& email,
                 int age,
                 const std::string& phone,
                 const std::string& address);
\end{lstlisting}

\begin{lstlisting}[language=C++]
// 使用结构体组织参数
struct UserInfo {
    std::string first_name;
    std::string last_name;
    std::string email;
    int age;
    std::string phone;
    std::string address;
};

void create_user(const UserInfo& info);
\end{lstlisting}

\subsection{规则 I.24：避免相邻的同类型参数}

\textbf{规则 I.24:} \textit{避免相邻的参数具有相同的类型且可以互换。}

相邻的同类型参数极易被搞混，编译器无法帮你检测这种错误。

\begin{lstlisting}[language=C++]
// 容易搞混参数顺序
void copy(char* dest, const char* src, size_t n);
// 调用时：copy(buf, input, len) 还是 copy(input, buf, len)？
\end{lstlisting}

\begin{lstlisting}[language=C++]
// 使用不同类型来防止搞混
void copy(std::span<char> dest, std::span<const char> src);
// 编译器确保不会搞混 dest 和 src
\end{lstlisting}

\subsection{规则 I.25：优先使用抽象类作为接口}

\textbf{规则 I.25:} \textit{优先使用抽象类作为接口。}

抽象类（纯虚基类）定义了清晰的接口契约，强制派生类实现特定方法，并且支持多态。

\begin{lstlisting}[language=C++]
// 好的接口设计
class ISensor {
public:
    virtual ~ISensor() = default;
    virtual double read() = 0;
    virtual std::string name() const = 0;
};

class TemperatureSensor : public ISensor {
public:
    double read() override;
    std::string name() const override;
};
\end{lstlisting}

\subsection{规则 I.26--I.27：ABI 稳定与 Pimpl 惯用法}

\textbf{规则 I.26:} \textit{如果你想跨"ABI 边界"保留一组实现，使用 Pimpl 惯用法。}

\textbf{规则 I.27:} \textit{为了稳定的 ABI，将类声明与实现分离。}

Pimpl（Pointer to Implementation）惯用法将类的实现细节隐藏在指针之后，使得修改实现不影响 ABI。

\begin{lstlisting}[language=C++]
// 头文件
class Widget {
public:
    Widget();
    ~Widget();
    Widget(const Widget& rhs);
    Widget& operator=(const Widget& rhs);
    void draw();
private:
    struct Impl;
    std::unique_ptr<Impl> pimpl_;
};

// 实现文件
struct Widget::Impl {
    int internal_data;
    std::string internal_string;
    // 修改这些不影响 Widget 的 ABI
};

Widget::Widget() : pimpl_(std::make_unique<Impl>()) {}
Widget::~Widget() = default;
void Widget::draw() { /* 使用 pimpl_ */ }
\end{lstlisting}

\subsection{规则 I.30：封装变化的部分}

\textbf{规则 I.30:} \textit{封装变化的部分。}

这是面向对象设计中最核心的原则之一。识别系统中可能变化的部分，将它们封装在接口之后，使得修改不影响其他部分。

\begin{lstlisting}[language=C++]
// 未封装变化——硬编码了数据库类型
class UserService {
    SQLiteDatabase db_;  // 直接依赖具体实现
public:
    User find(int id) { return db_.query<User>(id); }
};
\end{lstlisting}

\begin{lstlisting}[language=C++]
// 封装变化——通过接口解耦
class IDatabase {
public:
    virtual ~IDatabase() = default;
    virtual User find_user(int id) = 0;
};

class UserService {
    std::unique_ptr<IDatabase> db_;
public:
    explicit UserService(std::unique_ptr<IDatabase> db)
        : db_(std::move(db)) {}
    User find(int id) { return db_->find_user(id); }
};
\end{lstlisting}

\section*{练习题}

\begin{enumerate}
  \item 设计一个表示二维平面上"矩形"的类，使接口尽可能精确和强类型。考虑如何防止宽度和高度为负值。
  \item 将以下函数签名改进为更安全的接口：\texttt{void process(char* buf, int size, char* out, int out\_size)}。
  \item 使用 Pimpl 惯用法重构一个你熟悉的类，解释这样做的好处和代价。
  \item 为一个"消息队列"系统设计接口，使用抽象类来定义接口，使用工厂函数来创建实例。
  \item 解释为什么 \texttt{std::span} 比 \texttt{(pointer, size)} 对更安全，列出至少三个理由。
\end{enumerate}


% ------------------------------------------------------------
% Chapter 4: 函数设计
% ------------------------------------------------------------
\chapter{函数设计}

函数是 C++ 程序的基本构建块。好的函数设计使得代码易于理解、测试和维护。本章按照核心指南的分类，从定义、参数传递、返回值和其他方面来讨论函数设计的最佳实践。

\section{函数定义}

\subsection{规则 F.1：将每个合理的操作封装为函数}

\textbf{规则 F.1:} \textit{"合理"的操作应该打包为函数。}

\subsection{为什么重要}

将操作封装为函数提供了命名（自文档化）、参数化、可重用性和可测试性。但函数也不应过小——如果一个函数只有一行且只在一个地方使用，直接内联可能更清晰。

\subsection{错误示例}

\begin{lstlisting}[language=C++]
// 200 行的 main 函数，所有逻辑都在里面
int main() {
    // 读取配置
    // 初始化数据库
    // 启动网络服务
    // 处理请求循环
    // 清理资源
    // ... 全部写在一起 ...
}
\end{lstlisting}

\subsection{正确示例}

\begin{lstlisting}[language=C++]
int main() {
    auto config = load_config("app.conf");
    auto db = initialize_database(config);
    auto server = start_server(config, db);
    server.run();
    // 清理由 RAII 自动处理
}
\end{lstlisting}

\subsection{规则 F.2：函数应该执行单一操作}

\textbf{规则 F.2:} \textit{函数应该执行单一且完整的操作。}

如果一个函数名需要用"和"来描述（如"读取文件\textbf{和}解析数据\textbf{和}写入数据库"），它很可能应该被拆分为多个函数。

\begin{lstlisting}[language=C++]
// 做了太多事情
void process_order(Order& order) {
    // 验证订单
    // 计算价格
    // 检查库存
    // 扣减库存
    // 生成发票
    // 发送邮件通知
    // 更新统计信息
}
\end{lstlisting}

\begin{lstlisting}[language=C++]
// 每个函数做一件事
void process_order(Order& order) {
    validate_order(order);
    auto price = calculate_price(order);
    if (check_and_reserve_inventory(order)) {
        auto invoice = generate_invoice(order, price);
        send_notification(order, invoice);
        update_statistics(order);
    }
}
\end{lstlisting}

\subsection{规则 F.3：保持函数短小}

\textbf{规则 F.3:} \textit{保持函数短小。}

长函数难以理解、难以测试、难以维护。虽然没有绝对的行数限制，但一个函数如果超过一屏（约 40--60 行），通常意味着它需要拆分。

\subsection{规则 F.4：如果函数可能在编译时求值，使其成为 constexpr}

\textbf{规则 F.4:} \textit{如果函数有可能在编译时求值，就将其声明为 \texttt{constexpr}。}

\begin{lstlisting}[language=C++]
// 编译时求值
constexpr int fibonacci(int n) {
    if (n <= 1) return n;
    int a = 0, b = 1;
    for (int i = 2; i <= n; ++i) {
        int temp = a + b;
        a = b;
        b = temp;
    }
    return b;
}

constexpr int fib10 = fibonacci(10);  // 编译时计算
static_assert(fib10 == 55);
\end{lstlisting}

\subsection{规则 F.5：小型函数使用 inline}

\textbf{规则 F.5:} \textit{对于非常小型的、性能关键的函数，使用 \texttt{inline}。}

在类定义中定义的成员函数隐式地是 \texttt{inline} 的。对于头文件中的小型辅助函数，\texttt{inline} 可以避免多重定义错误。但现代编译器的内联优化不依赖于 \texttt{inline} 关键字——它只是一个链接提示。

\subsection{规则 F.6：使用 noexcept 标记不抛出异常的函数}

\textbf{规则 F.6:} \textit{使用 \texttt{noexcept} 标记不会抛出异常的函数。}

\texttt{noexcept} 不仅是对读者的文档——它还给编译器优化提供了信息，并且影响移动语义（\texttt{std::vector} 在重新分配时只使用 \texttt{noexcept} 的移动操作）。

\begin{lstlisting}[language=C++]
class Widget {
    int value_;
public:
    // 明确不会抛出异常
    int value() const noexcept { return value_; }

    // noexcept 的移动操作使得 vector 可以高效地重新分配
    Widget(Widget&& other) noexcept
        : value_(other.value_) { other.value_ = 0; }

    Widget& operator=(Widget&& other) noexcept {
        value_ = other.value_;
        other.value_ = 0;
        return *this;
    }
};
\end{lstlisting}

\subsection{规则 F.7：对于所有权转移，使用 unique\_ptr 而非裸指针}

\textbf{规则 F.7:} \textit{如果函数需要转移所有权，使用 \texttt{unique\_ptr} 而非裸指针或 \texttt{owner}。}

\begin{lstlisting}[language=C++]
// 不明确
Widget* make_widget();  // 需要 delete 吗？
\end{lstlisting}

\begin{lstlisting}[language=C++]
// 明确
std::unique_ptr<Widget> make_widget();  // 明确转移所有权
\end{lstlisting}

\subsection{规则 F.8：优先使用纯函数}

\textbf{规则 F.8:} \textit{优先使用纯函数。}

纯函数是指给定相同的输入总是返回相同的结果，并且没有副作用的函数。纯函数易于理解、测试、并行化和优化。

\begin{lstlisting}[language=C++]
// 纯函数
constexpr int add(int a, int b) noexcept {
    return a + b;
}

// 非纯函数——依赖外部状态
int g_offset = 0;
int add_with_offset(int a, int b) {
    return a + b + g_offset;
}
\end{lstlisting}

\subsection{规则 F.9：未使用的参数应该去掉}

\textbf{规则 F.9:} \textit{未使用的函数参数应该去掉。}

未使用的参数增加了认知负担，使读者困惑。如果因为接口兼容性必须保留参数，使用 \texttt{[[maybe\_unused]]} 属性或注释掉参数名。

\begin{lstlisting}[language=C++]
// 未使用的参数
void callback(int event_type, void* data, int flags) {
    // 只使用了 event_type
    handle(event_type);
}

// 正确：去掉未使用的参数，或标注
void callback(int event_type,
              void* /*data*/,
              int /*flags*/) {
    handle(event_type);
}
\end{lstlisting}

\subsection{规则 F.10--F.11：可复用操作与 Lambda}

\textbf{规则 F.10:} \textit{如果操作可以被重用，将其转化为独立的可复用操作。}

\textbf{规则 F.11:} \textit{使用 Lambda 来封装简单的局部操作。}

\begin{lstlisting}[language=C++]
// 使用 Lambda 封装局部逻辑
auto normalize = [](std::vector<double>& v) {
    double sum = std::accumulate(v.begin(), v.end(), 0.0);
    for (auto& x : v) x /= sum;
};

std::vector<double> data = {1.0, 2.0, 3.0};
normalize(data);
\end{lstlisting}

\section{参数传递}

\subsection{规则 F.15：按照惯例传递参数}

\textbf{规则 F.15:} \textit{按照惯例传递参数：廉价的用值传递，其他的用 const 引用传递。}

这是核心指南中最实用的规则之一。它定义了参数传递的标准方式：

\begin{itemize}
  \item 对于廉价的类型（如 \texttt{int}、\texttt{double}、\texttt{char}、小结构体、指针），使用普通的值传递。
  \item 对于其他类型，使用 \texttt{const T\&} 传递只读参数。
  \item 对于需要修改的参数，使用 \texttt{T\&} 传递（输入输出参数）。
  \item 对于需要转移所有权的参数，使用 \texttt{unique\_ptr<T>}。
  \item 对于需要共享所有权的参数，使用 \texttt{shared\_ptr<T>}。
\end{itemize}

\begin{lstlisting}[language=C++]
// 正确的参数传递方式
void process(int n,                    // 廉价类型：值传递
             const std::string& name,  // 只读大对象：const 引用
             std::vector<int>& output, // 输出参数：非 const 引用
             std::unique_ptr<Widget> widget);  // 所有权转移
\end{lstlisting}

\subsection{规则 F.16--F.17：输入参数}

\textbf{规则 F.16:} \textit{对于输入参数，廉价的用值传递，其他的用 const 引用传递。}

\textbf{规则 F.17:} \textit{对于"仅输入"的参数，如果调用者不需要看到修改结果，使用按值传递后移动。}

\begin{lstlisting}[language=C++]
// 对于需要拷贝的输入参数，按值传递然后移动
void set_name(std::string name) {  // 按值传递
    name_ = std::move(name);       // 移动赋值
}

// 调用
set_name("hello");           // 一次移动
set_name(std::move(my_str)); // 一次移动（不拷贝）
\end{lstlisting}

\subsection{规则 F.18--F.21：输入输出参数}

\textbf{规则 F.18:} \textit{对于"将被修改"的参数，使用非 const 引用传递。}

\textbf{规则 F.19:} \textit{对于"前向移动"参数，使用 \texttt{Forward} 类型参数。}

\textbf{规则 F.20:} \textit{对于"输出"值，优先使用返回值而非输出参数。}

\textbf{规则 F.21:} \textit{要返回多个"输出"值，优先返回结构体或使用元组。}

\begin{lstlisting}[language=C++]
// 不好：通过输出参数返回
void find_min_max(const std::vector<int>& v, int& min, int& max) {
    min = *std::min_element(v.begin(), v.end());
    max = *std::max_element(v.begin(), v.end());
}
\end{lstlisting}

\begin{lstlisting}[language=C++]
// 好：返回结构体
struct MinMax { int min; int max; };

MinMax find_min_max(const std::vector<int>& v) {
    return {
        *std::min_element(v.begin(), v.end()),
        *std::max_element(v.begin(), v.end())
    };
}

auto [min, max] = find_min_max(data);  // C++17 结构化绑定
\end{lstlisting}

\subsection{规则 F.22--F.26：各种参数类型}

\textbf{规则 F.22:} \textit{使用 \texttt{T*} 或 \texttt{T\&} 参数来表示单个对象（不表示数组）。}

\textbf{规则 F.23:} \textit{使用 \texttt{not\_null<T>} 来表示不允许 null 的指针。}

\textbf{规则 F.24:} \textit{使用 \texttt{span<T>} 或 \texttt{(T*, size\_t)} 来表示半开区间。}

\textbf{规则 F.25:} \textit{使用 \texttt{zstring} 或 \texttt{not\_null<zstring>} 来表示 C 风格字符串。}

\textbf{规则 F.26:} \textit{对于数据传递，使用 \texttt{unique\_ptr} 而非裸指针。}

\begin{lstlisting}[language=C++]
// 使用 span 传递数组视图
void process(std::span<const int> data);

// 使用 zstring 传递 C 风格字符串
void print(gsl::zstring msg);

// 使用 unique_ptr 传递所有权
void take_ownership(std::unique_ptr<Widget> w);
\end{lstlisting}

\section{返回值}

\subsection{规则 F.42：返回位置时不要返回指针}

\textbf{规则 F.42:} \textit{返回位置（索引、迭代器）时，不要返回指针。}

指针可能被误认为指向一个对象，而实际上它表示的是一个位置。使用迭代器或索引更清晰。

\begin{lstlisting}[language=C++]
// 不好：返回指针表示位置
int* find(std::vector<int>& v, int target) {
    for (auto& elem : v) {
        if (elem == target) return &elem;
    }
    return nullptr;
}
\end{lstlisting}

\begin{lstlisting}[language=C++]
// 好：返回迭代器
auto it = std::find(v.begin(), v.end(), target);
\end{lstlisting}

\subsection{规则 F.43：不要直接返回指向局部变量的指针（引用）}

\textbf{规则 F.43:} \textit{永远不要（直接或间接）返回指向局部变量的指针。}

这是 C++ 中最常见的错误之一。局部变量在函数返回后被销毁，指向它的指针成为悬垂指针。

\begin{lstlisting}[language=C++]
// 严重错误
int& get_value() {
    int local = 42;
    return local;  // 返回局部变量的引用——未定义行为！
}
\end{lstlisting}

\begin{lstlisting}[language=C++]
// 正确：返回值或使用动态分配
int get_value() {
    return 42;  // 返回值——安全
}

// 或使用 static（注意线程安全）
int& get_singleton_value() {
    static int value = 42;
    return value;
}
\end{lstlisting}

\subsection{规则 F.44--F.49：返回值的其他规则}

\textbf{规则 F.44:} \textit{当不需要拷贝时，返回 \texttt{T\&}。}

\textbf{规则 F.45:} \textit{不要返回 \texttt{T\&\&}（右值引用）来绑定到临时对象。}

\textbf{规则 F.46:} \textit{\texttt{main()} 的返回类型是 \texttt{int}。}

\textbf{规则 F.47:} \textit{赋值运算符返回 \texttt{T\&}。}

\textbf{规则 F.48:} \textit{不要 \texttt{return std::move(local)}。}

\textbf{规则 F.49:} \textit{不要返回 \texttt{const T}（值类型）。}

\begin{lstlisting}[language=C++]
// 错误：return std::move(local) 阻止了 NRVO
Widget make_widget() {
    Widget w;
    // ... 初始化 w ...
    return std::move(w);  // 阻止了编译器优化！
}

// 正确：直接返回
Widget make_widget() {
    Widget w;
    // ... 初始化 w ...
    return w;  // 编译器自动应用 NRVO 或移动语义
}
\end{lstlisting}

\begin{lstlisting}[language=C++]
// 错误：返回 const T 阻止了移动语义
const std::string get_name();  // const 返回值无意义

// 正确
std::string get_name();
\end{lstlisting}

\section{Lambda 与其他规则}

\subsection{规则 F.50：当需要闭包时使用 Lambda}

\textbf{规则 F.50:} \textit{当函数需要捕获局部状态时使用 Lambda。}

Lambda 是 C++11 引入的最有用的特性之一。它允许在调用点定义匿名函数，并捕获周围的变量。

\begin{lstlisting}[language=C++]
// Lambda 捕获局部变量
double threshold = 0.5;
auto above_threshold = [threshold](double value) {
    return value > threshold;
};

std::vector<double> data = {0.1, 0.6, 0.3, 0.8};
auto count = std::count_if(data.begin(), data.end(), above_threshold);
\end{lstlisting}

\subsection{规则 F.51--F.52：参数绑定与 Lambda 捕获}

\textbf{规则 F.51:} \textit{当需要绑定参数时，优先使用 Lambda 而非 \texttt{std::bind}。}

\textbf{规则 F.52:} \textit{谨慎使用按引用捕获。}

按引用捕获的 Lambda 在被调用时，引用的变量必须仍然存活。如果 Lambda 的生命周期可能超过被捕获的变量（例如在异步操作中），应该按值捕获。

\begin{lstlisting}[language=C++]
// 危险：按引用捕获局部变量
std::function<void()> make_callback() {
    int local = 42;
    return [&local]() {  // 悬垂引用！
        std::cout << local;
    };
}

// 安全：按值捕获
std::function<void()> make_callback_safe() {
    int local = 42;
    return [local]() {  // 拷贝到闭包中
        std::cout << local;
    };
}
\end{lstlisting}

\subsection{规则 F.53--F.56：其他函数设计规则}

\textbf{规则 F.53:} \textit{避免捕获非局部变量（包括 \texttt{this}）。}

\textbf{规则 F.54:} \textit{当使用 \texttt{this} 捕获时，显式写 \texttt{capture this}。}

\textbf{规则 F.55:} \textit{不要使用 \texttt{va\_arg} 宏。}

\textbf{规则 F.56:} \textit{避免过深的条件嵌套——使用提前返回和辅助函数。}

\begin{lstlisting}[language=C++]
// 错误：深层嵌套
void process(Data& d) {
    if (d.is_valid()) {
        if (d.has_data()) {
            if (d.data_size() > 0) {
                // 终于开始处理了
                handle(d);
            }
        }
    }
}
\end{lstlisting}

\begin{lstlisting}[language=C++]
// 正确：提前返回
void process(Data& d) {
    if (!d.is_valid()) return;
    if (!d.has_data()) return;
    if (d.data_size() <= 0) return;
    handle(d);
}
\end{lstlisting}

\section*{练习题}

\begin{enumerate}
  \item 将以下函数重构为符合核心指南的函数设计：
\begin{lstlisting}[language=C++]
void do_everything(char* input, int len, char* output,
                   int* out_len, int* error_code) {
    // 100 行代码...
}
\end{lstlisting}
  \item 解释为什么 \texttt{return std::move(local\_variable)} 是一个反模式，以及它如何影响 NRVO。
  \item 设计一个函数，它需要返回三个值（最小值、最大值和平均值）。使用结构体返回值，并展示调用代码。
  \item 编写一个 \texttt{constexpr} 函数来计算字符串的长度（不使用 \texttt{strlen}），并用 \texttt{static\_assert} 验证。
  \item 解释 \texttt{noexcept} 对 \texttt{std::vector} 的移动操作的影响。为什么移动构造函数应该标记为 \texttt{noexcept}？
\end{enumerate}


% ============================================================
% Part III: 类型系统
% ============================================================

\part{类型系统}

% ------------------------------------------------------------
% Chapter 5: 类与类层次
% ------------------------------------------------------------
\chapter{类与类层次}

类是 C++ 中组织数据和行为的基本单位。核心指南将类相关的规则分为几个类别：具体类型（concrete types）、类层次（class hierarchy）和运算符（operators）。本章按照这些类别来组织讨论。

\section{具体类型}

具体类型是最基本、最常用的类——它们不是基类，不用于多态，而是直接实例化使用。\texttt{std::vector}、\texttt{std::string}、\texttt{std::complex} 都是具体类型的例子。

\subsection{规则 C.1：用类来组织相关数据}

\textbf{规则 C.1:} \textit{将相关的数据组织到一个类中（而非使用裸的结构体或全局变量）。}

\begin{lstlisting}[language=C++]
// 不好：分散的数据
std::string name;
int age;
std::string email;
\end{lstlisting}

\begin{lstlisting}[language=C++]
// 好：组织到类中
class Person {
public:
    std::string name;
    int age;
    std::string email;
};
\end{lstlisting}

\subsection{规则 C.2：使用类来保证不变量}

\textbf{规则 C.2:} \textit{如果类有不变量，则不变量不应由调用者来维护。}

不变量是类的核心契约。如果不变量的维护依赖于调用者的"自觉"，那么它迟早会被违反。

\begin{lstlisting}[language=C++]
// 不变量可以被外部破坏
struct Date {
    int year, month, day;  // 不变量：month 1-12, day 1-31
};

Date d;
d.month = 13;  // 不变量被破坏！
\end{lstlisting}

\begin{lstlisting}[language=C++]
// 不变量由类自己维护
class Date {
    int year_, month_, day_;
public:
    Date(int y, int m, int d) {
        if (!is_valid(y, m, d))
            throw std::invalid_argument("invalid date");
        year_ = y; month_ = m; day_ = d;
    }
    void set_month(int m) {
        if (m < 1 || m > 12)
            throw std::invalid_argument("invalid month");
        month_ = m;
    }
    int month() const { return month_; }
    // ...
private:
    static bool is_valid(int y, int m, int d);
};
\end{lstlisting}

\subsection{规则 C.3--C.5：成员访问与默认操作}

\textbf{规则 C.3:} \textit{默认将成员声明为 private。}

将成员声明为 private 可以强制通过成员函数来访问，从而保证不变量。只在确实需要公开访问时才使用 \texttt{public}。对于简单的聚合体（如 \texttt{struct Point \{ int x, y; \}}），使用 public 成员是可以接受的。

\textbf{规则 C.4：使类的对象可以被安全地销毁。}

\textbf{规则 C.5：将辅助函数放在同一个命名空间中。}

\subsection{规则 C.6--C.10：默认操作与构造}

\textbf{规则 C.6:} \textit{遵循零规则——如果不需要自定义析构，就不要自定义。}

零规则是现代 C++ 最重要的设计原则之一。如果类的所有成员都能正确管理自己的资源（使用标准容器、智能指针等），那么类本身不需要自定义任何特殊成员函数。

\begin{lstlisting}[language=C++]
// 零规则——编译器生成的特殊成员函数就足够了
class Person {
    std::string name_;
    std::vector<Address> addresses_;
    std::unique_ptr<Profile> profile_;
    // 编译器自动生成的析构、拷贝、移动操作都是正确的
};
\end{lstlisting}

\textbf{规则 C.7：不要定义拷贝或移动操作，除非类确实需要。}

\textbf{规则 C.8：定义析构函数当且仅当类需要清理资源。}

\textbf{规则 C.9：最小化暴露——优先使用 private 而非 protected。}

\texttt{protected} 意味着派生类可以直接访问基类的内部实现，这增加了耦合。优先使用 private 成员加 protected 访问函数。

\textbf{规则 C.10：优先使用 \texttt{constexpr} 构造函数。}

\begin{lstlisting}[language=C++]
class Point {
    double x_, y_;
public:
    constexpr Point(double x, double y) : x_(x), y_(y) {}
    constexpr double x() const { return x_; }
    constexpr double y() const { return y_; }
};

constexpr Point origin{0.0, 0.0};  // 编译时构造
static_assert(origin.x() == 0.0);
\end{lstlisting}

\section{类层次}

类层次是 C++ 实现运行时多态的机制。核心指南对类层次的使用持谨慎态度——过度使用继承是许多 C++ 代码库中的常见问题。

\subsection{规则 C.35：有虚函数的基类必须有虚析构函数}

\textbf{规则 C.35:} \textit{有虚函数的类必须有虚析构函数。}

这是 C++ 中最经典的规则之一。如果通过基类指针删除派生类对象，而基类的析构函数不是虚函数，行为是未定义的。

\begin{lstlisting}[language=C++]
// 严重错误
class Base {
public:
    virtual void do_something() = 0;
    // 没有虚析构函数！
};

class Derived : public Base {
    std::vector<int> data_;
public:
    void do_something() override { /* ... */ }
};

std::unique_ptr<Base> p = std::make_unique<Derived>();
// 删除时只调用 Base 的析构函数——data_ 泄漏！
\end{lstlisting}

\begin{lstlisting}[language=C++]
// 正确
class Base {
public:
    virtual ~Base() = default;  // 虚析构函数
    virtual void do_something() = 0;
};
\end{lstlisting}

\subsection{规则 C.36--C.39：基类设计}

\textbf{规则 C.36:} \textit{基类析构函数要么是 public 且 virtual 的，要么是 protected 且非 virtual 的。}

如果析构函数是 protected 的，那么对象只能通过派生类来销毁，此时不需要 virtual。

\textbf{规则 C.37:} \textit{将基类设计为接口或具体类（不要混合）。}

\textbf{规则 C.38:} \textit{接口类不应有数据成员。}

\textbf{规则 C.39:} \textit{接口类不应有非纯虚函数。}

\begin{lstlisting}[language=C++]
// 好的接口类
class IDrawable {
public:
    virtual ~IDrawable() = default;
    virtual void draw(Canvas& c) const = 0;
    virtual BoundingBox bounds() const = 0;
    // 没有数据成员，没有非纯虚函数
};

// 好的具体基类（提供默认实现）
class Shape {
public:
    virtual ~Shape() = default;
    virtual void draw(Canvas& c) const = 0;
    // 提供通用的默认实现
    virtual Color color() const { return Color::black; }
private:
    Color color_;  // 具体基类可以有数据成员
};
\end{lstlisting}

\subsection{规则 C.40--C.45：构造函数与继承}

\textbf{规则 C.40:} \textit{如果类有虚函数，它就不应该有非虚的拷贝构造或拷贝赋值运算符。}

\textbf{规则 C.41:} \textit{构造函数应该创建完全初始化的对象。}

\textbf{规则 C.42:} \textit{不要在构造函数中调用虚函数。}

在构造函数中调用虚函数时，派生类的部分尚未构造，因此调用不会到达派生类的重写版本——这是一个常见的陷阱。

\begin{lstlisting}[language=C++]
class Base {
public:
    Base() {
        init();  // 调用的是 Base::init()，不是 Derived::init()！
    }
    virtual void init() { /* Base 的初始化 */ }
};

class Derived : public Base {
    std::vector<int> data_;
public:
    void init() override {
        data_.push_back(42);  // 不会被调用！
    }
};
\end{lstlisting}

\textbf{规则 C.43--C.45:} 关于默认构造、拷贝和移动操作的进一步规则。具体类型应该有明确的拷贝和移动语义，避免"部分移动"的状态。

\subsection{规则 C.50--C.55：工厂模式与多态}

\textbf{规则 C.50:} \textit{使用"虚函数构造器"（工厂模式）当需要多态地创建对象时。}

\begin{lstlisting}[language=C++]
class Shape {
public:
    virtual ~Shape() = default;
    virtual std::unique_ptr<Shape> clone() const = 0;
    virtual void draw() const = 0;
};

class Circle : public Shape {
    double radius_;
public:
    explicit Circle(double r) : radius_(r) {}
    std::unique_ptr<Shape> clone() const override {
        return std::make_unique<Circle>(*this);
    }
    void draw() const override { /* ... */ }
};

// 使用
std::vector<std::unique_ptr<Shape>> shapes;
shapes.push_back(std::make_unique<Circle>(5.0));
shapes.push_back(shapes.back()->clone());  // 多态拷贝
\end{lstlisting}

\section{运算符}

\subsection{规则 C.60：二元运算符应为非成员函数}

\textbf{规则 C.60:} \textit{使二元运算符为非成员函数。}

将二元运算符（如 \texttt{operator+}）定义为非成员函数可以确保两侧的操作数都能进行隐式转换。

\begin{lstlisting}[language=C++]
class Complex {
    double re_, im_;
public:
    Complex(double re = 0, double im = 0) : re_(re), im_(im) {}

    // 作为成员函数——左侧不能隐式转换
    // Complex operator+(const Complex& rhs) const;

    // 友元非成员函数——两侧都可以隐式转换
    friend Complex operator+(const Complex& a, const Complex& b) {
        return Complex(a.re_ + b.re_, a.im_ + b.im_);
    }

    Complex& operator+=(const Complex& rhs) {
        re_ += rhs.re_;
        im_ += rhs.im_;
        return *this;
    }
};

// 现在两种调用都可以
Complex c = Complex(1, 2) + Complex(3, 4);
Complex c2 = 1.0 + Complex(3, 4);  // 1.0 隐式转换为 Complex
\end{lstlisting}

\subsection{规则 C.61--C.65：惯用运算符}

\textbf{规则 C.61:} \textit{拷贝操作应该是语义上的拷贝。}

\textbf{规则 C.62:} \textit{使拷贝赋值运算符可以安全地自赋值。}

\textbf{规则 C.63--C.65:} 关于移动赋值、赋值运算符的返回值等规则。核心原则是：运算符应该遵循最小惊讶原则——\texttt{=} 应该做拷贝，\texttt{+} 应该做加法，等等。

\subsection{规则 C.70--C.75：转换运算符}

\textbf{规则 C.70:} \textit{避免复杂的转换运算符。使用命名函数替代隐式转换。}

\begin{lstlisting}[language=C++]
// 危险：隐式转换
class Money {
    double amount_;
public:
    operator double() const { return amount_; }  // 隐式转换！
};

Money m{100.0};
double d = m + 50.0;  // 意外地将 Money 转换为 double
\end{lstlisting}

\begin{lstlisting}[language=C++]
// 安全：显式转换
class Money {
    double amount_;
public:
    explicit operator double() const { return amount_; }
    double to_double() const { return amount_; }  // 命名函数更清晰
};

Money m{100.0};
double d = static_cast<double>(m) + 50.0;  // 显式转换
double d2 = m.to_double() + 50.0;          // 更清晰
\end{lstlisting}

\subsection{规则 C.80--C.85：继承的替代方案}

\textbf{规则 C.80:} \textit{如果只是想用"类似继承"的方式来复用实现，使用组合而非私有继承。}

\textbf{规则 C.81--C.85:} 关于何时使用继承、何时使用组合、何时使用模板的讨论。核心原则是：继承表示"is-a"关系，组合表示"has-a"关系。不要为了代码复用而滥用继承。

\section*{练习题}

\begin{enumerate}
  \item 设计一个 \texttt{Matrix} 类，遵循零规则（使用 \texttt{std::vector} 存储数据）。实现加法、减法和矩阵乘法运算符。
  \item 解释为什么"在构造函数中调用虚函数"是一个陷阱。给出一个具体的例子来说明问题。
  \item 设计一个图形编辑器中的类层次结构，包括 \texttt{Shape}（抽象基类）、\texttt{Circle}、\texttt{Rectangle} 和 \texttt{Triangle}。使用工厂模式创建对象。
  \item 实现一个 \texttt{Temperature} 类，支持摄氏度、华氏度和开尔文之间的转换。使用运算符重载使得 \texttt{Temperature} 对象可以进行比较和算术运算。
  \item 讨论"优先使用组合而非继承"这条原则。给出一个应该使用组合但被错误地使用继承的例子。
\end{enumerate}


% ------------------------------------------------------------
% Chapter 6: 枚举
% ------------------------------------------------------------
\chapter{枚举}

枚举类型用于表示一组相关的命名常量。C++ 提供了两种枚举：传统的无作用域枚举（\texttt{enum}）和 C++11 引入的有作用域枚举（\texttt{enum class}）。核心指南强烈推荐使用后者。

\section{规则 Enum.1：优先使用枚举而非宏}

\textbf{规则 Enum.1:} \textit{优先使用枚举而非宏来定义命名常量。}

\subsection{为什么重要}

宏常量没有类型信息，不参与作用域规则，并且难以调试。枚举提供了类型安全和作用域控制。

\subsection{错误示例}

\begin{lstlisting}[language=C++]
// 使用宏——没有类型安全
#define RED   0
#define GREEN 1
#define BLUE  2
// 宏是全局的，可能与其他宏冲突
// 没有类型检查：int x = RED + 100; 可以编译
\end{lstlisting}

\subsection{正确示例}

\begin{lstlisting}[language=C++]
// 使用枚举——有类型和作用域
enum class Color { red, green, blue };
Color c = Color::red;
// int x = c + 100;  // 编译错误——类型安全！
\end{lstlisting}

\section{规则 Enum.2：使用枚举表示相关命名常量集合}

\textbf{规则 Enum.2:} \textit{使用枚举来表示一组相关的命名常量。}

枚举将相关的常量组织在一起，使得代码更加清晰和可维护。

\begin{lstlisting}[language=C++]
// 不好：分散的常量
const int STATUS_OK = 0;
const int STATUS_ERROR = 1;
const int STATUS_PENDING = 2;
\end{lstlisting}

\begin{lstlisting}[language=C++]
// 好：使用枚举组织
enum class Status { ok, error, pending };
\end{lstlisting}

\section{规则 Enum.3：优先使用 enum class}

\textbf{规则 Enum.3:} \textit{优先使用 \texttt{enum class} 而非"普通"的 \texttt{enum}。}

\subsection{为什么重要}

传统的 \texttt{enum} 有三个主要问题：
\begin{enumerate}
  \item 枚举值泄漏到外层作用域，可能导致命名冲突。
  \item 枚举值可以隐式转换为整数，破坏了类型安全。
  \item 底层类型不明确。
\end{enumerate}

\texttt{enum class} 解决了所有这三个问题。

\subsection{错误示例}

\begin{lstlisting}[language=C++]
// 传统 enum——问题多多
enum Color { red, green, blue };
enum TrafficLight { red, yellow, green };  // 编译错误！red 和 green 重复

Color c = red;
int x = c;  // 隐式转换为 int——类型不安全
\end{lstlisting}

\subsection{正确示例}

\begin{lstlisting}[language=C++]
// enum class——类型安全
enum class Color { red, green, blue };
enum class TrafficLight { red, yellow, green };  // 没有问题

Color c = Color::red;
TrafficLight t = TrafficLight::red;
// int x = c;  // 编译错误——需要显式转换
int x = static_cast<int>(c);  // 显式转换——意图明确
\end{lstlisting}

\section{规则 Enum.4：为枚举定义操作}

\textbf{规则 Enum.4:} \textit{为枚举定义相关操作。}

\texttt{enum class} 不会自动提供递增、位运算等操作。如果需要这些操作，应该显式定义。

\begin{lstlisting}[language=C++]
enum class Day { mon, tue, wed, thu, fri, sat, sun };

// 定义递增操作
Day& operator++(Day& d) {
    d = (d == Day::sun) ? Day::mon
                        : static_cast<Day>(static_cast<int>(d) + 1);
    return d;
}

// 定义输出操作
std::ostream& operator<<(std::ostream& os, Day d) {
    switch (d) {
        case Day::mon: return os << "Monday";
        case Day::tue: return os << "Tuesday";
        case Day::wed: return os << "Wednesday";
        case Day::thu: return os << "Thursday";
        case Day::fri: return os << "Friday";
        case Day::sat: return os << "Saturday";
        case Day::sun: return os << "Sunday";
    }
    return os;
}
\end{lstlisting}

对于位标志枚举：

\begin{lstlisting}[language=C++]
enum class Permission : unsigned {
    none    = 0,
    read    = 1 << 0,
    write   = 1 << 1,
    execute = 1 << 2
};

Permission operator|(Permission a, Permission b) {
    return static_cast<Permission>(
        static_cast<unsigned>(a) | static_cast<unsigned>(b));
}

Permission operator&(Permission a, Permission b) {
    return static_cast<Permission>(
        static_cast<unsigned>(a) & static_cast<unsigned>(b));
}

Permission p = Permission::read | Permission::write;
bool can_read = (p & Permission::read) != Permission::none;
\end{lstlisting}

\section{规则 Enum.5：不要用全大写命名枚举值}

\textbf{规则 Enum.5:} \textit{不要对枚举值使用全大写命名。}

\subsection{为什么重要}

全大写命名是宏的惯例。使用全大写的枚举值容易与宏冲突，并且在代码中视觉上过于突出。由于 \texttt{enum class} 的值需要带作用域前缀（如 \texttt{Color::red}），使用小写或驼峰命名更加自然。

\begin{lstlisting}[language=C++]
// 不好：全大写——像宏
enum class Color { RED, GREEN, BLUE };
auto c = Color::RED;  // 视觉上太突出
\end{lstlisting}

\begin{lstlisting}[language=C++]
// 好：小写——与 enum class 的作用域配合自然
enum class Color { red, green, blue };
auto c = Color::red;  // 清晰自然
\end{lstlisting}

\section*{练习题}

\begin{enumerate}
  \item 将以下宏常量转换为 \texttt{enum class}：
\begin{lstlisting}[language=C++]
#define HTTP_OK 200
#define HTTP_NOT_FOUND 404
#define HTTP_SERVER_ERROR 500
\end{lstlisting}
  \item 为一个表示"方向"的 \texttt{enum class Direction \{ north, south, east, west \}} 定义递增运算符（顺时针旋转 90 度）和反向运算符（旋转 180 度）。
  \item 设计一个位标志枚举来表示文件权限（读、写、执行），并定义必要的位运算符。
  \item 解释为什么 \texttt{enum class} 比传统 \texttt{enum} 更安全。列出至少三个具体的安全问题。
  \item 为一个表示"月份"的 \texttt{enum class Month} 定义输出运算符和"下一个月份"运算符。
\end{enumerate}


% ------------------------------------------------------------
% Chapter 7: 常量与不可变性
% ------------------------------------------------------------
\chapter{常量与不可变性}

在 C++ 中，\texttt{const} 和 \texttt{constexpr} 是实现不可变性和编译时计算的关键工具。核心指南对 \texttt{const} 的使用持积极态度——默认使用 \texttt{const} 是写出安全、清晰代码的重要手段。

\section{规则 Con.1：默认使对象不可变}

\textbf{规则 Con.1:} \textit{默认情况下，使对象不可变（\texttt{const}）。}

\subsection{为什么重要}

不可变对象更容易推理——你不需要追踪它的修改历史。不可变对象天然是线程安全的。将变量默认声明为 \texttt{const} 可以防止意外的修改，并使代码意图更加清晰。

\subsection{错误示例}

\begin{lstlisting}[language=C++]
// 默认可变——容易意外修改
void process(const std::vector<int>& data) {
    int total = 0;
    for (int i = 0; i < data.size(); ++i) {
        total += data[i];
        data[i] = 0;  // 意外修改了输入数据！
    }
}
\end{lstlisting}

\subsection{正确示例}

\begin{lstlisting}[language=C++]
// 默认 const——防止意外修改
void process(const std::vector<int>& data) {
    int total = 0;
    for (const auto& elem : data) {
        total += elem;
        // elem = 0;  // 编译错误——const 引用不可修改
    }
}

// 局部变量也默认 const
void calculate() {
    const double pi = 3.14159265358979;
    const auto result = compute(pi);
    // pi = 3.0;  // 编译错误——防止意外修改
}
\end{lstlisting}

\subsection{要点}
养成在声明变量时首先写 \texttt{const} 的习惯。只有当确实需要修改变量时，才去掉 \texttt{const}。这个简单的习惯可以消除一大类 bug。

\section{规则 Con.2：默认使成员函数为 const}

\textbf{规则 Con.2:} \textit{默认情况下，使成员函数为 \texttt{const}。}

\subsection{为什么重要}

\texttt{const} 成员函数承诺不修改对象的状态。这个承诺使得：
\begin{itemize}
  \item \texttt{const} 对象可以调用该函数。
  \item 读者可以立即知道该函数不会改变对象状态。
  \item 编译器可以强制这个承诺。
\end{itemize}

\subsection{错误示例}

\begin{lstlisting}[language=C++]
class Widget {
    int value_;
public:
    int get_value() { return value_; }  // 不是 const！
    // const Widget 对象无法调用 get_value()
};

void print(const Widget& w) {
    // w.get_value();  // 编译错误！
}
\end{lstlisting}

\subsection{正确示例}

\begin{lstlisting}[language=C++]
class Widget {
    int value_;
public:
    int get_value() const { return value_; }  // const 成员函数
};

void print(const Widget& w) {
    std::cout << w.get_value();  // OK
}
\end{lstlisting}

\subsection{要点}
写完一个成员函数后，问自己："这个函数需要修改对象的状态吗？"如果答案是否定的，就加上 \texttt{const}。如果你不确定，先加上 \texttt{const}——如果编译器报错说需要修改某个成员，再考虑将该成员声明为 \texttt{mutable} 或去掉函数的 \texttt{const}。

\section{规则 Con.3：默认传递 const 指针和引用}

\textbf{规则 Con.3:} \textit{默认情况下，传递指向 \texttt{const} 的指针和引用。}

\subsection{为什么重要}

如果函数不需要修改参数，就应该使用 \texttt{const} 引用或 \texttt{const} 指针。这向调用者保证他们的数据不会被修改，并且允许函数接受临时对象和 \texttt{const} 对象的参数。

\subsection{错误示例}

\begin{lstlisting}[language=C++]
// 非 const 引用——调用者不知道数据是否会被修改
double compute_average(std::vector<int>& data);

std::vector<int> get_data();
double avg = compute_average(get_data());  // 编译错误！临时对象不能绑定到非 const 引用
\end{lstlisting}

\subsection{正确示例}

\begin{lstlisting}[language=C++]
// const 引用——明确不修改
double compute_average(const std::vector<int>& data);

double avg = compute_average(get_data());  // OK
\end{lstlisting}

\section{规则 Con.4：使用 const 定义构造后不变的值}

\textbf{规则 Con.4:} \textit{使用 \texttt{const} 来定义在构造后值不改变的对象。}

\subsection{为什么重要}

许多变量在初始化之后就不应该再被修改。使用 \texttt{const} 可以将这个意图编码到类型系统中。

\subsection{错误示例}

\begin{lstlisting}[language=C++]
// 可以意外修改
int max_connections = 100;
// ... 200 行代码之后 ...
max_connections = 50;  // 意外修改！
\end{lstlisting}

\subsection{正确示例}

\begin{lstlisting}[language=C++]
// 不可修改
const int max_connections = 100;
// max_connections = 50;  // 编译错误
\end{lstlisting}

\section{规则 Con.5：使用 constexpr 表示编译时可计算的值}

\textbf{规则 Con.5:} \textit{使用 \texttt{constexpr} 来定义可以在编译时计算的值。}

\subsection{为什么重要}

\texttt{constexpr} 变量在编译时就已经有了确定的值，这意味着：
\begin{itemize}
  \item 零运行时开销。
  \item 可以用于需要编译时常量的上下文（如数组大小、模板参数、\texttt{case} 标签）。
  \item 编译器可以进行更积极的优化。
\end{itemize}

\subsection{错误示例}

\begin{lstlisting}[language=C++]
// 运行时计算——不必要的开销
double circle_area(double radius) {
    double pi = 3.14159265358979;  // 每次调用都"初始化"
    return pi * radius * radius;
}
\end{lstlisting}

\subsection{正确示例}

\begin{lstlisting}[language=C++]
// 编译时计算
constexpr double pi = 3.14159265358979;

constexpr double circle_area(double radius) {
    return pi * radius * radius;
}

constexpr double area = circle_area(5.0);  // 编译时计算
static_assert(circle_area(1.0) > 3.0);     // 编译时验证

// 可以用于需要编译时常量的上下文
constexpr int table_size = 256;
std::array<int, table_size> lookup_table;
\end{lstlisting}

\subsection{constexpr 与 const 的区别}

理解 \texttt{const} 和 \texttt{constexpr} 的区别很重要：

\begin{itemize}
  \item \texttt{const} 表示"运行时不可修改"——值在运行时确定后就不能改变。
  \item \texttt{constexpr} 表示"编译时可计算"——值在编译时就已经确定。
\end{itemize}

\begin{lstlisting}[language=C++]
const int runtime_const = get_value();    // 运行时确定，之后不可修改
constexpr int compile_const = 42;         // 编译时确定

// constexpr 函数可以在运行时调用
void process(int n) {
    constexpr int block_size = 256;  // 编译时常量
    int chunk_count = n / block_size;
}
\end{lstlisting}

\section*{练习题}

\begin{enumerate}
  \item 在以下代码中添加所有可能的 \texttt{const}，并解释每处添加的理由：
\begin{lstlisting}[language=C++]
class StringProcessor {
    std::string prefix_;
public:
    StringProcessor(std::string prefix) : prefix_(prefix) {}
    std::string process(std::string input) {
        std::string result = prefix_ + input;
        std::transform(result.begin(), result.end(),
                       result.begin(), ::toupper);
        return result;
    }
    std::string get_prefix() { return prefix_; }
};
\end{lstlisting}
  \item 编写一个 \texttt{constexpr} 函数来计算阶乘，并验证 \texttt{factorial(12) == 479001600}。
  \item 解释为什么成员函数默认应该是 \texttt{const} 的。\texttt{mutable} 关键字在什么情况下是合理的？
  \item 设计一个 \texttt{constexpr} 的 \texttt{Point2D} 类，支持距离计算（使用 \texttt{constexpr} 的平方根近似）。
  \item 讨论"默认 const"原则在多线程编程中的优势。
\end{enumerate}


% ------------------------------------------------------------
% Chapter 8: 模板与泛型编程
% ------------------------------------------------------------
\chapter{模板与泛型编程}

模板是 C++ 泛型编程的基础。通过模板，我们可以编写不依赖于具体类型的通用代码。C++20 引入的概念（Concepts）极大地改善了模板编程的体验——更清晰的接口、更好的错误信息、更强的表达力。

\section{模板抽象与概念基础}

\subsection{规则 T.1：使用模板来提高代码的抽象层次}

\textbf{规则 T.1:} \textit{使用模板来提高代码的抽象层次。}

\subsection{为什么重要}

模板使得我们可以编写一次代码而用于多种类型。这不仅减少了代码重复，更重要的是它提高了代码的抽象层次——我们可以讨论"任意可排序的类型"而非"int 的排序"和"string 的排序"。

\subsection{错误示例}

\begin{lstlisting}[language=C++]
// 为每种类型重复代码
int max_int(int a, int b) { return a > b ? a : b; }
double max_double(double a, double b) { return a > b ? a : b; }
std::string max_string(const std::string& a, const std::string& b) {
    return a > b ? a : b;
}
\end{lstlisting}

\subsection{正确示例}

\begin{lstlisting}[language=C++]
// 一个模板适用于所有可比较类型
template<typename T>
T max_val(T a, T b) {
    return a > b ? a : b;
}

auto m1 = max_val(3, 5);
auto m2 = max_val(3.14, 2.71);
auto m3 = max_val(std::string("abc"), std::string("xyz"));
\end{lstlisting}

\subsection{规则 T.2：使用概念（Concepts）来指定模板参数的要求}

\textbf{规则 T.2:} \textit{使用概念（Concepts）来指定模板参数的要求。}

概念是 C++20 最重要的特性之一。它为模板参数提供了命名化的约束集合，使得模板接口更加清晰，错误信息更加可读。

\begin{lstlisting}[language=C++]
// 没有概念——错误信息难以理解
template<typename T>
T add(T a, T b) {
    return a + b;
}
// add(std::vector<int>{}, std::vector<int>{});
// 产生数十行难以理解的错误信息
\end{lstlisting}

\begin{lstlisting}[language=C++]
// 使用概念——错误信息清晰
template<std::regular T>
T add(T a, T b) {
    return a + b;
}
// add(std::vector<int>{}, std::vector<int>{});
// 错误：vector<int> 不满足 regular 概念

// 自定义概念
template<typename T>
concept Numeric = std::is_arithmetic_v<T>;

template<Numeric T>
T safe_divide(T a, T b) {
    if (b == T{0}) throw std::domain_error("division by zero");
    return a / b;
}
\end{lstlisting}

\subsection{规则 T.3--T.10：类型要求与概念设计}

\textbf{规则 T.3:} \textit{使用概念来指定类型的精确要求（而非使用 \texttt{requires} 子句或 \texttt{static\_assert}）。}

虽然 \texttt{requires} 子句和 \texttt{static\_assert} 也可以约束模板参数，但概念更加可重用、可组合，并且提供更好的错误信息。

\textbf{规则 T.10：为所有模板参数指定概念。}

没有概念约束的模板参数可以接受任何类型，这使得接口文档不完整，错误信息难以理解。

\begin{lstlisting}[language=C++]
// 不好：无约束的模板参数
template<typename T, typename U>
auto combine(T t, U u) {
    return t + u;  // T 和 U 需要什么操作？不清楚
}
\end{lstlisting}

\begin{lstlisting}[language=C++]
// 好：有概念约束
template<std::regular T, std::regular U>
    requires std::convertible_to<T, U> || std::convertible_to<U, T>
auto combine(T t, U u) {
    return t + u;
}
\end{lstlisting}

\section{模板工厂与类型函数}

\subsection{规则 T.40：使用工厂函数来创建多态对象}

\textbf{规则 T.40:} \textit{使用工厂函数来创建需要多态行为的对象。}

工厂函数可以封装对象的创建细节，返回基类指针或智能指针，使得调用者不需要知道具体的派生类。

\begin{lstlisting}[language=C++]
// 工厂函数模板
template<typename Derived, typename... Args>
std::unique_ptr<Base> make_shape(Args&&... args) {
    return std::make_unique<Derived>(std::forward<Args>(args)...);
}

auto circle = make_shape<Circle>(5.0);
auto rect = make_shape<Rectangle>(3.0, 4.0);
\end{lstlisting}

\subsection{规则 T.41--T.45：类型函数}

\textbf{规则 T.41:} \textit{使用类型特征（type traits）来查询类型的属性。}

\textbf{规则 T.42:} \textit{使用别名模板（alias templates）来简化类型的使用。}

\begin{lstlisting}[language=C++]
// 类型特征
static_assert(std::is_integral_v<int>);
static_assert(std::is_base_of_v<Base, Derived>);
static_assert(std::is_copy_constructible_v<Widget>);

// 别名模板
template<typename T>
using Vec = std::vector<T>;

template<typename T>
using Unique = std::unique_ptr<T>;

Vec<int> numbers;  // 比 std::vector<int> 更简洁
Unique<Widget> w = std::make_unique<Widget>();
\end{lstlisting}

\section{模板命名与概念定义}

\subsection{规则 T.80--T.85：模板命名与声明}

\textbf{规则 T.80:} \textit{对于泛型代码，优先使用概念约束的模板而非 \texttt{requires} 子句。}

\textbf{规则 T.81--T.83:} 关于概念的定义和组合。概念应该遵循命名惯例（通常使用小写或描述性的名称），应该原子化（每个概念检查一个语义属性），并且应该可组合（通过 \texttt{\&\&} 和 \texttt{||} 组合）。

\begin{lstlisting}[language=C++]
// 原子概念
template<typename T>
concept Incrementable = requires(T a) {
    { ++a } -> std::same_as<T&>;
    { a++ } -> std::same_as<T>;
};

template<typename T>
concept Decrementable = requires(T a) {
    { --a } -> std::same_as<T&>;
    { a-- } -> std::same_as<T>;
};

// 组合概念
template<typename T>
concept BidirectionalIterator = Incrementable<T> && Decrementable<T>;
\end{lstlisting}

\subsection{规则 T.100--T.110：概念的设计原则}

\textbf{规则 T.100:} \textit{使用概念来表达语义类别（而非仅仅是语法要求）。}

好的概念不仅描述类型"能做什么"（语法），还描述类型"代表什么"（语义）。例如，\texttt{std::regular} 不仅要求类型可以拷贝和比较，还要求这些操作具有"值语义"——拷贝后的对象与原对象等价。

\begin{lstlisting}[language=C++]
// 仅语法要求——不够
template<typename T>
concept HasPlus = requires(T a, T b) {
    { a + b } -> std::same_as<T>;
};

// 语义要求——更好
// std::regular 包含了：
// - 可默认构造
// - 可拷贝构造和赋值
// - 可比较（== 和 !=）
// - 可交换（swap）
// - 这些操作具有正确的语义
template<std::regular T>
T add(T a, T b) { return a + b; }
\end{lstlisting}

\subsection{规则 T.120--T.140：模板的高级使用}

\textbf{规则 T.120:} \textit{当需要编译时多态时使用模板。}

模板提供编译时多态（静态多态），虚函数提供运行时多态（动态多态）。当类型在编译时已知且性能至关重要时，使用模板。

\begin{lstlisting}[language=C++]
// 编译时多态——零开销
template<typename Container>
auto sum(const Container& c) {
    typename Container::value_type total{};
    for (const auto& elem : c) {
        total += elem;
    }
    return total;
}

// 运行时多态——有虚函数开销
// 当需要存储不同类型的对象时使用
\end{lstlisting}

\textbf{规则 T.140：如果函数不需要被重写，不要使用虚函数。}

如果一个操作不需要在运行时改变行为，就不应该使用虚函数。模板或概念可以提供更好的性能。

\subsection{变量模板与 constexpr if}

C++14 的变量模板和 C++17 的 \texttt{constexpr if} 进一步增强了模板编程的表达力：

\begin{lstlisting}[language=C++]
// 变量模板
template<typename T>
constexpr T pi = T(3.14159265358979323846);

double area_d = pi<double> * 5.0 * 5.0;
float area_f = pi<float> * 5.0f * 5.0f;

// constexpr if——编译时分支
template<typename T>
std::string to_string(const T& value) {
    if constexpr (std::is_arithmetic_v<T>) {
        return std::to_string(value);
    } else if constexpr (requires { value.to_string(); }) {
        return value.to_string();
    } else {
        return "unknown";
    }
}
\end{lstlisting}

\section*{练习题}

\begin{enumerate}
  \item 定义一个概念 \texttt{Printable}，要求类型支持 \texttt{operator<<} 输出到 \texttt{std::ostream}。使用该概念约束一个泛型的 \texttt{print} 函数。
  \item 实现一个泛型的 \texttt{Stack<T>} 容器，使用 \texttt{std::vector<T>} 作为底层存储。要求 \texttt{T} 满足 \texttt{std::regular} 概念。
  \item 使用 \texttt{constexpr if} 实现一个泛型的 \texttt{serialize} 函数，对于算术类型直接写入字节，对于 \texttt{std::string} 先写入长度再写入内容，对于 \texttt{std::vector<T>} 递归序列化每个元素。
  \item 定义概念 \texttt{Container} 来约束标准容器的通用接口（\texttt{begin()}、\texttt{end()}、\texttt{size()}、\texttt{empty()}）。使用该概念编写一个泛型的 \texttt{print\_container} 函数。
  \item 比较编译时多态（模板）和运行时多态（虚函数）的优缺点。在什么场景下应该选择哪种方式？
\end{enumerate}


% ============================================================
% Part IV: 资源管理与错误处理
% ============================================================

\part{资源管理与错误处理}

% ============================================================
% Chapter 9: 资源管理 (Resource Management - "R" section)
% ============================================================

\chapter{资源管理}

资源管理是现代C++编程中最核心的课题之一。所谓"资源"，不仅包括动态分配的内存，还涵盖文件句柄、网络套接字、数据库连接、互斥锁、图形设备上下文等一切需要在获取后最终释放的系统要素。C++ Core Guidelines 的"R"系列规则围绕一个中心思想展开：用语言机制和类型系统来保证资源的安全管理，从而杜绝泄漏、悬空指针和未定义行为。

\section{通用资源管理原则 (R.1--R.10)}

\subsection*{\textbf{规则 R.1:} 资源获取即初始化（RAII）}

RAII（Resource Acquisition Is Initialization）是C++中管理资源的基本技术。其核心思想是：将资源的生命周期绑定到一个局部对象的生命周期。当对象构造时获取资源，当对象离开作用域（无论正常退出还是异常退出）时，析构函数自动释放资源。这样就不需要手动释放资源，也天然地保证了异常安全。

\textbf{错误示例：}

\begin{lstlisting}[language=C++]
void write_data(const string& filename) {
    FILE* fp = fopen(filename.c_str(), "w");
    if (!fp) return;
    // 如果这里抛出异常，fp 永远不会被关闭
    fprintf(fp, "hello\n");
    fclose(fp);  // 如果上面抛异常，此行不会执行
}
\end{lstlisting}

\textbf{正确示例：}

\begin{lstlisting}[language=C++]
void write_data(const string& filename) {
    ofstream fp(filename);
    if (!fp) return;
    // 即使这里抛出异常，ofstream 的析构函数也会关闭文件
    fp << "hello\n";
}  // fp 在这里自动关闭
\end{lstlisting}

\textbf{要点：} RAII 是C++中异常安全的基石。任何需要成对操作（获取/释放、加锁/解锁、打开/关闭）的场景，都应该用 RAII 来管理。标准库中的 \texttt{unique\_ptr}、\texttt{shared\_ptr}、\texttt{lock\_guard}、\texttt{fstream} 等都是 RAII 的典型应用。

\subsection*{\textbf{规则 R.2:} 优先使用 \texttt{unique\_ptr} 而非 \texttt{shared\_ptr}}

当不需要共享所有权时，应当使用 \texttt{unique\_ptr}。\texttt{unique\_ptr} 的开销与裸指针几乎相同——它不维护引用计数，不产生额外的内存分配，且可以被优化为与裸指针完全相同的机器码。而 \texttt{shared\_ptr} 需要维护控制块和原子引用计数，带来不可忽视的运行时开销。

\textbf{错误示例：}

\begin{lstlisting}[language=C++]
// 不需要共享所有权，却使用了 shared_ptr
shared_ptr<Sensor> create_sensor() {
    return make_shared<Sensor>();
}
// 调用者无法知道是否需要共享
\end{lstlisting}

\textbf{正确示例：}

\begin{lstlisting}[language=C++]
// 明确表达独占所有权
unique_ptr<Sensor> create_sensor() {
    return make_unique<Sensor>();
}
// 如果调用者确实需要共享，可以显式转换
shared_ptr<Sensor> sp = create_sensor(); // 隐式转换
\end{lstlisting}

\textbf{要点：} 返回值类型本身就是一种契约。\texttt{unique\_ptr} 告诉调用者"你获得了独占权"，而 \texttt{shared\_ptr} 则意味着"你获得了一份共享权"。选择错误的智能指针类型会误导读者并引入不必要的开销。

\subsection*{\textbf{规则 R.3:} 裸指针是可疑的}

裸指针（raw pointer）本身不表达所有权语义。当你看到一个裸指针时，你无法知道它是否拥有资源、是否需要释放、是否可以删除。因此，裸指针应当仅用于"不拥有资源"的观察场景，其他情况下应当使用智能指针或引用。

\textbf{错误示例：}

\begin{lstlisting}[language=C++]
Widget* create_widget(int n) {
    return new Widget(n);  // 谁负责 delete？
}

void process(Widget* w) {  // w 是拥有还是观察？不清楚
    w->do_something();
}
\end{lstlisting}

\textbf{正确示例：}

\begin{lstlisting}[language=C++]
unique_ptr<Widget> create_widget(int n) {
    return make_unique<Widget>(n);  // 明确：转移所有权
}

void process(Widget& w) {  // 明确：仅观察，不拥有
    w.do_something();
}
\end{lstlisting}

\textbf{要点：} 如果一个函数参数是裸指针，读者必须阅读函数体才能知道是否需要释放。这增加了认知负担并容易出错。用引用或智能指针替代裸指针可以让接口自文档化。

\subsection*{\textbf{规则 R.4:} 裸句柄是可疑的}

与裸指针类似，裸句柄（如文件描述符 \texttt{int fd}、Windows \texttt{HANDLE}、数据库连接句柄等）不表达所有权和生命周期信息。应当将它们封装在 RAII 类中。

\textbf{错误示例：}

\begin{lstlisting}[language=C++]
int fd = open("data.txt", O_RDONLY);
if (fd < 0) { /* 错误处理 */ }
// 如果这里抛出异常，fd 永远不会被关闭
read(fd, buffer, size);
close(fd);
\end{lstlisting}

\textbf{正确示例：}

\begin{lstlisting}[language=C++]
class FileDescriptor {
    int fd_;
public:
    explicit FileDescriptor(const char* path)
        : fd_(open(path, O_RDONLY)) {
        if (fd_ < 0) throw runtime_error("open failed");
    }
    ~FileDescriptor() { if (fd_ >= 0) close(fd_); }
    FileDescriptor(const FileDescriptor&) = delete;
    FileDescriptor& operator=(const FileDescriptor&) = delete;
    int get() const { return fd_; }
};

void read_data() {
    FileDescriptor fd("data.txt");
    read(fd.get(), buffer, size);
}  // 自动关闭
\end{lstlisting}

\textbf{要点：} 任何非类类型的资源句柄都应当被封装。封装不仅保证资源释放，还能防止复制导致的 double-close 问题。

\subsection*{\textbf{规则 R.5--R.10:} 具体的 RAII 模式}

这组规则进一步细化了 RAII 的使用模式：

\textbf{R.5:} 对象的所有权应当是明确的。如果一个对象被分配在自由存储（堆）上，应当用智能指针或容器管理它。

\textbf{R.6:} 避免裸 \texttt{new} 和裸 \texttt{delete} 表达式。它们应当被封装在 RAII 对象内部。

\textbf{R.7:} 避免裸 \texttt{malloc} 和 \texttt{free}。在C++中应当使用 \texttt{new}/\texttt{delete}（通过智能指针间接使用）或标准容器。

\textbf{R.8:} 避免裸 \texttt{new} 和 \texttt{delete} 表达式——即使在成员函数中。

\textbf{R.9:} 不要在裸指针上调用 \texttt{delete}。如果你需要释放资源，说明你不应该使用裸指针。

\textbf{R.10:} 避免 \texttt{new T{}} 和 \texttt{delete p} 的直接使用。

\textbf{错误示例：}

\begin{lstlisting}[language=C++]
class Container {
    Widget* data_;
    size_t size_;
public:
    Container(size_t n) : data_(new Widget[n]), size_(n) {}
    ~Container() { delete[] data_; }
    // 如果忘记拷贝构造/赋值运算符，浅拷贝导致 double free
};
\end{lstlisting}

\textbf{正确示例：}

\begin{lstlisting}[language=C++]
class Container {
    vector<Widget> data_;
public:
    Container(size_t n) : data_(n) {}
    // 拷贝、移动、析构均由 vector 自动处理
};
\end{lstlisting}

\textbf{要点：} 标准容器（\texttt{vector}、\texttt{string}、\texttt{unique\_ptr} 等）本身就是 RAII 对象。优先使用它们可以消除大量手动资源管理代码，同时获得正确的拷贝和移动语义。

\section{所有权与智能指针 (R.20--R.35)}

\subsection*{\textbf{规则 R.20:} 使用 \texttt{unique\_ptr} 或 \texttt{shared\_ptr} 表示所有权}

当函数需要转移动态分配对象的所有权时，返回值应当是智能指针类型。这是向调用者传达所有权语义的最清晰方式。

\textbf{错误示例：}

\begin{lstlisting}[language=C++]
Widget* make_widget() {
    return new Widget();  // 调用者必须记得 delete
}
\end{lstlisting}

\textbf{正确示例：}

\begin{lstlisting}[language=C++]
unique_ptr<Widget> make_widget() {
    return make_unique<Widget>();  // 所有权明确转移
}
\end{lstlisting}

\textbf{要点：} 返回智能指针是C++11之后表达所有权转移的标准做法。裸指针返回值在现代C++中被视为代码异味（code smell）。

\subsection*{\textbf{规则 R.21:} 不需要共享所有权时优先使用 \texttt{unique\_ptr}}

\texttt{unique\_ptr} 的语义简单明确：一个对象只有一个所有者。当所有者销毁时，对象也被销毁。\texttt{shared\_ptr} 的引用计数带来额外开销，且共享所有权可能导致对象生命周期超出预期（逻辑泄漏）。

\textbf{错误示例：}

\begin{lstlisting}[language=C++]
// 不必要的引用计数开销
shared_ptr<Connection> open_connection() {
    return make_shared<Connection>();
}
// 调用者可能误以为需要共享
\end{lstlisting}

\textbf{正确示例：}

\begin{lstlisting}[language=C++]
unique_ptr<Connection> open_connection() {
    return make_unique<Connection>();
}
// 调用者知道：我拥有这个连接
\end{lstlisting}

\textbf{要点：} \texttt{unique\_ptr} 是零开销抽象。它不产生额外的内存分配（不像 \texttt{shared\_ptr} 的控制块），不产生原子操作开销。只有在确实需要共享时才使用 \texttt{shared\_ptr}。

\subsection*{\textbf{规则 R.22:} 使用 \texttt{make\_unique} 创建 \texttt{unique\_ptr}}

\texttt{make\_unique} 比直接 \texttt{new} 再传入 \texttt{unique\_ptr} 构造函数更安全、更高效。它避免了异常安全问题（两个参数求值顺序不确定时，\texttt{new} 可能泄漏），且只需一次内存分配。

\textbf{错误示例：}

\begin{lstlisting}[language=C++]
unique_ptr<Widget> p(new Widget(42));  // 两次分配，异常不安全
\end{lstlisting}

\textbf{正确示例：}

\begin{lstlisting}[language=C++]
auto p = make_unique<Widget>(42);  // 一次分配，异常安全
\end{lstlisting}

\textbf{要点：} \texttt{make\_unique} 是C++14引入的，C++11中可以自行实现。它消除了重复的类型名，减少了代码中的 \texttt{new} 出现次数。

\subsection*{\textbf{规则 R.23:} 使用 \texttt{make\_shared} 创建 \texttt{shared\_ptr}}

\texttt{make\_shared} 将对象和控制块分配在同一段内存中，减少了一次内存分配，提高了缓存局部性。

\textbf{错误示例：}

\begin{lstlisting}[language=C++]
shared_ptr<Widget> p(new Widget(42));  // 两次分配
\end{lstlisting}

\textbf{正确示例：}

\begin{lstlisting}[language=C++]
auto p = make_shared<Widget>(42);  // 一次分配
\end{lstlisting}

\textbf{要点：} \texttt{make\_shared} 的缺点是对象内存不能在其引用计数归零后立即释放（因为控制块和对象在同一分配中），但在绝大多数场景下，性能优势大于这一微小延迟。

\subsection*{\textbf{规则 R.24--R.30:} 更多所有权模式}

\textbf{R.24:} 不要使用 \texttt{std::auto\_ptr}。它已被C++17移除，其拷贝语义（移动而非拷贝）违反直觉。

\textbf{R.25:} 不要使用 \texttt{shared\_ptr} 作为"万能指针"。如果一个对象不需要共享所有权，不要用 \texttt{shared\_ptr}。

\textbf{R.26:} 避免开销巨大的 \texttt{shared\_ptr} 误用。例如，不要用 \texttt{shared\_ptr} 管理不需要动态分配的对象。

\textbf{R.27:} 使用 \texttt{std::owner\_less} 或弱指针打破循环引用。\texttt{shared\_ptr} 的循环引用会导致内存泄漏。

\textbf{R.28:} 避免使用 \texttt{shared\_ptr} 形成循环引用。循环引用中的对象永远不会被释放。

\textbf{R.29:} \texttt{shared\_ptr} 的引用计数是原子操作——注意性能影响。在多线程环境中，每次拷贝和销毁 \texttt{shared\_ptr} 都会产生原子操作的开销。

\textbf{R.30:} 不要让 \texttt{shared\_ptr} 的生命周期管理变得过于复杂。如果所有权关系不清晰，考虑重新设计。

\textbf{错误示例（循环引用）：}

\begin{lstlisting}[language=C++]
struct Node {
    shared_ptr<Node> next;
    shared_ptr<Node> prev;
    ~Node() { cout << "destroyed\n"; }  // 永远不会被调用！
};

auto a = make_shared<Node>();
auto b = make_shared<Node>();
a->next = b;
b->prev = a;
// a 和 b 离开作用域后，引用计数均为1，永远不会归零
\end{lstlisting}

\textbf{正确示例：}

\begin{lstlisting}[language=C++]
struct Node {
    shared_ptr<Node> next;
    weak_ptr<Node> prev;  // 用 weak_ptr 打破循环
    ~Node() { cout << "destroyed\n"; }  // 现在会正确调用
};
\end{lstlisting}

\textbf{要点：} 循环引用是 \texttt{shared\_ptr} 最常见的陷阱。在设计数据结构时，应当从一开始就考虑所有权方向，用 \texttt{weak\_ptr} 表示反向引用。

\subsection*{\textbf{规则 R.31--R.35:} 参数传递与所有权转移}

\textbf{R.31:} 如果你需要从一个函数中移出一个对象，使用 \texttt{unique\_ptr} 作为返回类型。

\textbf{R.32:} 使用 \texttt{unique\_ptr<T\&\&>} 是不合法的——智能指针不能引用局部变量。

\textbf{R.33:} 如果函数需要获取对象的所有权，参数应当使用 \texttt{unique\_ptr}（移动语义）或 \texttt{shared\_ptr}（共享语义）。

\textbf{R.34:} 如果函数不需要所有权，使用裸指针、引用或 \texttt{string\_view} 等观察类型。

\textbf{R.35:} \texttt{shared\_ptr\&} 作为参数通常是设计错误——它意味着函数可能改变调用者的 \texttt{shared\_ptr}。

\textbf{错误示例：}

\begin{lstlisting}[language=C++]
// 不必要的 shared_ptr 拷贝
void process(shared_ptr<Widget>& w) {  // 引用 shared_ptr
    w->do_something();
}
// 调用者困惑：为什么参数是 shared_ptr 的引用？
\end{lstlisting}

\textbf{正确示例：}

\begin{lstlisting}[language=C++]
// 仅观察
void process(Widget& w) {
    w.do_something();
}

// 需要共享所有权
void store(shared_ptr<Widget> w) {  // 值传递，拷贝
    data_ = move(w);
}
\end{lstlisting}

\textbf{要点：} 函数参数的类型就是所有权的文档。\texttt{unique\_ptr} 值参数表示"我将拿走所有权"，\texttt{shared\_ptr} 值参数表示"我将共享所有权"，引用或裸指针表示"我只是观察"。

\section*{练习题}

\begin{enumerate}
\item 将以下函数改写为 RAII 风格，确保异常安全：
\begin{lstlisting}[language=C++]
void copy_file(const char* src, const char* dst) {
    FILE* in = fopen(src, "rb");
    FILE* out = fopen(dst, "wb");
    char buf[4096];
    size_t n;
    while ((n = fread(buf, 1, sizeof(buf), in)) > 0)
        fwrite(buf, 1, n, out);
    fclose(in);
    fclose(out);
}
\end{lstlisting}

\item 解释为什么在以下代码中 \texttt{shared\_ptr} 比 \texttt{unique\_ptr} 更合适，并给出实现：一个观察者模式，其中多个观察者同时观察同一个主题对象，主题对象需要在所有观察者都断开后才能销毁。

\item 找出以下代码中的循环引用问题并修复：
\begin{lstlisting}[language=C++]
struct Task {
    string name;
    shared_ptr<Task> parent;
    vector<shared_ptr<Task>> children;
};
\end{lstlisting}

\item 设计一个 \texttt{SocketGuard} 类，封装 POSIX 套接字描述符，实现 RAII 语义。要求支持移动构造和移动赋值，禁止拷贝。

\item 分析以下代码的性能问题，并给出改进方案：
\begin{lstlisting}[language=C++]
void process(vector<shared_ptr<Event>>& events) {
    for (auto& e : events) {
        shared_ptr<Event> copy = e;  // 不必要的拷贝
        copy->handle();
    }
}
\end{lstlisting}
\end{enumerate}

% ============================================================
% Chapter 10: 错误处理 (Error Handling - "E" section)
% ============================================================

\chapter{错误处理}

错误处理是软件可靠性的基石。C++提供了异常机制作为错误处理的主要手段，但异常的使用需要遵循严格的规则，否则可能引入比返回值检查更多的混乱。Core Guidelines 的"E"系列规则旨在帮助开发者建立一致、安全、高效的错误处理策略。

\section{错误处理策略 (E.1--E.10)}

\subsection*{\textbf{规则 E.1:} 制定错误处理策略并一致地遵循}

一个项目应当有明确的错误处理策略：何时使用异常？何时使用错误码？何时终止程序？没有一致的策略会导致代码中混合使用多种错误处理方式，增加维护难度和出错概率。

\textbf{错误示例：}

\begin{lstlisting}[language=C++]
// 同一个项目中混合使用多种方式
int divide(int a, int b) {
    if (b == 0) return -1;         // 错误码
    return a / b;
}

double divide(double a, double b) {
    if (b == 0.0) throw runtime_error("div by zero");  // 异常
    return a / b;
}

// 调用者不知道应该检查返回值还是捕获异常
\end{lstlisting}

\textbf{正确示例：}

\begin{lstlisting}[language=C++]
// 统一策略：用异常报告错误
double divide(double a, double b) {
    if (b == 0.0) throw invalid_argument("division by zero");
    return a / b;
}

// 调用者知道需要 try-catch
try {
    auto result = divide(1.0, 0.0);
} catch (const invalid_argument& e) {
    cerr << "Error: " << e.what() << endl;
}
\end{lstlisting}

\textbf{要点：} 团队应当在项目初期就制定错误处理策略并写入编码规范。通常的策略是：用异常处理"不应该发生"的错误，用返回值处理"预期内"的错误（如用户输入验证），用 \texttt{assert} 处理"不可能发生"的逻辑错误。

\subsection*{\textbf{规则 E.2:} 抛出对象来表示函数无法完成其任务}

当函数无法完成其声明的任务时，应当抛出一个异常对象。异常对象应当携带足够的信息，让调用者了解错误的性质。

\textbf{错误示例：}

\begin{lstlisting}[language=C++]
void connect_to_server(const string& host) {
    if (connection_failed)
        throw 42;  // 抛出整数，毫无信息
    // ...
    if (timeout)
        throw "timeout";  // 抛出字符串字面量，类型不安全
}
\end{lstlisting}

\textbf{正确示例：}

\begin{lstlisting}[language=C++]
class ConnectionError : public runtime_error {
public:
    explicit ConnectionError(const string& msg)
        : runtime_error(msg) {}
};

void connect_to_server(const string& host) {
    if (connection_failed)
        throw ConnectionError("Cannot connect to " + host);
    // ...
}
\end{lstlisting}

\textbf{要点：} 异常对象应当继承自标准异常类（\texttt{std::exception} 或其子类），这样调用者可以用统一的 \texttt{catch} 块处理，同时获得 \texttt{what()} 方法提供的错误描述。

\subsection*{\textbf{规则 E.3:} 仅将异常用于错误处理}

异常不应当被用于正常的控制流。异常的设计目的是处理"异常"情况——即函数无法完成其正常任务的情况。将异常用于正常流程不仅性能低下（异常抛出和捕获的开销远大于条件分支），还会使代码逻辑难以理解。

\textbf{错误示例：}

\begin{lstlisting}[language=C++]
// 用异常实现正常的查找逻辑
int find(const vector<int>& v, int target) {
    for (size_t i = 0; i < v.size(); ++i) {
        if (v[i] == target) throw FoundAt(i);  // 异常用于正常流程
    }
    throw NotFound();
}
\end{lstlisting}

\textbf{正确示例：}

\begin{lstlisting}[language=C++]
// 使用返回值或可选类型
optional<size_t> find(const vector<int>& v, int target) {
    for (size_t i = 0; i < v.size(); ++i) {
        if (v[i] == target) return i;
    }
    return nullopt;
}
\end{lstlisting}

\textbf{要点：} "异常"顾名思义，应当用于异常情况。如果某个"错误"经常发生（如查找失败），它就不是异常，应当用 \texttt{optional}、\texttt{expected}（C++23）或错误码来处理。

\subsection*{\textbf{规则 E.4:} 围绕目标设计错误处理策略}

错误处理策略应当围绕以下目标设计：(1) 不泄漏资源；(2) 不使程序处于不一致状态；(3) 向调用者提供足够的错误信息；(4) 允许调用者从错误中恢复（如果可能）。

\textbf{要点：} 在设计错误处理策略时，应当问自己：如果这个操作失败，程序的状态是什么？资源是否被正确释放？调用者能否知道发生了什么？调用者能否采取补救措施？

\subsection*{\textbf{规则 E.5:} 让构造函数成为构造函数}

构造函数的职责是建立类的不变量。如果构造函数无法建立不变量，它应当抛出异常。不要试图用"无效状态"标记来表示构造失败——这会导致对象处于不一致状态，后续操作可能产生未定义行为。

\textbf{错误示例：}

\begin{lstlisting}[language=C++]
class Connection {
    bool valid_ = false;
public:
    Connection(const string& host) {
        if (!connect(host)) {
            valid_ = false;  // 构造"成功"但对象无效
            return;
        }
        valid_ = true;
    }
    void send(const string& data) {
        if (!valid_) { /* 未定义行为或需要每次检查 */ }
        // ...
    }
};
\end{lstlisting}

\textbf{正确示例：}

\begin{lstlisting}[language=C++]
class Connection {
public:
    explicit Connection(const string& host) {
        if (!connect(host))
            throw ConnectionError("Failed to connect to " + host);
        // 如果构造成功，不变量已建立
    }
    void send(const string& data) {
        // 不需要检查 valid_，因为对象一定有效
    }
};
\end{lstlisting}

\textbf{要点：} 构造函数抛出异常意味着对象从未存在过。这比返回一个"无效对象"安全得多，因为调用者不需要在每个操作前检查对象是否有效。

\subsection*{\textbf{规则 E.6:} 使用 RAII 防止泄漏}

这与规则 R.1 相呼应。在错误处理的语境下，RAII 的重要性更加突出：当异常被抛出时，栈展开（stack unwinding）会销毁所有已构造的局部对象。如果资源没有被 RAII 对象管理，这些资源将永久泄漏。

\textbf{错误示例：}

\begin{lstlisting}[language=C++]
void process() {
    int* buffer = new int[1000];
    // ... 可能抛出异常
    delete[] buffer;  // 异常时不会执行
}
\end{lstlisting}

\textbf{正确示例：}

\begin{lstlisting}[language=C++]
void process() {
    vector<int> buffer(1000);
    // ... 即使抛出异常，buffer 的析构函数也会释放内存
}
\end{lstlisting}

\textbf{要点：} RAII 是C++中实现异常安全的最重要工具。遵循"零泄漏"原则：任何代码路径（包括异常路径）都不应泄漏资源。

\subsection*{\textbf{规则 E.7:} 声明前置条件}

前置条件是调用者在调用函数之前必须满足的条件。如果前置条件不满足，函数行为未定义。前置条件应当被明确记录，并在可能的情况下用断言检查。

\textbf{错误示例：}

\begin{lstlisting}[language=C++]
// 调用者不知道需要什么条件
double sqrt(double x) {
    // 如果 x < 0 会怎样？文档没说
    return std::sqrt(x);
}
\end{lstlisting}

\textbf{正确示例：}

\begin{lstlisting}[language=C++]
// 前置条件明确
// 前置条件: x >= 0
double sqrt(double x) {
    Expects(x >= 0);  // GSL 断言
    return std::sqrt(x);
}
\end{lstlisting}

\textbf{要点：} 前置条件、后置条件和不变量合称"设计契约"（Design by Contract）。明确声明它们可以帮助调用者正确使用函数，帮助实现者做出合理假设。

\section{详细的错误处理规则 (E.10--E.30)}

\subsection*{\textbf{规则 E.10:} 声明后置条件}

后置条件是函数保证在正常返回时成立的条件。后置条件帮助调用者了解函数的行为，帮助测试者编写测试用例。

\textbf{错误示例：}

\begin{lstlisting}[language=C++]
vector<int> sort(vector<int> v) {
    // 调用者不知道返回值有什么性质
    // 是排序后的吗？是原样返回吗？
    std::sort(v.begin(), v.end());
    return v;
}
\end{lstlisting}

\textbf{正确示例：}

\begin{lstlisting}[language=C++]
// 后置条件: 返回值是 v 的排列，且元素非递减
vector<int> sort(vector<int> v) {
    std::sort(v.begin(), v.end());
    Ensures(is_sorted(v.begin(), v.end()));
    return v;
}
\end{lstlisting}

\textbf{要点：} 后置条件可以用 \texttt{Ensures} 宏（来自 GSL）在调试模式下检查。在发布模式下这些检查可以被移除以避免性能损失。

\subsection*{\textbf{规则 E.12:} 当函数无法完成任务时，使用 \texttt{noexcept}}

如果函数在失败时应当终止程序而非抛出异常（例如析构函数、移动操作），应当将其标记为 \texttt{noexcept}。

\textbf{错误示例：}

\begin{lstlisting}[language=C++]
~Widget() {
    cleanup();  // 如果 cleanup 抛出异常，程序终止
    // 但析构函数默认是 noexcept 的，所以异常会导致 std::terminate
}
\end{lstlisting}

\textbf{正确示例：}

\begin{lstlisting}[language=C++]
~Widget() noexcept {
    try {
        cleanup();
    } catch (...) {
        // 记录日志但不抛出
        log_error("cleanup failed in destructor");
    }
}
\end{lstlisting}

\textbf{要点：} 析构函数、移动构造函数、移动赋值运算符应当是 \texttt{noexcept} 的。这不仅是为了正确性，也是为了让标准容器在移动时使用更高效的移动操作而非拷贝。

\subsection*{\textbf{规则 E.13:} 不要直接抛出正在被捕获的异常}

如果你需要重新抛出异常，使用 \texttt{throw;} （不带表达式）来重新抛出当前异常，而不是抛出一个新异常。

\textbf{错误示例：}

\begin{lstlisting}[language=C++]
try {
    // ...
} catch (const MyError& e) {
    log(e.what());
    throw e;  // 可能切片，丢失派生类信息
}
\end{lstlisting}

\textbf{正确示例：}

\begin{lstlisting}[language=C++]
try {
    // ...
} catch (const MyError& e) {
    log(e.what());
    throw;  // 重新抛出原始异常，保留完整类型
}
\end{lstlisting}

\textbf{要点：} \texttt{throw;} 重新抛出的是原始异常对象，保留了完整的动态类型信息。\texttt{throw e;} 可能会发生对象切片。

\subsection*{\textbf{规则 E.14--E.17:} 异常安全级别}

C++中的异常安全分为三个级别：

\textbf{基本保证（Basic guarantee）：} 异常发生后，所有对象仍处于有效状态，没有资源泄漏。

\textbf{强保证（Strong guarantee）：} 异常发生后，程序状态回滚到调用函数之前。即操作要么完全成功，要么完全失败（"提交或回滚"语义）。

\textbf{无异常保证（No-throw guarantee）：} 函数永不抛出异常。

\textbf{错误示例：}

\begin{lstlisting}[language=C++]
void push_all(vector<int>& v, const vector<int>& src) {
    for (int x : src)
        v.push_back(x);  // 如果中间抛出异常，v 处于部分更新状态
}
\end{lstlisting}

\textbf{正确示例：}

\begin{lstlisting}[language=C++]
// 强保证：要么全部成功，要么 v 不变
void push_all(vector<int>& v, const vector<int>& src) {
    vector<int> tmp = v;  // 拷贝
    for (int x : src)
        tmp.push_back(x);
    v = move(tmp);  // 移动赋值是 noexcept 的
}
\end{lstlisting}

\textbf{要点：} 尽可能提供强保证。使用 copy-and-swap 惯用法是实现强保证的经典技术。

\subsection*{\textbf{规则 E.18:} 使用 RAII 管理资源，避免手动清理}

这与 R.1 和 E.6 一致。在异常处理语境下，特别强调 \texttt{finally} 模式的替代方案。C++没有 \texttt{finally} 块，但 RAII 提供了更优雅的等价方案。

\textbf{错误示例：}

\begin{lstlisting}[language=C++]
void process() {
    lock_mutex(m);
    try {
        do_work();
    } catch (...) {
        unlock_mutex(m);  // 手动清理
        throw;
    }
    unlock_mutex(m);  // 手动清理
}
\end{lstlisting}

\textbf{正确示例：}

\begin{lstlisting}[language=C++]
void process() {
    lock_guard<mutex> lock(m);  // RAII
    do_work();
}  // 自动解锁，无论是否抛出异常
\end{lstlisting}

\textbf{要点：} RAII 消除了对 \texttt{finally} 块的需求。每一个需要"无论怎样都要执行"的操作，都应当封装在 RAII 对象中。

\subsection*{\textbf{规则 E.19:} 使用 \texttt{finally} 对象替代 \texttt{finally} 块}

如果确实需要一个"作用域退出"操作（不便于用 RAII 类封装的），可以使用 \texttt{finally} 辅助对象。

\textbf{正确示例：}

\begin{lstlisting}[language=C++]
void process() {
    auto cleanup = finally([] {
        release_special_resource();
    });
    do_work();
}  // cleanup 在作用域退出时自动调用
\end{lstlisting}

\textbf{要点：} \texttt{finally} 不是标准库的一部分，但 GSL 提供了实现。它是 RAII 的补充，适用于一次性或临时的清理操作。

\subsection*{\textbf{规则 E.25:} 如果不能抛出异常，模拟 RAII}

在禁止使用异常的环境中（如某些嵌入式系统或游戏引擎），仍然需要管理资源。此时应当模拟 RAII 的行为——使用析构函数来释放资源，即使不使用异常。

\textbf{要点：} 即使没有异常，RAII 的原则仍然适用。资源的获取和释放应当绑定到对象的生命周期。区别仅在于错误报告方式：用返回值或错误码替代异常。

\subsection*{\textbf{规则 E.26:} 不要在异常处理中进行可能失败的操作}

\texttt{catch} 块中的代码应当尽量简单，避免执行可能再次抛出异常的操作。如果 \texttt{catch} 块中的操作再次抛出异常，而该异常未被捕获，程序将调用 \texttt{std::terminate}。

\textbf{错误示例：}

\begin{lstlisting}[language=C++]
try {
    process();
} catch (const Error& e) {
    send_notification(e);  // 如果发送失败呢？
    throw;
}
\end{lstlisting}

\textbf{正确示例：}

\begin{lstlisting}[language=C++]
try {
    process();
} catch (const Error& e) {
    try {
        send_notification(e);
    } catch (...) {
        // 吞掉次要错误，确保原始异常被传播
    }
    throw;
}
\end{lstlisting}

\textbf{要点：} 异常处理代码应当是"安全港"——它本身不应当引入新的失败模式。

\subsection*{\textbf{规则 E.27:} 使用 \texttt{expected} 或 \texttt{optional} 替代异常}

对于"预期内"的错误（如查找失败、解析失败），使用 \texttt{std::optional}（C++17）或 \texttt{std::expected}（C++23）比异常更合适。这些类型将错误作为值的一部分，强制调用者处理错误情况。

\textbf{错误示例：}

\begin{lstlisting}[language=C++]
// 查找失败是正常情况，不应用异常
int find(const vector<int>& v, int target) {
    for (size_t i = 0; i < v.size(); ++i)
        if (v[i] == target) return i;
    throw NotFound();  // 查找失败不是"异常"情况
}
\end{lstlisting}

\textbf{正确示例：}

\begin{lstlisting}[language=C++]
optional<size_t> find(const vector<int>& v, int target) {
    for (size_t i = 0; i < v.size(); ++i)
        if (v[i] == target) return i;
    return nullopt;  // 调用者必须检查返回值
}
\end{lstlisting}

\textbf{要点：} 异常适合处理"不应该发生"的错误，\texttt{optional}/\texttt{expected} 适合处理"预期内可能发生"的错误。区分这两者是设计良好错误处理策略的关键。

\subsection*{\textbf{规则 E.28:} 避免基于错误码的错误处理}

错误码容易被忽略，且会污染正常的返回值。如果调用者忘记检查错误码，错误将被静默忽略。

\textbf{错误示例：}

\begin{lstlisting}[language=C++]
int result = process();
// 调用者忘记检查 result
do_next_step();  // 在错误状态下继续执行
\end{lstlisting}

\textbf{正确示例：}

\begin{lstlisting}[language=C++]
// 异常方式：不可能忽略错误
process();  // 如果失败，异常会中断正常流程

// 或 optional 方式：编译器可以警告未使用的返回值
auto result = process();
if (!result) { /* 必须处理 */ }
\end{lstlisting}

\textbf{要点：} 如果必须使用错误码，至少使用 \texttt{[[nodiscard]]} 属性标记返回值，让编译器在未检查时发出警告。

\section*{练习题}

\begin{enumerate}
\item 设计一个文件处理类 \texttt{FileProcessor}，要求：构造函数打开文件（失败时抛出异常），析构函数关闭文件，提供读取方法。确保该类是异常安全的，并满足值语义（可拷贝、可移动）。

\item 以下代码有什么问题？如何修复？
\begin{lstlisting}[language=C++]
class Database {
    Connection* conn_;
public:
    Database(const string& url) {
        conn_ = connect(url);
        if (!conn_) return;  // 构造"成功"但无效
    }
    void query(const string& sql) {
        conn_->execute(sql);  // 如果 conn_ 为空？
    }
    ~Database() {
        if (conn_) conn_->close();
    }
};
\end{lstlisting}

\item 为以下函数添加适当的异常安全保证（基本保证或强保证）：
\begin{lstlisting}[language=C++]
void transfer(Account& from, Account& to, double amount) {
    from.balance -= amount;
    to.balance += amount;
}
\end{lstlisting}

\item 解释为什么以下析构函数是危险的，并给出修复方案：
\begin{lstlisting}[language=C++]
~Logger() {
    flush_buffer();  // 可能抛出异常
    close_file();    // 可能抛出异常
}
\end{lstlisting}

\item 设计一个错误处理策略，适用于一个网络服务器应用。考虑以下场景：(a) 客户端连接超时；(b) 请求解析失败；(c) 数据库查询失败；(d) 内存分配失败。对每种场景，说明应当使用异常、错误码还是终止程序。
\end{enumerate}

% ============================================================
% Part V: 表达式与性能
% ============================================================

\part{表达式与性能}

% ============================================================
% Chapter 11: 表达式与语句
% ============================================================

\chapter{表达式与语句}

表达式和语句是程序的基本构建块。C++ Core Guidelines 的"ES"（Expressive Statements）系列规则旨在帮助开发者编写更清晰、更安全、更少出错的表达式和语句。这些规则涵盖了从变量命名到循环结构的方方面面，其核心思想是：让代码直接表达意图，减少隐式行为和未定义行为的可能性。

\section{通用表达式规则 (ES.1--ES.10)}

\subsection*{\textbf{规则 ES.1:} 适当时使用标准库而非手写代码}

标准库经过广泛的测试和优化，使用标准库代码通常比手写代码更正确、更高效、更易维护。例如，使用 \texttt{std::sort} 而非手写排序算法，使用 \texttt{std::find} 而非手写循环查找。

\textbf{错误示例：}

\begin{lstlisting}[language=C++]
// 手写查找
bool found = false;
for (int i = 0; i < v.size(); ++i) {
    if (v[i] == target) { found = true; break; }
}

// 手写最大值
int max_val = v[0];
for (int i = 1; i < v.size(); ++i) {
    if (v[i] > max_val) max_val = v[i];
}
\end{lstlisting}

\textbf{正确示例：}

\begin{lstlisting}[language=C++]
// 使用标准库
auto it = find(v.begin(), v.end(), target);
bool found = (it != v.end());

int max_val = *max_element(v.begin(), v.end());
\end{lstlisting}

\textbf{要点：} 标准库算法通常有经过验证的正确性保证，且编译器可以对标准库算法进行特殊优化。手写代码不仅容易出错（边界条件、空容器等），还可能不如标准库高效。

\subsection*{\textbf{规则 ES.2:} 适当时使用 \texttt{auto} 避免类型名重复}

\texttt{auto} 可以避免类型名的冗余书写，防止意外的类型转换，并使代码更加泛型化。但 \texttt{auto} 不应被过度使用——当类型信息对读者有帮助时，应当显式写出类型。

\textbf{错误示例：}

\begin{lstlisting}[language=C++]
map<string, vector<pair<int, int>>>::iterator it = m.begin();
// 类型名冗长且重复
\end{lstlisting}

\textbf{正确示例：}

\begin{lstlisting}[language=C++]
auto it = m.begin();
// 或者更好的做法：使用范围 for
for (auto& [key, value] : m) {
    // 直接使用结构化绑定
}
\end{lstlisting}

\textbf{要点：} \texttt{auto} 的核心价值在于：(1) 避免冗余；(2) 避免窄化转换；(3) 使代码泛型化。但不要用 \texttt{auto} 隐藏重要信息——如果读者需要知道确切类型才能理解代码，就应当显式写出类型。

\subsection*{\textbf{规则 ES.3:} 不要重复类型名——使用别名或 \texttt{auto}}

这与 ES.2 相关。重复类型名不仅冗长，而且在类型改变时需要修改多处代码。

\textbf{错误示例：}

\begin{lstlisting}[language=C++]
MyLongTypeName<int, string>* p = new MyLongTypeName<int, string>();
\end{lstlisting}

\textbf{正确示例：}

\begin{lstlisting}[language=C++]
auto p = make_unique<MyLongTypeName<int, string>>();
// 或使用类型别名
using ShortName = MyLongTypeName<int, string>;
auto p = make_unique<ShortName>();
\end{lstlisting}

\textbf{要点：} DRY（Don't Repeat Yourself）原则同样适用于类型名。

\subsection*{\textbf{规则 ES.5:} 避免魔数——使用命名常量}

魔数（magic number）是代码中直接出现的数字字面量，其含义不直观。使用命名常量可以自文档化代码，并使修改更加安全。

\textbf{错误示例：}

\begin{lstlisting}[language=C++]
if (status == 4) { /* 4 是什么？ */ }
double area = 3.14159 * r * r;
for (int i = 0; i < 256; ++i) { /* 256 是什么？ */ }
\end{lstlisting}

\textbf{正确示例：}

\begin{lstlisting}[language=C++]
constexpr int status_error = 4;
if (status == status_error) { /* 含义清晰 */ }

constexpr double pi = 3.14159265358979;
double area = pi * r * r;

constexpr int buffer_size = 256;
for (int i = 0; i < buffer_size; ++i) { /* 含义清晰 */ }
\end{lstlisting}

\textbf{要点：} 使用 \texttt{constexpr} 定义编译期常量。它们不仅自文档化，还能让编译器进行优化。避免使用 \texttt{\#define} 定义常量——它没有作用域，不参与类型检查。

\subsection*{\textbf{规则 ES.10:} 声明变量时同时初始化}

未初始化的变量是C++中最常见的错误来源之一。声明变量时立即初始化可以消除一整类错误。

\textbf{错误示例：}

\begin{lstlisting}[language=C++]
int x;        // 未初始化
string s;     // 默认初始化（空字符串），但可能不是期望的值
// ... 多行代码 ...
x = 42;       // 如果中间使用了 x，就是未定义行为
\end{lstlisting}

\textbf{正确示例：}

\begin{lstlisting}[language=C++]
int x = 42;
string s = "hello";
// 或者使用直接初始化
vector<int> v = {1, 2, 3};
\end{lstlisting}

\textbf{要点：} 如果变量的初始值需要计算，可以使用 lambda 立即调用：
\begin{lstlisting}[language=C++]
const int x = [&] {
    int result = 0;
    for (int i = 0; i < n; ++i) result += compute(i);
    return result;
}();
\end{lstlisting}

\section{变量规则 (ES.11--ES.23)}

\subsection*{\textbf{规则 ES.11:} 使用 \texttt{const} 保护不可变量}

如果一个变量在初始化后不再被修改，应当声明为 \texttt{const}。这不仅是一种文档化手段，还能让编译器捕获意外的修改。

\textbf{错误示例：}

\begin{lstlisting}[language=C++]
int max_size = 100;
// ... 100行代码 ...
max_size = 50;  // 意外修改！
\end{lstlisting}

\textbf{正确示例：}

\begin{lstlisting}[language=C++]
const int max_size = 100;
// ... 100行代码 ...
max_size = 50;  // 编译错误！
\end{lstlisting}

\textbf{要点：} \texttt{const} 是C++中表达"不修改"意图的基本工具。默认使用 \texttt{const}，只在确实需要修改时才去掉它。这被称为"const by default"原则。

\subsection*{\textbf{规则 ES.12:} 不要在类作用域外使用 \texttt{using namespace}}

在头文件中使用 \texttt{using namespace} 会污染所有包含该头文件的翻译单元的命名空间，可能导致名称冲突。

\textbf{错误示例：}

\begin{lstlisting}[language=C++]
// 在头文件中
#include <vector>
using namespace std;  // 所有包含此头文件的文件都受影响

class MyContainer {
    vector<int> data;  // 看起来正常，但...
};
\end{lstlisting}

\textbf{正确示例：}

\begin{lstlisting}[language=C++]
#include <vector>

class MyContainer {
    std::vector<int> data;  // 明确使用 std::
};

// 在 .cpp 文件的函数作用域内可以使用
void foo() {
    using namespace std;  // 只影响当前函数
}
\end{lstlisting}

\textbf{要点：} 头文件中的 \texttt{using namespace} 是C++中最常见的反模式之一。它可能在不经意间改变大量代码的含义。

\subsection*{\textbf{规则 ES.20:} 避免全局变量}

全局变量（包括命名空间作用域的变量和静态成员变量）是代码可维护性的敌人。它们使得代码难以理解（任何函数都可能修改它们）、难以测试（状态在测试之间泄漏）、难以并行化（数据竞争）。

\textbf{错误示例：}

\begin{lstlisting}[language=C++]
int global_counter = 0;  // 全局变量

void process() {
    global_counter++;  // 任何函数都可以修改
}
\end{lstlisting}

\textbf{正确示例：}

\begin{lstlisting}[language=C++]
// 使用类封装状态
class Processor {
    int counter_ = 0;
public:
    void process() { counter_++; }
    int count() const { return counter_; }
};
\end{lstlisting}

\textbf{要点：} 如果确实需要全局状态（如配置、日志器），使用单例模式或依赖注入，并确保线程安全。全局 \texttt{const} 变量（编译期常量）是可以接受的，因为它们不会改变。

\section{声明与初始化 (ES.21--ES.30)}

\subsection*{\textbf{规则 ES.21:} 不要在声明变量之前引入初始化}

变量的声明和初始化应当尽可能接近。不要在远离使用位置的地方声明变量——这会增加读者追踪变量状态的认知负担。

\textbf{错误示例：}

\begin{lstlisting}[language=C++]
int i, j, k, x, y, z;  // 在函数开头声明所有变量（C风格）
// ... 50行代码 ...
i = 1; j = 2;
\end{lstlisting}

\textbf{正确示例：}

\begin{lstlisting}[language=C++]
int i = 1;
int j = 2;
// 在使用前立即声明和初始化
\end{lstlisting}

\textbf{要点：} C++允许在任意位置声明变量。利用这一特性，将变量声明推迟到第一次使用处，并同时初始化。

\subsection*{\textbf{规则 ES.22:} 不要在声明中混合变量和逗号}

在一条声明语句中声明多个变量会降低可读性，并容易导致指针/引用声明的错误。

\textbf{错误示例：}

\begin{lstlisting}[language=C++]
int* p, q;  // p 是指针，q 不是！容易误解
int a = 1, b = 2, c = 3;  // 冗长
\end{lstlisting}

\textbf{正确示例：}

\begin{lstlisting}[language=C++]
int* p = nullptr;
int q = 0;  // 分开声明，清晰明了

int a = 1;
int b = 2;
int c = 3;
\end{lstlisting}

\textbf{要点：} 每条声明语句只声明一个变量。这不仅提高可读性，还避免了 \texttt{int* p, q} 这类经典陷阱。

\subsection*{\textbf{规则 ES.23:} 使用 \texttt{\{\}} 初始化语法}

花括号初始化（列表初始化）可以防止窄化转换，且语法统一。它是C++中最安全的初始化方式。

\textbf{错误示例：}

\begin{lstlisting}[language=C++]
int x = 3.14;     // 隐式截断为 3，无警告
double d = 1000;  // 可能丢失精度（在某些平台上）
\end{lstlisting}

\textbf{正确示例：}

\begin{lstlisting}[language=C++]
int x{3.14};      // 编译错误！窄化转换
double d{1000};   // OK，无窄化

vector<int> v{1, 2, 3};  // 统一的初始化语法
\end{lstlisting}

\textbf{要点：} 花括号初始化是"默认"初始化方式。只有在明确需要时才使用圆括号初始化（如调用特定构造函数）。

\section{表达式规则 (ES.31--ES.40)}

\subsection*{\textbf{规则 ES.31:} 不要用C风格强制类型转换}

C风格的强制类型转换（\texttt{(T)expr}）功能强大但危险——它可以执行任何类型的转换，编译器无法区分有意转换和无意转换。C++提供了更精确的转换关键字。

\textbf{错误示例：}

\begin{lstlisting}[language=C++]
double d = 3.14;
int i = (int)d;           // C风格转换
Base* b = new Derived;
Derived* d2 = (Derived*)b; // C风格转换，可能不安全
\end{lstlisting}

\textbf{正确示例：}

\begin{lstlisting}[language=C++]
double d = 3.14;
int i = static_cast<int>(d);  // 明确：数值转换

Base* b = new Derived;
Derived* d2 = dynamic_cast<Derived*>(b);  // 安全：运行时检查
\end{lstlisting}

\textbf{要点：} C++的四种转换关键字各有用途：\texttt{static\_cast}（安全的隐式转换的逆操作）、\texttt{dynamic\_cast}（安全的多态向下转换）、\texttt{const\_cast}（去除 \texttt{const}）、\texttt{reinterpret\_cast}（低层位模式转换，极少使用）。使用它们可以在代码搜索中轻松找到所有转换点。

\subsection*{\textbf{规则 ES.40:} 避免危险表达式}

某些表达式虽然合法但容易导致错误。例如，在同一表达式中多次修改同一变量（序列点问题），或使用复杂的嵌套三元表达式。

\textbf{错误示例：}

\begin{lstlisting}[language=C++]
int x = ++i + i++;  // 未定义行为！
int y = (x > 0) ? (x > 10 ? x : -x) : 0;  // 难以阅读
\end{lstlisting}

\textbf{正确示例：}

\begin{lstlisting}[language=C++]
int x = i++;
x += i;  // 分步操作，清晰

int y;
if (x > 0)
    y = (x > 10) ? x : -x;
else
    y = 0;
\end{lstlisting}

\textbf{要点：} 表达式的复杂性应当与读者的理解能力匹配。如果一个表达式需要多次阅读才能理解，就应当拆分为多条语句。

\section{语句规则 (ES.41--ES.70)}

\subsection*{\textbf{规则 ES.41:} 如果不能确定运算符优先级，使用括号}

C++的运算符优先级规则复杂且容易记忆错误。如果不确定，使用括号明确优先级。

\textbf{错误示例：}

\begin{lstlisting}[language=C++]
if (a & b == c)  // 意图是 (a & b) == c 还是 a & (b == c)？
int result = a + b << 2;  // + 和 << 哪个优先？
\end{lstlisting}

\textbf{正确示例：}

\begin{lstlisting}[language=C++]
if ((a & b) == c)  // 明确
int result = a + (b << 2);  // 明确
\end{lstlisting}

\textbf{要点：} 编译器不会因为你多加了括号而惩罚你（编译器会优化掉多余的括号），但会因为优先级错误而产生难以调试的bug。

\subsection*{\textbf{规则 ES.43:} 避免未定义行为的表达式}

C++中有许多未定义行为（UB）陷阱：有符号整数溢出、空指针解引用、越界访问、数据竞争等。这些规则的核心是：不要依赖未定义行为，即使它在某个编译器上"碰巧能工作"。

\textbf{错误示例：}

\begin{lstlisting}[language=C++]
int arr[5] = {1, 2, 3, 4, 5};
int x = arr[10];  // 未定义行为：越界访问

int* p = nullptr;
int y = *p;       // 未定义行为：空指针解引用
\end{lstlisting}

\textbf{正确示例：}

\begin{lstlisting}[language=C++]
array<int, 5> arr = {1, 2, 3, 4, 5};
int x = arr.at(10);  // 抛出异常，而非未定义行为

int* p = get_pointer();
if (p) int y = *p;   // 检查后再使用
\end{lstlisting}

\textbf{要点：} 未定义行为是C++中最危险的特性。它意味着"任何事情都可能发生"——程序可能崩溃、产生错误结果、或在调试版本中工作但在发布版本中失败。使用 \texttt{at()} 而非 \texttt{[]}，使用智能指针而非裸指针，使用边界检查，都是避免 UB 的有效手段。

\subsection*{\textbf{规则 ES.44:} 不要依赖宏}

宏（\texttt{\#define}）不参与作用域规则，不进行类型检查，错误信息难以理解。应当使用 \texttt{constexpr}、\texttt{const}、\texttt{inline}、\texttt{enum class} 等语言特性替代宏。

\textbf{错误示例：}

\begin{lstlisting}[language=C++]
#define MAX_SIZE 100
#define SQUARE(x) ((x) * (x))
#define DEBUG

int arr[MAX_SIZE];
int y = SQUARE(a + b);  // 宏展开可能有问题
\end{lstlisting}

\textbf{正确示例：}

\begin{lstlisting}[language=C++]
constexpr int max_size = 100;
constexpr auto square = [](auto x) { return x * x; };

int arr[max_size];
int y = square(a + b);  // 类型安全，可调试
\end{lstlisting}

\textbf{要点：} 宏是C的遗产，在C++中几乎总有更好的替代方案。唯一的例外是头文件保护符（include guard），但 \texttt{\#pragma once} 也是广泛支持的替代方案。

\subsection*{\textbf{规则 ES.45:} 避免"魔数"和"魔字符"}

这与 ES.5 相关。特别强调字符字面量也应当避免直接使用。

\textbf{错误示例：}

\begin{lstlisting}[language=C++]
if (ch == '\x03') { /* 这是什么字符？ */ }
buffer[1024] = '\0';
\end{lstlisting}

\textbf{正确示例：}

\begin{lstlisting}[language=C++]
constexpr char ETX = '\x03';  // End of Text
if (ch == ETX) { /* 含义清晰 */ }

constexpr size_t buffer_size = 1024;
buffer[buffer_size - 1] = '\0';
\end{lstlisting}

\textbf{要点：} 每个字面量都应当有名字。名字不仅传达含义，还使得修改只需要在一处进行。

\subsection*{\textbf{规则 ES.46:} 避免窄化转换}

窄化转换（narrowing conversion）是指在转换过程中可能丢失信息的转换，如 \texttt{double} 到 \texttt{int}、\texttt{int} 到 \texttt{short} 等。这些转换往往是bug的来源。

\textbf{错误示例：}

\begin{lstlisting}[language=C++]
double d = 3.14;
int i = d;          // 隐式窄化：丢失小数部分
int j = 1000;
short s = j;        // 隐式窄化：可能溢出
\end{lstlisting}

\textbf{正确示例：}

\begin{lstlisting}[language=C++]
double d = 3.14;
int i = static_cast<int>(d);  // 明确：我知道我在截断
// 或使用 gsl::narrow_cast 进行运行时检查
int j = 1000;
short s = gsl::narrow<short>(j);  // 运行时检查溢出
\end{lstlisting}

\textbf{要点：} 使用花括号初始化可以防止隐式窄化转换。如果确实需要窄化，使用 \texttt{static\_cast} 明确表达意图。

\subsection*{\textbf{规则 ES.50:} 不要把 \texttt{const} 转换掉}

\texttt{const\_cast} 是C++中最危险的转换之一。去除一个对象的 \texttt{const} 属性并修改它，是未定义行为（如果原始对象是 \texttt{const} 的）。

\textbf{错误示例：}

\begin{lstlisting}[language=C++]
const int x = 42;
int& rx = const_cast<int&>(x);
rx = 100;  // 未定义行为！x 是 const 的
\end{lstlisting}

\textbf{正确示例：}

\begin{lstlisting}[language=C++]
// 如果需要修改，一开始就不要声明为 const
int x = 42;
x = 100;  // OK

// 如果接口返回 const 但你确实需要修改，说明设计有问题
// 考虑重新设计接口
\end{lstlisting}

\textbf{要点：} \texttt{const\_cast} 的唯一合法用途是与不支持 \texttt{const} 的旧式API交互。即使如此，也应当确保不修改原始对象。

\subsection*{\textbf{规则 ES.55:} 使用范围 \texttt{for} 而非下标循环}

范围 \texttt{for}（range-based for）比下标循环更安全、更简洁，且不容易出错（如越界、错误的循环变量）。

\textbf{错误示例：}

\begin{lstlisting}[language=C++]
for (int i = 0; i < v.size(); ++i) {
    cout << v[i] << endl;
}
\end{lstlisting}

\textbf{正确示例：}

\begin{lstlisting}[language=C++]
for (const auto& elem : v) {
    cout << elem << endl;
}
\end{lstlisting}

\textbf{要点：} 范围 \texttt{for} 适用于所有标准容器、数组、以及任何提供了 \texttt{begin()}/\texttt{end()} 的类型。只有在需要索引时才使用下标循环。

\subsection*{\textbf{规则 ES.60:} 避免 \texttt{new} 和 \texttt{delete} 表达式}

直接使用的 \texttt{new} 和 \texttt{delete} 是资源泄漏和悬空指针的主要来源。应当使用智能指针和容器来管理动态内存。

\textbf{错误示例：}

\begin{lstlisting}[language=C++]
Widget* w = new Widget();
// ... 可能忘记 delete，或异常导致 delete 不执行
delete w;
\end{lstlisting}

\textbf{正确示例：}

\begin{lstlisting}[language=C++]
auto w = make_unique<Widget>();
// 自动释放，无需手动 delete
\end{lstlisting}

\textbf{要点：} 在现代C++中，直接出现 \texttt{new} 和 \texttt{delete} 的代码应当被视为代码异味。它们只应当出现在智能指针和容器的实现内部。

\subsection*{\textbf{规则 ES.65:} 不要解引用空指针}

空指针解引用是C++中最常见的崩溃原因之一。在解引用指针之前，应当检查其是否为空。更好的做法是使用引用替代可能为空的指针。

\textbf{错误示例：}

\begin{lstlisting}[language=C++]
Widget* w = find_widget("foo");
w->do_something();  // 如果 find_widget 返回 nullptr？
\end{lstlisting}

\textbf{正确示例：}

\begin{lstlisting}[language=C++]
Widget* w = find_widget("foo");
if (w) w->do_something();

// 或者更好的做法：如果不应该为空，使用引用
Widget& w = get_widget("foo");  // 保证非空
w.do_something();
\end{lstlisting}

\textbf{要点：} 如果一个指针可能为空，使用 \texttt{optional<reference\_wrapper<T>>} 或在使用前检查。如果指针不应为空，使用引用。

\subsection*{\textbf{规则 ES.70:} 优先使用 \texttt{switch} 而非 \texttt{if-else} 链}

当需要在多个离散值之间选择时，\texttt{switch} 语句比长链的 \texttt{if-else} 更清晰，且编译器可以对 \texttt{switch} 进行优化（跳转表）。此外，编译器可以在 \texttt{switch} 缺少 \texttt{case} 时发出警告。

\textbf{错误示例：}

\begin{lstlisting}[language=C++]
if (op == '+') add();
else if (op == '-') subtract();
else if (op == '*') multiply();
else if (op == '/') divide();
// 容易遗漏 case，编译器无法警告
\end{lstlisting}

\textbf{正确示例：}

\begin{lstlisting}[language=C++]
switch (op) {
case '+': add(); break;
case '-': subtract(); break;
case '*': multiply(); break;
case '/': divide(); break;
// 编译器可以在 enum 新增值时警告缺少 case
}
\end{lstlisting}

\textbf{要点：} 使用 \texttt{enum class} 配合 \texttt{switch} 可以获得编译器的穷举检查。这是类型安全和代码清晰度的双赢。

\section*{练习题}

\begin{enumerate}
\item 将以下代码中的所有魔数替换为命名常量，并解释每个常量的含义：
\begin{lstlisting}[language=C++]
void render(int mode) {
    if (mode == 1) { /* 普通模式 */ }
    else if (mode == 2) { /* 高质量模式 */ }
    for (int i = 0; i < 1920; ++i)
        for (int j = 0; j < 1080; ++j)
            buffer[i * 1920 + j] = compute(i, j);
}
\end{lstlisting}

\item 以下代码有哪些未定义行为？如何修复？
\begin{lstlisting}[language=C++]
int arr[5] = {1, 2, 3, 4, 5};
int sum = 0;
for (int i = 0; i <= 5; ++i)
    sum += arr[i];
int avg = sum / 5;
\end{lstlisting}

\item 重写以下代码，使用范围 \texttt{for} 和 \texttt{auto}：
\begin{lstlisting}[language=C++]
map<string, vector<int>> data;
for (map<string, vector<int>>::iterator it = data.begin();
     it != data.end(); ++it) {
    for (int i = 0; i < it->second.size(); ++i) {
        cout << it->first << ": " << it->second[i] << endl;
    }
}
\end{lstlisting}

\item 解释以下代码中 \texttt{auto} 推导出的类型是什么，并讨论是否应该使用 \texttt{auto}：
\begin{lstlisting}[language=C++]
auto x = 3;
auto y = 3.0;
auto z = "hello";
auto w = {1, 2, 3};
\end{lstlisting}

\item 使用 \texttt{enum class} 和 \texttt{switch} 重写以下代码：
\begin{lstlisting}[language=C++]
void handle_event(int event_type) {
    if (event_type == 0) { /* mouse click */ }
    else if (event_type == 1) { /* key press */ }
    else if (event_type == 2) { /* window resize */ }
    else if (event_type == 3) { /* window close */ }
}
\end{lstlisting}
\end{enumerate}

% ============================================================
% Chapter 12: 性能 (Performance - "Per" section)
% ============================================================

\chapter{性能}

性能优化是软件开发中的重要课题，但也是容易走入歧途的领域。C++ Core Guidelines 的"Per"系列规则传达的核心信息是：不要过早优化，不要凭直觉优化，而应当基于测量结果进行有针对性的优化。同时，有些"优化"实际上是反优化（pessimization），它们使代码更复杂却不带来性能提升，甚至降低性能。本章的规则旨在帮助开发者避免这些陷阱。

\section{基本原则 (Per.1--Per.3)}

\subsection*{\textbf{规则 Per.1:} 不要过早优化}

Donald Knuth 的名言"过早优化是万恶之源"在C++中尤其重要。现代编译器的优化能力远超大多数程序员的手动优化。在没有测量数据的情况下进行优化，不仅浪费时间，还可能引入bug、降低代码可读性，甚至降低性能。

\textbf{错误示例：}

\begin{lstlisting}[language=C++]
// "优化"：使用裸数组替代 vector，因为"vector 慢"
int* data = new int[10000];
// 手动管理内存，容易出错，且可能不如 vector 快
// （因为 vector 的实现可以利用编译器优化）
\end{lstlisting}

\textbf{正确示例：}

\begin{lstlisting}[language=C++]
// 先使用正确的抽象
vector<int> data(10000);
// 如果性能不够，用 profiler 找到瓶颈，再针对性优化
\end{lstlisting}

\textbf{要点：} 正确的优化流程是：(1) 编写清晰、正确的代码；(2) 测量性能；(3) 找到瓶颈；(4) 针对瓶颈优化；(5) 再次测量，确认优化有效。

\subsection*{\textbf{规则 Per.2:} 不要凭直觉优化——等待测量结果}

程序员对性能的直觉往往是错误的。缓存局部性、分支预测、编译器优化等因素使得代码的实际性能难以凭直觉判断。只有通过性能分析工具（profiler）才能获得可靠的数据。

\textbf{错误示例：}

\begin{lstlisting}[language=C++]
// "我觉得 string 拼接很慢，用 stringstream 替代"
// 实际上在小字符串场景下，+ 运算符可能更快
string result;
for (const auto& s : parts) {
    stringstream ss;
    ss << s;
    result += ss.str();  // 更复杂，不一定更快
}
\end{lstlisting}

\textbf{正确示例：}

\begin{lstlisting}[language=C++]
// 先写简单的版本
string result;
for (const auto& s : parts) {
    result += s;
}
// 如果性能不够，用 profiler 确认瓶颈，再优化
\end{lstlisting}

\textbf{要点：} 不要猜测性能瓶颈。使用 \texttt{perf}、\texttt{gprof}、\texttt{Valgrind/Callgrind}、Visual Studio Profiler 等工具进行实际测量。

\subsection*{\textbf{规则 Per.3:} 不要为了"效率"而牺牲代码的清晰性}

清晰的代码更容易被编译器优化，更容易被人类审查和优化。为了微小的性能提升而引入复杂的代码，往往是得不偿失的。

\textbf{要点：} 编译器擅长优化简单的、模式化的代码。复杂的、充满技巧的代码反而可能阻碍编译器优化。

\section{具体性能规则 (Per.4--Per.7)}

\subsection*{\textbf{规则 Per.4:} 记住编译器的优化能力}

现代C++编译器（GCC、Clang、MSVC）具有强大的优化能力：内联展开、循环展开、向量化、常量折叠、死代码消除等。许多程序员认为"慢"的操作，编译器可能已经优化掉了。

\textbf{要点：} 以下操作通常会被编译器优化：(1) 小对象的拷贝（内联后消除）；(2) 临时对象的创建（返回值优化 RVO/NRVO）；(3) 常量表达式的计算（编译期计算）；(4) 简单的循环（自动向量化）。不要手动"优化"这些场景。

\subsection*{\textbf{规则 Per.5:} 不要使用 \texttt{const} 引用传递小类型——按值传递}

对于小类型（如 \texttt{int}、\texttt{double}、小 \texttt{struct}），按值传递比 \texttt{const} 引用传递更高效。引用在底层是指针，解引用指针的开销可能大于直接拷贝小对象。

\textbf{错误示例：}

\begin{lstlisting}[language=C++]
void set_value(const int& value) {  // 引用传递 int
    x_ = value;  // 需要解引用指针
}
\end{lstlisting}

\textbf{正确示例：}

\begin{lstlisting}[language=C++]
void set_value(int value) {  // 按值传递 int
    x_ = value;  // 直接使用寄存器中的值
}
\end{lstlisting}

\textbf{要点：} 一般规则：对于大小不超过指针的类型（通常8字节以下），按值传递。对于更大的类型，使用 \texttt{const} 引用传递。

\subsection*{\textbf{规则 Per.6:} 不要使用 \texttt{string} 的 \texttt{const} 引用传递字面量——使用 \texttt{string\_view}}

\texttt{string\_view}（C++17）是一个不拥有字符串的轻量视图，它避免了字符串拷贝，同时可以接受字面量、\texttt{string}、\texttt{const char*} 等多种类型。

\textbf{错误示例：}

\begin{lstlisting}[language=C++]
void print(const string& s) {
    cout << s << endl;
}
print("hello");  // 创建了临时 string 对象！
\end{lstlisting}

\textbf{正确示例：}

\begin{lstlisting}[language=C++]
void print(string_view s) {
    cout << s << endl;
}
print("hello");  // 零拷贝
print(string("hello"));  // 也OK
\end{lstlisting}

\textbf{要点：} \texttt{string\_view} 是只读字符串参数的最佳选择。它的大小与指针相同（一个指针加一个长度），不拥有内存，不分配内存。但要注意：\texttt{string\_view} 不保证空终止，不能替代需要 \texttt{null-terminated} 字符串的API。

\subsection*{\textbf{规则 Per.7:} 设计数据结构以优化缓存局部性}

现代CPU的内存访问速度远低于CPU速度，缓存命中率是性能的关键因素。设计数据结构时，应当考虑缓存局部性——频繁一起访问的数据应当在内存中相邻。

\textbf{错误示例：}

\begin{lstlisting}[language=C++]
// 链表：每个节点分散在堆上，缓存不友好
struct Node {
    int data;
    Node* next;
};
// 遍历时缓存命中率极低
\end{lstlisting}

\textbf{正确示例：}

\begin{lstlisting}[language=C++]
// 数组：元素连续存储，缓存友好
vector<int> data;
// 遍历时缓存命中率高
\end{lstlisting}

\textbf{要点：} 数据导向设计（Data-Oriented Design）的核心思想是：按照CPU访问数据的方式来组织数据，而非按照面向对象的抽象层次。数组优于链表，结构体数组（AoS）和数组结构体（SoA）的选择取决于访问模式。

\section{更多性能建议}

\subsection*{移动语义与性能}

移动语义（C++11）是C++中最重要的性能优化特性之一。通过移动而非拷贝，可以避免大量不必要的内存分配和拷贝。

\textbf{错误示例：}

\begin{lstlisting}[language=C++]
vector<string> create_names() {
    vector<string> result;
    // ... 填充 result ...
    return result;  // C++11之前：拷贝整个 vector
}
\end{lstlisting}

\textbf{正确示例：}

\begin{lstlisting}[language=C++]
vector<string> create_names() {
    vector<string> result;
    // ... 填充 result ...
    return result;  // C++11及之后：自动移动（或RVO消除拷贝）
}
\end{lstlisting}

\textbf{要点：} 确保自定义类型正确实现了移动构造函数和移动赋值运算符。标准容器已经正确实现了移动语义。

\subsection*{避免不必要的内存分配}

动态内存分配（\texttt{new}/\texttt{malloc}）是相对昂贵的操作。在性能敏感的代码路径中，应当尽量减少内存分配。

\textbf{错误示例：}

\begin{lstlisting}[language=C++]
// 在循环中反复创建和销毁 string
for (int i = 0; i < 1000000; ++i) {
    string s = to_string(i);  // 每次循环分配内存
    process(s);
}
\end{lstlisting}

\textbf{正确示例：}

\begin{lstlisting}[language=C++]
// 复用 string 缓冲区
string s;
s.reserve(32);  // 预分配足够空间
for (int i = 0; i < 1000000; ++i) {
    s = to_string(i);  // 复用已分配的内存
    process(s);
}
\end{lstlisting}

\textbf{要点：} \texttt{reserve()} 是减少 \texttt{vector} 和 \texttt{string} 重新分配的有效手段。对象池（object pool）模式可以在极端性能要求下减少分配开销。

\subsection*{编译期计算}

\texttt{constexpr} 允许将计算从运行时移到编译时，消除运行时开销。

\textbf{正确示例：}

\begin{lstlisting}[language=C++]
// 编译期计算阶乘
constexpr int factorial(int n) {
    if (n <= 1) return 1;
    return n * factorial(n - 1);
}

constexpr int fact5 = factorial(5);  // 编译期计算为 120
\end{lstlisting}

\textbf{要点：} 尽可能将函数声明为 \texttt{constexpr}。编译器会在可能的情况下在编译期求值，在不可能时自动退回到运行时求值。

\section*{练习题}

\begin{enumerate}
\item 以下代码中哪些是"过早优化"？哪些是合理的优化？
\begin{lstlisting}[language=C++]
// (a) 使用裸数组替代 vector
// (b) 使用 string::reserve 预分配空间
// (c) 使用位移替代乘除法
// (d) 使用移动语义减少拷贝
\end{lstlisting}

\item 编写一个基准测试程序，比较以下两种字符串拼接方式的性能：(a) 使用 \texttt{operator+}；(b) 使用 \texttt{stringstream}。分析结果并解释原因。

\item 解释为什么以下代码中按值传递比 \texttt{const} 引用更快：
\begin{lstlisting}[language=C++]
void foo(int x);        // 按值
void bar(const int& x); // const 引用
\end{lstlisting}

\item 设计一个缓存友好的数据结构，用于存储和查询二维网格上的点。解释你的设计如何利用缓存局部性。

\item 将以下函数改写为 \texttt{constexpr}，使其能在编译期求值：
\begin{lstlisting}[language=C++]
int fibonacci(int n) {
    if (n <= 1) return n;
    return fibonacci(n-1) + fibonacci(n-2);
}
\end{lstlisting}
\end{enumerate}

% ============================================================
% Part VI: 并发与现代实践
% ============================================================

\part{并发与现代实践}

% ============================================================
% Chapter 13: 并发与并行
% ============================================================

\chapter{并发与并行}

并发编程是现代软件开发中最具挑战性的领域之一。C++11引入了内存模型和线程库，为跨平台并发编程提供了语言层面的支持。Core Guidelines 的"CP"（Concurrency and Parallelism）系列规则旨在帮助开发者编写正确、高效的并发代码，避免数据竞争、死锁和其他并发相关的错误。

\section{通用并发规则 (CP.1--CP.10)}

\subsection*{\textbf{规则 CP.1:} 假设你的代码将作为并发程序的一部分运行}

即使当前代码是单线程的，将来也可能被多线程环境使用。因此，应当避免全局可变状态，确保函数是线程安全的（或明确标记为非线程安全）。

\textbf{错误示例：}

\begin{lstlisting}[language=C++]
// 全局可变状态——多线程下必然出错
int call_count = 0;

void process() {
    call_count++;  // 数据竞争
    // ...
}
\end{lstlisting}

\textbf{正确示例：}

\begin{lstlisting}[language=C++]
// 使用原子变量
atomic<int> call_count{0};

void process() {
    call_count++;  // 线程安全
    // ...
}
\end{lstlisting}

\textbf{要点：} 编写线程安全代码的成本远低于将非线程安全代码改造为线程安全的成本。从一开始就考虑并发安全性是更经济的选择。

\subsection*{\textbf{规则 CP.2:} 避免数据竞争——确保共享可变状态的访问是串行的}

数据竞争（data race）是并发编程中最常见也最危险的错误。当两个或多个线程同时访问同一内存位置，且至少有一个是写操作时，就发生了数据竞争。数据竞争是未定义行为。

\textbf{错误示例：}

\begin{lstlisting}[language=C++]
int counter = 0;

void increment() {
    for (int i = 0; i < 1000; ++i)
        counter++;  // 数据竞争！
}

// 两个线程调用 increment()，结果不确定
\end{lstlisting}

\textbf{正确示例：}

\begin{lstlisting}[language=C++]
mutex mtx;
int counter = 0;

void increment() {
    for (int i = 0; i < 1000; ++i) {
        lock_guard<mutex> lock(mtx);
        counter++;
    }
}
\end{lstlisting}

\textbf{要点：} 消除数据竞争的三种方法：(1) 互斥锁（\texttt{mutex}）；(2) 原子操作（\texttt{atomic}）；(3) 避免共享（每个线程拥有自己的数据）。

\subsection*{\textbf{规则 CP.3:} 最小化共享数据的可见范围}

共享数据的范围越小，并发控制越简单，出错的可能性越低。如果一个变量只在一个函数内使用，不要让它成为全局变量或成员变量。

\textbf{错误示例：}

\begin{lstlisting}[language=C++]
class Worker {
    mutex mtx_;
    vector<Task> tasks_;  // 所有方法都能访问
public:
    void add_task(Task t) {
        lock_guard<mutex> lock(mtx_);
        tasks_.push_back(move(t));
    }
    void process_all() {
        lock_guard<mutex> lock(mtx_);
        for (auto& t : tasks_) t.execute();
        tasks_.clear();
    }
};
\end{lstlisting}

\textbf{正确示例：}

\begin{lstlisting}[language=C++]
// 将共享数据限制在最小范围内
class Worker {
    mutex mtx_;
    vector<Task> tasks_;  // 只被 add_task 和 process_all 使用
    // 其他成员不需要锁
};
\end{lstlisting}

\textbf{要点：} 封装是并发安全的基础。将可变状态封装在类内部，通过受控的接口访问，可以大幅减少数据竞争的可能性。

\subsection*{\textbf{规则 CP.4:} 思考并发，而非串行}

在编写并发代码时，必须始终考虑多个线程同时执行的情况。不能假设某个操作是"原子的"——除非它确实是（如单个 \texttt{atomic} 操作）。

\textbf{错误示例：}

\begin{lstlisting}[language=C++]
// "检查然后行动"——不是原子的！
if (!map.contains(key)) {
    map[key] = value;  // 另一个线程可能已经插入了相同的 key
}
\end{lstlisting}

\textbf{正确示例：}

\begin{lstlisting}[language=C++]
// 使用互斥锁保证原子性
{
    lock_guard<mutex> lock(mtx);
    if (!map.contains(key)) {
        map[key] = value;
    }
}
\end{lstlisting}

\textbf{要点：} "检查然后行动"（check-then-act）是并发编程中最常见的错误模式。检查和行动必须在同一个临界区内。

\subsection*{\textbf{规则 CP.8:} 不要使用 \texttt{atomic\_flag} 实现自旋锁}

\texttt{atomic\_flag} 是最底层的原子类型，它不提供等待机制。使用它实现自旋锁会浪费CPU周期。应当使用 \texttt{std::mutex} 或 \texttt{std::atomic} 配合 \texttt{wait}/\texttt{notify}。

\textbf{要点：} 自旋锁只在极短的临界区（几十纳秒级别）中有优势。对于大多数场景，\texttt{std::mutex} 是更好的选择，因为它在无法获取锁时会让出CPU。

\section{线程与锁的 RAII (CP.20--CP.30)}

\subsection*{\textbf{规则 CP.20:} 使用 RAII 管理线程生命周期}

\texttt{std::thread} 的析构函数在 \texttt{joinable()} 为真时会调用 \texttt{std::terminate}。因此，必须确保每个 \texttt{thread} 对象在析构前被 \texttt{join} 或 \texttt{detach}。使用 RAII 封装可以确保安全。

\textbf{错误示例：}

\begin{lstlisting}[language=C++]
void process() {
    thread t(do_work);
    if (some_condition) return;  // t 仍然 joinable！
    // 析构时调用 std::terminate
    t.join();
}
\end{lstlisting}

\textbf{正确示例：}

\begin{lstlisting}[language=C++]
// 使用 scope_guard 或确保所有路径都 join
void process() {
    thread t(do_work);
    // 使用 RAII 确保 join
    auto guard = std::scope_exit([&t] { t.join(); });
    if (some_condition) return;  // guard 确保 join
}

// 或者更简单：使用 jthread（C++20）
void process() {
    jthread t(do_work);  // 析构时自动 join
    if (some_condition) return;  // 安全
}
\end{lstlisting}

\textbf{要点：} C++20的 \texttt{jthread} 是管理线程生命周期的推荐方式。它在析构时自动 \texttt{request\_stop} 并 \texttt{join}。

\subsection*{\textbf{规则 CP.21:} 使用 \texttt{lock\_guard} 或 \texttt{unique\_lock} 管理互斥锁}

手动调用 \texttt{lock()} 和 \texttt{unlock()} 是危险的——任何提前返回或异常都会导致锁未释放。RAII 锁管理器确保锁在作用域退出时自动释放。

\textbf{错误示例：}

\begin{lstlisting}[language=C++]
mutex mtx;

void process() {
    mtx.lock();
    if (error_condition) return;  // 忘记 unlock！
    // ...
    mtx.unlock();
}
\end{lstlisting}

\textbf{正确示例：}

\begin{lstlisting}[language=C++]
void process() {
    lock_guard<mutex> lock(mtx);
    if (error_condition) return;  // 自动 unlock
    // ...
}  // 自动 unlock
\end{lstlisting}

\textbf{要点：} \texttt{lock\_guard} 是最简单的RAII锁，适用于锁持有整个作用域的场景。\texttt{unique\_lock} 更灵活，支持延迟加锁、解锁、重新加锁，适用于需要更细粒度控制的场景（如条件变量）。

\subsection*{\textbf{规则 CP.22:} 不要持有锁调用未知代码}

如果在持有锁的情况下调用了回调函数或虚函数，你不知道该代码是否会尝试获取同一个锁（导致死锁）或其他锁（导致锁顺序反转）。

\textbf{错误示例：}

\begin{lstlisting}[language=C++]
void process(Widget& w, function<void(Widget&)> callback) {
    lock_guard<mutex> lock(mtx);
    w.update();
    callback(w);  // 如果 callback 尝试获取 mtx？死锁！
}
\end{lstlisting}

\textbf{正确示例：}

\begin{lstlisting}[language=C++]
void process(Widget& w, function<void(Widget&)> callback) {
    {
        lock_guard<mutex> lock(mtx);
        w.update();
    }  // 释放锁
    callback(w);  // 在无锁状态下调用
}
\end{lstlisting}

\textbf{要点：} 持有锁时只操作你知道的代码。如果需要调用外部代码，先释放锁，或使用"拷贝-操作"模式（在锁内拷贝数据，在锁外操作拷贝）。

\subsection*{\textbf{规则 CP.23:} 不要从持有锁的作用域中 \texttt{return}}

这与 CP.21 相关。使用 RAII 锁管理器后，\texttt{return} 是安全的（锁会自动释放）。但如果不使用 RAII，\texttt{return} 前忘记 \texttt{unlock} 是常见错误。

\textbf{要点：} 始终使用 RAII 锁管理器。这样 \texttt{return}、\texttt{break}、\texttt{continue}、异常抛出都是安全的。

\subsection*{\textbf{规则 CP.24:} 检查 \texttt{thread} 对象是否 \texttt{joinable}}

在 \texttt{join} 或 \texttt{detach} 之前，应当检查线程是否 \texttt{joinable}。对不可 \texttt{joinable} 的线程调用 \texttt{join} 或 \texttt{detach} 是未定义行为。

\textbf{正确示例：}

\begin{lstlisting}[language=C++]
thread t(do_work);
// ...
if (t.joinable()) {
    t.join();
}
\end{lstlisting}

\textbf{要点：} 已经被 \texttt{join}、\texttt{detach} 或默认构造的 \texttt{thread} 对象是不可 \texttt{joinable} 的。

\section{最小化锁竞争 (CP.40--CP.50)}

\subsection*{\textbf{规则 CP.40:} 最小化临界区}

锁持有的时间越短，其他线程等待的时间越短，程序的并发性越好。将不需要共享访问的计算移出临界区。

\textbf{错误示例：}

\begin{lstlisting}[language=C++]
void process(Data& data) {
    lock_guard<mutex> lock(mtx);
    data.update();        // 需要锁
    auto result = compute(data);  // 不需要锁
    save_to_file(result);  // 不需要锁
}
\end{lstlisting}

\textbf{正确示例：}

\begin{lstlisting}[language=C++]
void process(Data& data) {
    Data snapshot;
    {
        lock_guard<mutex> lock(mtx);
        data.update();
        snapshot = data;  // 拷贝需要的数据
    }  // 释放锁
    auto result = compute(snapshot);  // 无锁计算
    save_to_file(result);             // 无锁IO
}
\end{lstlisting}

\textbf{要点：} 临界区应当只包含必须串行化的操作。计算、IO等不需要共享访问的操作应当在临界区外进行。

\subsection*{\textbf{规则 CP.41:} 最小化锁的持有时间}

这与 CP.40 类似，但更强调时间维度。即使临界区内的操作确实需要锁，也应当尽量减少操作的数量和复杂度。

\textbf{要点：} 一个好的经验法则是：临界区内只做"赋值"和"简单计算"，不做IO、不做网络请求、不做复杂算法。

\subsection*{\textbf{规则 CP.42:} 不要在没有锁的情况下 \texttt{wait}}

使用条件变量时，必须在持有锁的情况下调用 \texttt{wait}。\texttt{wait} 会原子地释放锁并等待，被唤醒后重新获取锁。

\textbf{错误示例：}

\begin{lstlisting}[language=C++]
// 不在锁内 wait——可能错过通知
if (!ready) {
    cv.wait(lk);  // 如果没有锁保护，可能在检查后、wait前收到通知
}
\end{lstlisting}

\textbf{正确示例：}

\begin{lstlisting}[language=C++]
unique_lock<mutex> lk(mtx);
cv.wait(lk, [] { return ready; });  // 原子地检查并等待
\end{lstlisting}

\textbf{要点：} 条件变量的 \texttt{wait} 必须在 \texttt{unique\_lock} 保护下调用。使用带谓词版本的 \texttt{wait} 可以同时防止虚假唤醒（spurious wakeup）和通知丢失。

\subsection*{\textbf{规则 CP.43:} 注意锁的获取顺序}

当需要获取多个锁时，所有线程必须以相同的顺序获取它们，否则可能发生死锁。

\textbf{错误示例：}

\begin{lstlisting}[language=C++]
// 线程1: 先锁A，再锁B
void thread1() {
    lock_guard<mutex> lockA(mtxA);
    lock_guard<mutex> lockB(mtxB);
    // ...
}

// 线程2: 先锁B，再锁A——死锁！
void thread2() {
    lock_guard<mutex> lockB(mtxB);
    lock_guard<mutex> lockA(mtxA);
    // ...
}
\end{lstlisting}

\textbf{正确示例：}

\begin{lstlisting}[language=C++]
// 使用 std::lock 同时获取多个锁
void thread1() {
    unique_lock<mutex> lockA(mtxA, defer_lock);
    unique_lock<mutex> lockB(mtxB, defer_lock);
    lock(lockA, lockB);  // 无死锁地获取两个锁
}

void thread2() {
    unique_lock<mutex> lockA(mtxA, defer_lock);
    unique_lock<mutex> lockB(mtxB, defer_lock);
    lock(lockA, lockB);  // 同样的顺序
}
\end{lstlisting}

\textbf{要点：} \texttt{std::lock} 和 \texttt{std::scoped\_lock}（C++17）可以无死锁地获取多个锁。如果必须手动获取多个锁，定义全局的锁获取顺序并严格遵守。

\subsection*{\textbf{规则 CP.44:} 记得给 \texttt{lock\_guard} 的 \texttt{mutex} 参数命名}

这是一个微妙但重要的陷阱：匿名临时 \texttt{lock\_guard} 会立即析构，等于没有加锁。

\textbf{错误示例：}

\begin{lstlisting}[language=C++]
void process() {
    lock_guard<mutex>(mtx);  // 匿名临时对象！立即析构
    // 实际上没有加锁
    shared_data.modify();
}
\end{lstlisting}

\textbf{正确示例：}

\begin{lstlisting}[language=C++]
void process() {
    lock_guard<mutex> guard(mtx);  // 命名对象，持有到作用域结束
    shared_data.modify();
}
\end{lstlisting}

\textbf{要点：} 这是C++中"最令人烦恼的解析"（most vexing parse）的变体。始终给 RAII 对象命名。一些编译器（如Clang）可以在开启警告时检测这个问题。

\subsection*{\textbf{规则 CP.50:} 使用 \texttt{scoped\_lock} 替代 \texttt{lock\_guard}}

\texttt{scoped\_lock}（C++17）是 \texttt{lock\_guard} 的升级版，支持同时管理多个互斥锁，并使用 \texttt{std::lock} 算法避免死锁。

\textbf{正确示例：}

\begin{lstlisting}[language=C++]
void transfer(Account& from, Account& to, double amount) {
    // 同时锁定两个账户的互斥锁，无死锁风险
    scoped_lock lock(from.mtx, to.mtx);
    from.balance -= amount;
    to.balance += amount;
}
\end{lstlisting}

\textbf{要点：} 在新代码中，应当默认使用 \texttt{scoped\_lock} 而非 \texttt{lock\_guard}。它更安全、更灵活。

\section{Future/Promise 模式 (CP.60--CP.70)}

\subsection*{\textbf{规则 CP.60:} 使用 \texttt{future} 获取异步操作的结果}

\texttt{std::future} 和 \texttt{std::async} 提供了高级的异步编程接口，比直接使用线程更简单、更安全。

\textbf{错误示例：}

\begin{lstlisting}[language=C++]
// 手动管理线程和共享变量
int result;
mutex mtx;
thread t([&] {
    lock_guard<mutex> lock(mtx);
    result = compute();
});
t.join();
cout << result;
\end{lstlisting}

\textbf{正确示例：}

\begin{lstlisting}[language=C++]
auto result = async(launch::async, compute);
cout << result.get();  // 等待并获取结果
\end{lstlisting}

\textbf{要点：} \texttt{async} + \texttt{future} 消除了手动管理线程、互斥锁和共享变量的需要。它自动处理异常传播——如果异步操作抛出异常，\texttt{get()} 会重新抛出。

\subsection*{\textbf{规则 CP.61:} 使用 \texttt{promise} 建立 \texttt{future}}

\texttt{std::promise} 和 \texttt{std::future} 是一对配合使用的工具。\texttt{promise} 用于设置值，\texttt{future} 用于获取值。

\textbf{正确示例：}

\begin{lstlisting}[language=C++]
promise<int> prom;
future<int> fut = prom.get_future();

thread t([&prom] {
    int result = compute();
    prom.set_value(result);  // 设置值
});

cout << fut.get();  // 等待并获取值
t.join();
\end{lstlisting}

\textbf{要点：} \texttt{promise/future} 适合生产者-消费者模式和跨线程通信。每个 \texttt{promise} 只能设置一次值，多次设置会抛出异常。

\subsection*{\textbf{规则 CP.62:} 使用 \texttt{ packaged\_task} 包装可调用对象}

\texttt{packaged\_task} 将一个可调用对象与 \texttt{future} 绑定，使得可调用对象的结果可以通过 \texttt{future} 获取。如果可调用对象抛出异常，异常会通过 \texttt{future} 传播。

\textbf{正确示例：}

\begin{lstlisting}[language=C++]
packaged_task<int(int, int)> task([](int a, int b) {
    return a + b;
});
future<int> fut = task.get_future();
thread(move(task), 3, 4);  // 在另一个线程中执行
cout << fut.get();  // 输出 7
\end{lstlisting}

\textbf{要点：} \texttt{packaged\_task} 适合在线程池中分发任务。它将任务与结果通道绑定，简化了异步任务的管理。

\subsection*{\textbf{规则 CP.100:} 不要使用 \texttt{volatile} 进行线程同步}

\texttt{volatile} 在C++中的语义与Java或C\#不同。C++的 \texttt{volatile} 不提供任何线程同步保证，不创建内存屏障。对于原子操作，应当使用 \texttt{std::atomic}。

\textbf{错误示例：}

\begin{lstlisting}[language=C++]
volatile bool ready = false;

// 线程1
while (!ready) {}  // 自旋等待——不保证看到线程2的写入

// 线程2
data = compute();
ready = true;  // 不保证在线程1之前可见
\end{lstlisting}

\textbf{正确示例：}

\begin{lstlisting}[language=C++]
atomic<bool> ready{false};

// 线程1
while (!ready.load(memory_order_acquire)) {}
// 保证能看到 data 的值

// 线程2
data = compute();
ready.store(true, memory_order_release);  // 保证 data 在此之前可见
\end{lstlisting}

\textbf{要点：} C++中的 \texttt{volatile} 仅用于特殊内存地址（如内存映射IO），不用于线程同步。线程同步应当使用 \texttt{std::atomic} 和内存序。

\section*{练习题}

\begin{enumerate}
\item 以下代码有什么并发问题？如何修复？
\begin{lstlisting}[language=C++]
class Counter {
    int count_ = 0;
public:
    void increment() { count_++; }
    int get() const { return count_; }
};
\end{lstlisting}

\item 实现一个线程安全的单例模式，使用 \texttt{std::call\_once} 和 \texttt{std::once\_flag}。

\item 解释以下代码为什么可能死锁，并给出两种不同的修复方案：
\begin{lstlisting}[language=C++]
void transfer(Account& a, Account& b, double amount) {
    lock_guard<mutex> lockA(a.mtx);
    lock_guard<mutex> lockB(b.mtx);
    a.balance -= amount;
    b.balance += amount;
}
// 线程1: transfer(A, B, 100);
// 线程2: transfer(B, A, 200);
\end{lstlisting}

\item 使用 \texttt{async} 和 \texttt{future} 重写以下代码：
\begin{lstlisting}[language=C++]
vector<int> results(1000);
vector<thread> threads;
for (int i = 0; i < 1000; ++i) {
    threads.emplace_back([&results, i] {
        results[i] = compute(i);
    });
}
for (auto& t : threads) t.join();
\end{lstlisting}

\item 设计一个生产者-消费者队列，使用 \texttt{mutex}、\texttt{condition\_variable} 和 \texttt{queue}。要求支持多生产者和多消费者。
\end{enumerate}

% ============================================================
% Chapter 14: C风格编程与源文件
% ============================================================

\chapter{C风格编程与源文件}

C++脱胎于C，两者共享大量语法和底层特性。然而，C++提供了更高级的抽象机制（类、模板、异常、RAII等），使得许多C风格的编程模式在C++中变得不必要甚至有害。Core Guidelines 的"CPL"（C-style programming）系列规则帮助开发者识别并消除C风格的代码，转而使用更安全的C++特性。"SF"（Source files）系列规则则关注源文件的组织和编译。

\section{C风格编程规则 (CPL.1--CPL.10)}

\subsection*{\textbf{规则 CPL.1:} 最小化C风格代码}

C++中应当优先使用C++特性而非C特性。C风格的代码（如裸数组、C字符串、\texttt{malloc}、\texttt{printf}）缺乏类型安全和资源管理保证，在C++中有更好的替代方案。

\textbf{错误示例：}

\begin{lstlisting}[language=C++]
// C风格
char* buf = (char*)malloc(256);
sprintf(buf, "Hello, %s! You are %d years old.", name, age);
printf("%s\n", buf);
free(buf);
\end{lstlisting}

\textbf{正确示例：}

\begin{lstlisting}[language=C++]
// C++风格
string buf = format("Hello, {}! You are {} years old.", name, age);
cout << buf << endl;
\end{lstlisting}

\textbf{要点：} C++替代C风格的对照表：\texttt{malloc/free} $\rightarrow$ \texttt{new/delete}（通过智能指针）或容器；\texttt{printf} $\rightarrow$ \texttt{cout} 或 \texttt{format}；裸数组 $\rightarrow$ \texttt{array}/\texttt{vector}；C字符串 $\rightarrow$ \texttt{string}/\texttt{string\_view}；\texttt{typedef} $\rightarrow$ \texttt{using}。

\subsection*{\textbf{规则 CPL.2:} 避免C风格的联合（union）}

C风格的联合体（union）不提供类型安全——它不知道当前存储的是哪种类型。读取一个不是最近写入的成员是未定义行为。C++11引入了"可辨识联合"（discriminated union）\texttt{std::variant}。

\textbf{错误示例：}

\begin{lstlisting}[language=C++]
union Value {
    int i;
    double d;
    char* s;
};

Value v;
v.i = 42;
cout << v.d;  // 未定义行为！读取的不是最近写入的成员
\end{lstlisting}

\textbf{正确示例：}

\begin{lstlisting}[language=C++]
variant<int, double, string> value;
value = 42;
cout << get<double>(value);  // 抛出 bad_variant_access
// 安全的方式
if (auto* p = get_if<int>(&value))
    cout << *p;
\end{lstlisting}

\textbf{要点：} \texttt{std::variant} 在运行时跟踪当前存储的类型，提供类型安全的访问。它是C风格联合的安全替代。

\subsection*{\textbf{规则 CPL.3:} 避免C风格的枚举}

C风格的枚举（\texttt{enum}）将枚举值暴露在 enclosing scope 中，且允许隐式转换为整数。这导致名称冲突和类型不安全的比较。C++11的 \texttt{enum class} 解决了这些问题。

\textbf{错误示例：}

\begin{lstlisting}[language=C++]
enum Color { red, green, blue };
enum TrafficLight { red, yellow, green };  // 错误：red 和 green 重复定义！

Color c = red;
int x = c;  // 隐式转换为整数——类型不安全
\end{lstlisting}

\textbf{正确示例：}

\begin{lstlisting}[language=C++]
enum class Color { red, green, blue };
enum class TrafficLight { red, yellow, green };  // OK，不冲突

Color c = Color::red;
int x = static_cast<int>(c);  // 需要显式转换
\end{lstlisting}

\textbf{要点：} \texttt{enum class} 提供作用域隔离（枚举值不会泄漏到外部作用域）和强类型（不能隐式转换为整数）。在新代码中应当始终使用 \texttt{enum class}。

\subsection*{\textbf{规则 CPL.4:} 避免可变参数函数（\texttt{...}）}

C风格的可变参数函数（如 \texttt{printf}）不是类型安全的。编译器无法检查传入参数的类型和数量是否匹配。C++提供了更好的替代方案。

\textbf{错误示例：}

\begin{lstlisting}[language=C++]
void log(const char* fmt, ...) {
    va_list args;
    va_start(args, fmt);
    vprintf(fmt, args);  // 如果格式字符串与参数不匹配？未定义行为
    va_end(args);
}

log("Value: %d, Name: %s", 42);  // 缺少参数，未定义行为
\end{lstlisting}

\textbf{正确示例：}

\begin{lstlisting}[language=C++]
// 使用模板可变参数
template<typename... Args>
void log(string_view fmt, Args&&... args) {
    cout << format(fmt, forward<Args>(args)...);
}

log("Value: {}, Name: {}", 42, "hello");  // 类型安全
\end{lstlisting}

\textbf{要点：} C++的可变参数模板（variadic templates）是类型安全的替代方案。\texttt{std::format}（C++20）提供了类型安全的字符串格式化。

\subsection*{\textbf{规则 CPL.5:} 不要使用 \texttt{setjmp}/\texttt{longjmp}}

\texttt{setjmp}/\texttt{longjmp} 是C风格的非局部跳转机制。它们绕过析构函数，破坏RAII，导致资源泄漏和未定义行为。在C++中应当使用异常进行非局部控制流转移。

\textbf{错误示例：}

\begin{lstlisting}[language=C++]
jmp_buf env;

void deep_function() {
    // ... 一些操作 ...
    longjmp(env, 1);  // 跳过所有析构函数！
}

void caller() {
    Widget w;  // w 的析构函数不会被调用
    if (setjmp(env) == 0) {
        deep_function();
    }
}
\end{lstlisting}

\textbf{正确示例：}

\begin{lstlisting}[language=C++]
void deep_function() {
    // ... 一些操作 ...
    throw MyError();  // 栈展开会正确调用析构函数
}

void caller() {
    Widget w;  // w 的析构函数会被正确调用
    try {
        deep_function();
    } catch (const MyError&) {
        // 处理错误
    }
}
\end{lstlisting}

\textbf{要点：} \texttt{setjmp}/\texttt{longjmp} 完全不了解C++的对象模型。它们与RAII、异常处理、构造/析构语义完全不兼容。在C++中绝对不应当使用。

\subsection*{\textbf{规则 CPL.6--CPL.10:} 更多C风格替代规则}

\textbf{CPL.6:} 使用 \texttt{std::array} 或 \texttt{std::vector} 替代C风格数组。C风格数组不知道自己的大小，会退化为指针，容易越界访问。

\textbf{CPL.7:} 使用 \texttt{std::string} 或 \texttt{std::string\_view} 替代C风格字符串（\texttt{char*}）。C风格字符串需要手动管理内存，容易缓冲区溢出。

\textbf{CPL.8:} 使用 \texttt{using} 替代 \texttt{typedef}。\texttt{using} 支持模板别名，语法更清晰。

\textbf{错误示例：}

\begin{lstlisting}[language=C++]
typedef void (*FuncPtr)(int, int);
typedef vector<int> IntVec;
\end{lstlisting}

\textbf{正确示例：}

\begin{lstlisting}[language=C++]
using FuncPtr = void (*)(int, int);
using IntVec = vector<int>;

// using 支持模板别名
template<typename T>
using Vec = vector<T>;
\end{lstlisting}

\textbf{CPL.9:} 使用 \texttt{constexpr} 替代 \texttt{\#define} 定义编译期常量。\texttt{constexpr} 有作用域、有类型、可调试。

\textbf{CPL.10:} 使用 \texttt{static\_cast}、\texttt{dynamic\_cast} 等替代C风格强制类型转换。

\textbf{要点：} 每一条CPL规则的核心思想是：C++提供了比C更安全、更表达力强的替代方案。使用它们不仅使代码更安全，也使代码更清晰地表达意图。

\section{源文件规则 (SF.1--SF.10)}

\subsection*{\textbf{规则 SF.1:} 使用 \texttt{.cpp} 后缀命名要被编译的文件}

源文件的扩展名应当让编译器和构建系统知道如何处理它。\texttt{.cpp} 是最广泛使用的C++源文件扩展名。

\textbf{要点：} 保持一致的文件命名约定。\texttt{.cpp}、\texttt{.cc}、\texttt{.cxx} 都是常见的选择，但在一个项目中应当统一。

\subsection*{\textbf{规则 SF.2:} 头文件应当有保护符（include guard）}

头文件可能被多次包含（通过不同的包含路径）。没有保护符会导致重复定义错误。

\textbf{错误示例：}

\begin{lstlisting}[language=C++]
// widget.h —— 没有保护符
class Widget {
    // ...
};
\end{lstlisting}

\textbf{正确示例：}

\begin{lstlisting}[language=C++]
// widget.h
#ifndef WIDGET_H
#define WIDGET_H

class Widget {
    // ...
};

#endif

// 或使用 #pragma once（广泛支持但非标准）
#pragma once

class Widget {
    // ...
};
\end{lstlisting}

\textbf{要点：} 每个头文件都应当有保护符。使用 \texttt{\#ifndef} 风格的保护符是标准兼容的做法；\texttt{\#pragma once} 更简洁但非标准。两者选一种并在项目中统一。

\subsection*{\textbf{规则 SF.3:} 使用头文件只用于声明，不用于定义}

头文件中只应当包含声明（函数声明、类声明、模板声明），不应当包含非内联函数的定义或全局变量的定义。否则，多个源文件包含同一头文件会导致链接错误（ODR违反）。

\textbf{错误示例：}

\begin{lstlisting}[language=C++]
// utils.h
void utility_function() {  // 非内联函数的定义
    // ...
}

int global_counter = 0;  // 全局变量的定义
\end{lstlisting}

\textbf{正确示例：}

\begin{lstlisting}[language=C++]
// utils.h
inline void utility_function() {  // inline 允许在头文件中定义
    // ...
}

extern int global_counter;  // 声明，定义在 .cpp 中

// 或 C++17 的 inline 变量
inline int global_counter = 0;  // C++17 允许
\end{lstlisting}

\textbf{要点：} 头文件中的定义（非 \texttt{inline}）会导致多重定义错误。记住：声明可以出现多次，定义只能出现一次（ODR，One Definition Rule）。

\subsection*{\textbf{规则 SF.4:} 在使用前包含头文件}

头文件应当按照一定的顺序包含：先包含当前模块对应的头文件，然后按标准库、第三方库、项目内头文件的顺序包含。

\textbf{正确示例：}

\begin{lstlisting}[language=C++]
// widget.cpp
#include "widget.h"      // 当前模块的头文件（最先包含）
#include <vector>         // 标准库
#include <string>         // 标准库
#include <boost/...>      // 第三方库
#include "utils.h"        // 项目内头文件
\end{lstlisting}

\textbf{要点：} 将当前模块的头文件放在最前面包含，可以检测到头文件是否自包含（是否依赖于其他头文件的隐式包含）。

\subsection*{\textbf{规则 SF.5:} \texttt{.cpp} 文件必须包含其对应的头文件}

这与 SF.4 相关。\texttt{.cpp} 文件应当首先包含它声明接口的头文件，以确保头文件是自包含的。

\textbf{要点：} 如果 \texttt{widget.cpp} 不首先包含 \texttt{widget.h}，那么 \texttt{widget.h} 可能依赖于其他头文件的隐式包含，这在其他文件单独包含 \texttt{widget.h} 时会失败。

\subsection*{\textbf{规则 SF.6:} 使用 \texttt{using namespace} 时遵循作用域规则}

这在 ES.12 中已经讨论。在头文件中绝不使用 \texttt{using namespace}；在 \texttt{.cpp} 文件中谨慎使用；在函数作用域内使用是安全的。

\subsection*{\textbf{规则 SF.7:} 不要在包含头文件时使用 \texttt{using} 声明}

在头文件中使用 \texttt{using} 声明会污染包含该头文件的所有翻译单元。

\textbf{错误示例：}

\begin{lstlisting}[language=C++]
// mylib.h
using std::vector;  // 所有包含此头文件的文件都受影响

class MyContainer {
    vector<int> data;
};
\end{lstlisting}

\textbf{正确示例：}

\begin{lstlisting}[language=C++]
// mylib.h
class MyContainer {
    std::vector<int> data;  // 明确使用 std::
};
\end{lstlisting}

\textbf{要点：} 头文件中的 \texttt{using} 声明和 \texttt{using namespace} 都是危险的。它们的影响范围是所有包含该头文件的文件，可能远超作者的预期。

\subsection*{\textbf{规则 SF.8:} 使用 \texttt{\#include} 包含标准库头文件}

使用标准库头文件时不要加路径前缀或 \texttt{.h} 后缀。C++标准库头文件不带 \texttt{.h} 后缀（如 \texttt{<vector>} 而非 \texttt{<vector.h>}）。

\textbf{错误示例：}

\begin{lstlisting}[language=C++]
#include <vector.h>     // 错误！C++标准库头文件不带 .h
#include <stdio.h>      // C风格头文件，应当用 <cstdio>
#include "mylib/vector" // 不需要路径前缀
\end{lstlisting}

\textbf{正确示例：}

\begin{lstlisting}[language=C++]
#include <vector>       // 正确
#include <cstdio>       // C++风格的C头文件
#include "mylib.h"      // 项目内头文件使用引号
\end{lstlisting}

\textbf{要点：} C++标准库头文件将C库函数放在 \texttt{std} 命名空间中（如 \texttt{std::printf} 在 \texttt{<cstdio>} 中）。使用C++风格的头文件可以获得更好的类型安全和重载支持。

\subsection*{\textbf{规则 SF.10:} 使用模板时，将模板定义放在头文件中}

模板的定义（而非仅声明）必须在每个使用它的翻译单元中可见。这意味着模板的定义通常放在头文件中。

\textbf{错误示例：}

\begin{lstlisting}[language=C++]
// container.h
template<typename T>
class Container {
public:
    void add(T value);  // 只有声明
};

// container.cpp
template<typename T>
void Container<T>::add(T value) {  // 定义在 .cpp 中
    // ...
}
// 链接错误！其他文件使用 Container<int> 时找不到 add 的定义
\end{lstlisting}

\textbf{正确示例：}

\begin{lstlisting}[language=C++]
// container.h
template<typename T>
class Container {
public:
    void add(T value) {  // 定义在头文件中
        // ...
    }
};

// 或者使用显式实例化
// container.cpp
template class Container<int>;  // 显式实例化
\end{lstlisting}

\textbf{要点：} 模板定义需要在每个使用它的翻译单元中可见。最简单的方式是将定义放在头文件中。如果模板很大，可以使用 \texttt{inline} 关键字或将实现放在 \texttt{\_impl.h} 文件中。

\section*{练习题}

\begin{enumerate}
\item 将以下C风格代码改写为现代C++风格：
\begin{lstlisting}[language=C++]
typedef struct {
    char name[64];
    int age;
    double salary;
} Employee;

void print_employee(Employee* e) {
    printf("Name: %s, Age: %d, Salary: %.2f\n",
           e->name, e->age, e->salary);
}

Employee* create_employee(const char* name, int age, double salary) {
    Employee* e = (Employee*)malloc(sizeof(Employee));
    strcpy(e->name, name);
    e->age = age;
    e->salary = salary;
    return e;
}
\end{lstlisting}

\item 解释以下代码为什么会导致链接错误，并给出修复方案：
\begin{lstlisting}[language=C++]
// config.h
int max_connections = 100;
string default_host = "localhost";
\end{lstlisting}

\item 为以下头文件添加正确的 include guard：
\begin{lstlisting}[language=C++]
// matrix.h
class Matrix {
    vector<vector<double>> data_;
public:
    Matrix(int rows, int cols);
    double& at(int r, int c);
};
\end{lstlisting}

\item 解释为什么以下代码违反了ODR（One Definition Rule），并给出修复方案：
\begin{lstlisting}[language=C++]
// utils.h
void helper() {
    cout << "helper called" << endl;
}
\end{lstlisting}

\item 设计一个项目的头文件组织结构。包括：(a) 目录结构；(b) 命名约定；(c) include guard 策略；(d) include 顺序规则。
\end{enumerate}

% ============================================================
% Chapter 15: 标准库 (Standard Library - "SL" section)
% ============================================================

\chapter{标准库}

C++标准库是语言最强大的武器之一。它提供了经过广泛测试和优化的容器、算法、字符串处理、IO、并发等基础设施。Core Guidelines 的"SL"系列规则的核心信息是：优先使用标准库，而不是自己造轮子。标准库代码是跨平台的、类型安全的、性能经过调优的，且所有C++开发者都熟悉它。

\section{通用标准库规则 (SL.1--SL.10)}

\subsection*{\textbf{规则 SL.1:} 使用标准库}

不要重新发明轮子。标准库提供了大量经过验证的组件：容器（\texttt{vector}、\texttt{map}、\texttt{unordered\_map} 等）、算法（\texttt{sort}、\texttt{find}、\texttt{transform} 等）、字符串（\texttt{string}、\texttt{string\_view}）、智能指针（\texttt{unique\_ptr}、\texttt{shared\_ptr}）、并发原语（\texttt{thread}、\texttt{mutex}、\texttt{atomic}）等。使用标准库可以节省开发时间、减少bug、提高代码可维护性。

\textbf{错误示例：}

\begin{lstlisting}[language=C++]
// 自己实现链表
template<typename T>
struct ListNode {
    T data;
    ListNode* next;
    ListNode* prev;
};

template<typename T>
class MyList {
    ListNode<T>* head_;
    // ... 几百行容易出错的实现
};
\end{lstlisting}

\textbf{正确示例：}

\begin{lstlisting}[language=C++]
// 使用标准库
list<int> my_list;
my_list.push_back(42);
\end{lstlisting}

\textbf{要点：} 标准库的每个组件都经过了数十年的审查、优化和实际使用验证。自己实现的版本几乎不可能在所有边界情况下都正确。

\subsection*{\textbf{规则 SL.2:} 优先使用标准库组件而非第三方组件}

当标准库提供了所需功能时，不应当引入第三方库。标准库是免费的（无需额外依赖）、跨平台的、编译器集成的。引入第三方库增加了构建复杂度、依赖管理成本和潜在的兼容性问题。

\textbf{错误示例：}

\begin{lstlisting}[language=C++]
// 使用第三方库实现基本功能
#include <boost/optional.hpp>  // C++17 已有 std::optional
#include <boost/variant.hpp>   // C++17 已有 std::variant
\end{lstlisting}

\textbf{正确示例：}

\begin{lstlisting}[language=C++]
#include <optional>
#include <variant>
\end{lstlisting}

\textbf{要点：} 第三方库在标准库不提供所需功能时是合理的选择（如JSON解析、网络框架、GUI库等）。但对于标准库已经覆盖的功能，应当优先使用标准库。

\subsection*{\textbf{规则 SL.3:} 不要给标准库添加特化或重载}

给 \texttt{std} 命名空间添加自己的特化或重载是未定义行为（除了标准明确允许的特化，如 \texttt{std::hash}、\texttt{std::swap} 等）。这样做可能与未来标准库版本冲突。

\textbf{错误示例：}

\begin{lstlisting}[language=C++]
namespace std {
    // 未定义行为！给 vector 添加自定义方法
    template<typename T>
    void vector<T>::custom_method() { /* ... */ }
}
\end{lstlisting}

\textbf{正确示例：}

\begin{lstlisting}[language=C++]
// 使用自由函数
template<typename T>
void custom_method(vector<T>& v) {
    // ...
}

// 或使用继承（如果确实需要扩展）
template<typename T>
class MyVector : public vector<T> {
    // 注意：vector 没有虚析构函数，不推荐继承
};
\end{lstlisting}

\textbf{要点：} 标准允许的特化包括：\texttt{std::hash}（用于自定义类型的哈希）、\texttt{std::swap}（用于自定义类型的交换）、\texttt{std::iterator\_traits} 等。其他情况下不要向 \texttt{std} 命名空间添加内容。

\subsection*{\textbf{规则 SL.4:} 安全地使用标准库类型}

标准库类型在正确使用下是安全的。但某些误用（如迭代器失效、越界访问）仍然可能导致未定义行为。了解并遵守每种类型的安全使用规则是必要的。

\textbf{错误示例：}

\begin{lstlisting}[language=C++]
vector<int> v = {1, 2, 3, 4, 5};
auto it = v.begin();
v.push_back(6);  // 可能使 it 失效！
cout << *it;     // 未定义行为
\end{lstlisting}

\textbf{正确示例：}

\begin{lstlisting}[language=C++]
vector<int> v = {1, 2, 3, 4, 5};
v.push_back(6);
auto it = v.begin();  // 在修改之后获取迭代器
cout << *it;          // 安全
\end{lstlisting}

\textbf{要点：} 了解迭代器失效规则是安全使用标准库容器的关键。\texttt{vector} 在重新分配内存时会使所有迭代器失效；\texttt{map} 在插入时不会使已有迭代器失效；\texttt{deque} 在两端插入时会使所有迭代器失效。

\subsection*{\textbf{规则 SL.5:} 优先使用 \texttt{vector} 而非其他容器}

除非有明确的理由使用其他容器，否则默认使用 \texttt{vector}。\texttt{vector} 具有最好的缓存局部性（连续内存），最快的随机访问，且在尾部插入的均摊时间复杂度为 $O(1)$。在许多场景下，\texttt{vector} 的性能优于看似更"合适"的容器（如 \texttt{list}、\texttt{map}）。

\textbf{错误示例：}

\begin{lstlisting}[language=C++]
// "我需要频繁在中间插入，所以用 list"
list<int> data;
// 实际上由于缓存不友好，list 在大多数场景下比 vector 慢
\end{lstlisting}

\textbf{正确示例：}

\begin{lstlisting}[language=C++]
// 默认使用 vector
vector<int> data;
// 只有当 profiler 确认 list 更快时才改用 list
\end{lstlisting}

\textbf{要点：} "vector is king"是现代C++性能优化的经验法则。在绝大多数场景下，\texttt{vector} 的连续内存布局带来的缓存优势超过了其"低效"操作（如中间插入）的理论复杂度劣势。

\subsection*{\textbf{规则 SL.6:} 使用 \texttt{array} 替代C风格数组}

\texttt{std::array} 是C风格数组的安全替代。它知道自身大小，不会退化为指针，支持标准算法和范围 \texttt{for}。

\textbf{错误示例：}

\begin{lstlisting}[language=C++]
int arr[10];
// arr 退化为 int*，丢失大小信息
// 无法使用 arr.size()
// 容易越界访问
\end{lstlisting}

\textbf{正确示例：}

\begin{lstlisting}[language=C++]
array<int, 10> arr;
cout << arr.size();  // 10
for (auto& x : arr) { /* ... */ }  // 范围 for
arr.at(15) = 42;  // 运行时边界检查
\end{lstlisting}

\textbf{要点：} \texttt{std::array} 与C风格数组有相同的内存布局和性能（零开销抽象），但提供了类型安全的接口。

\subsection*{\textbf{规则 SL.7:} 使用 \texttt{string} 替代C风格字符串}

\texttt{std::string} 自动管理内存，知道自身长度，支持拼接、比较、查找等操作。它消除了缓冲区溢出、内存泄漏等C风格字符串的常见问题。

\textbf{错误示例：}

\begin{lstlisting}[language=C++]
char buf[256];
strcpy(buf, user_input);  // 如果 user_input 超过 255 字符？缓冲区溢出！
strcat(buf, more_data);   // 如果总长度超过 255？缓冲区溢出！
\end{lstlisting}

\textbf{正确示例：}

\begin{lstlisting}[language=C++]
string buf = user_input;  // 自动管理内存
buf += more_data;         // 自动扩展
\end{lstlisting}

\textbf{要点：} \texttt{std::string} 是现代C++中字符串操作的默认选择。只有在与C API交互时才需要 \texttt{c\_str()}。

\subsection*{\textbf{规则 SL.8:} 避免C风格的IO（\texttt{printf}/\texttt{scanf}）}

C风格的IO函数（\texttt{printf}、\texttt{scanf} 等）不是类型安全的。格式字符串与参数不匹配是未定义行为。C++提供了类型安全的IO流（\texttt{iostream}）和 \texttt{std::format}（C++20）。

\textbf{错误示例：}

\begin{lstlisting}[language=C++]
int x = 42;
printf("%s\n", x);  // 格式不匹配！未定义行为
\end{lstlisting}

\textbf{正确示例：}

\begin{lstlisting}[language=C++]
int x = 42;
cout << x << endl;  // 类型安全

// 或 C++20
cout << format("{}\n", x);  // 类型安全
\end{lstlisting}

\textbf{要点：} \texttt{printf} 的性能通常优于 \texttt{iostream}，但 \texttt{std::format}（C++20）兼具类型安全和高性能。在C++20可用时，\texttt{std::format} 是最佳选择。

\subsection*{\textbf{规则 SL.9:} 不要使用 \texttt{std::auto\_ptr}}

\texttt{auto\_ptr} 已被C++17移除。它的拷贝语义（拷贝时实际执行移动）违反直觉，容易导致bug。使用 \texttt{unique\_ptr} 替代。

\textbf{要点：} 如果代码中还有 \texttt{auto\_ptr}，应当立即替换为 \texttt{unique\_ptr}。\texttt{unique\_ptr} 的移动语义是显式的（使用 \texttt{std::move}），不会在拷贝时偷偷移动。

\subsection*{\textbf{规则 SL.10:} 使用 \texttt{std::sort} 而非C风格 \texttt{qsort}}

\texttt{std::sort} 是类型安全的，通常比 \texttt{qsort} 更快（因为可以内联比较函数），且支持自定义比较器和迭代器范围。

\textbf{错误示例：}

\begin{lstlisting}[language=C++]
int compare_ints(const void* a, const void* b) {
    return *(int*)a - *(int*)b;  // 可能溢出
}

int arr[] = {5, 3, 1, 4, 2};
qsort(arr, 5, sizeof(int), compare_ints);
\end{lstlisting}

\textbf{正确示例：}

\begin{lstlisting}[language=C++]
vector<int> arr = {5, 3, 1, 4, 2};
sort(arr.begin(), arr.end());  // 类型安全，更快

// 自定义比较
sort(arr.begin(), arr.end(), greater<int>());
\end{lstlisting}

\textbf{要点：} \texttt{std::sort} 使用 introsort（内省排序），时间复杂度保证为 $O(n \log n)$。由于模板内联，它通常比 \texttt{qsort} 快2-3倍。

\section{容器与算法}

\subsection*{选择合适的容器}

标准库提供了多种容器，每种适用于不同的场景：

\begin{itemize}
\item \textbf{\texttt{vector}:} 默认选择。连续内存，缓存友好，随机访问 $O(1)$，尾部插入均摊 $O(1)$。
\item \textbf{\texttt{deque}:} 需要两端高效插入/删除时选择。
\item \textbf{\texttt{list}/\texttt{forward\_list}:} 需要频繁在中间插入/删除且已有迭代器时选择（但通常 \texttt{vector} 仍然更快）。
\item \textbf{\texttt{set}/\texttt{map}:} 需要有序遍历或 $O(\log n)$ 查找时选择。基于红黑树。
\item \textbf{\texttt{unordered\_set}/\texttt{unordered\_map}:} 需要 $O(1)$ 平均查找时选择。基于哈希表。
\item \textbf{\texttt{array}:} 固定大小的数组。零开销，缓存友好。
\end{itemize}

\subsection*{使用标准算法}

标准库提供了大量算法（定义在 \texttt{<algorithm>} 中），应当优先使用它们而非手写循环：

\textbf{错误示例：}

\begin{lstlisting}[language=C++]
// 手写查找最大值
int max_val = v[0];
for (size_t i = 1; i < v.size(); ++i)
    if (v[i] > max_val) max_val = v[i];

// 手写计数
int count = 0;
for (const auto& x : v)
    if (x > threshold) count++;
\end{lstlisting}

\textbf{正确示例：}

\begin{lstlisting}[language=C++]
int max_val = *max_element(v.begin(), v.end());

int count = count_if(v.begin(), v.end(),
    [threshold](int x) { return x > threshold; });
\end{lstlisting}

\textbf{要点：} 标准算法有经过验证的正确性，可以被编译器优化，且通过名称直接表达意图。常用的算法包括：\texttt{sort}、\texttt{find}、\texttt{find\_if}、\texttt{count\_if}、\texttt{transform}、\texttt{for\_each}、\texttt{copy}、\texttt{remove\_if}、\texttt{accumulate}、\texttt{any\_of}、\texttt{all\_of}、\texttt{none\_of} 等。

\section*{练习题}

\begin{enumerate}
\item 以下场景应当选择哪种标准容器？解释你的选择：
\begin{enumerate}
\item 存储一个配置文件中的键值对，需要按名称快速查找。
\item 维护一个按时间排序的事件列表，需要按时间顺序遍历。
\item 实现一个URL缓存，需要 $O(1)$ 的查找性能。
\item 存储一个固定大小的颜色查找表（256个条目）。
\end{enumerate}

\item 使用标准算法重写以下代码：
\begin{lstlisting}[language=C++]
vector<string> names = {"Alice", "Bob", "Charlie", "David"};
vector<string> result;
for (const auto& name : names) {
    if (name.length() > 3) {
        result.push_back(name);
    }
}
sort(result.begin(), result.end());
\end{lstlisting}

\item 解释以下代码中的迭代器失效问题，并给出修复方案：
\begin{lstlisting}[language=C++]
vector<int> v = {1, 2, 3, 4, 5};
for (auto it = v.begin(); it != v.end(); ++it) {
    if (*it % 2 == 0) {
        v.erase(it);
    }
}
\end{lstlisting}

\item 实现一个自定义类型的 \texttt{std::hash} 特化，使其可以用于 \texttt{unordered\_map}：
\begin{lstlisting}[language=C++]
struct Point {
    int x, y;
    bool operator==(const Point& other) const {
        return x == other.x && y == other.y;
    }
};
\end{lstlisting}

\item 比较 \texttt{vector} 和 \texttt{list} 在以下操作中的性能差异，并解释原因：
\begin{enumerate}
\item 顺序遍历100万个元素
\item 在头部插入1000个元素
\item 随机访问第50万个元素
\item 在中间位置插入1000个元素（已有迭代器指向插入位置）
\end{enumerate}
\end{enumerate}


% ============================================================
% 附录
% ============================================================
\appendix

% ------------------------------------------------------------
% 附录 A: 术语表
% ------------------------------------------------------------
\chapter{核心术语表}

本附录收录了本书及 C\texttt{++} Core Guidelines 中的关键术语，供读者查阅。

\begin{longtable}{>{\bfseries}p{4cm} p{10cm}}
\toprule
\textbf{术语} & \textbf{释义} \\
\midrule
\endfirsthead

\toprule
\textbf{术语} & \textbf{释义} \\
\midrule
\endhead

\bottomrule
\endfoot

RAII & Resource Acquisition Is Initialization（资源获取即初始化）。将资源的生命周期绑定到局部对象的生命周期，利用构造/析构机制自动管理资源。\\[6pt]
类型安全 (Type Safety) & 保证每个对象仅按照其类型所允许的方式被访问和操作，避免类型混淆导致的未定义行为。\\[6pt]
边界安全 (Bounds Safety) & 保证对数组、容器等序列的访问不超出其有效范围，防止缓冲区溢出。\\[6pt]
生命周期安全 (Lifetime Safety) & 保证对对象的引用或指针不会在对象被销毁后继续使用，避免悬空引用。\\[6pt]
所有权 (Ownership) & 明确谁负责在何时释放资源。现代 C\texttt{++} 通过智能指针（\texttt{unique\_ptr}、\texttt{shared\_ptr}）在类型系统中表达所有权语义。\\[6pt]
移动语义 (Move Semantics) & 通过右值引用（\texttt{\&\&}）和移动构造/赋值，将资源从即将销毁的对象高效地转移到新对象，避免不必要的深拷贝。\\[6pt]
零规则 (Rule of Zero) & 如果一个类不需要自定义析构函数，则通常也不需要自定义拷贝/移动构造和赋值运算符。优先使用 RAII 成员让编译器自动生成的特殊成员函数完成工作。\\[6pt]
五规则 (Rule of Five) & 如果自定义了析构函数、拷贝构造、拷贝赋值、移动构造、移动赋值中的任何一个，通常需要定义全部五个。\\[6pt]
SFINAE & Substitution Failure Is Not An Error（替换失败不是错误）。模板元编程中的核心机制，允许在模板参数替换失败时优雅地回退而非报错。\\[6pt]
概念 (Concepts) & C\texttt{++}20 引入的语言特性，用于对模板参数施加约束，使模板接口更加清晰、错误信息更加友好。\\[6pt]
不变式 (Invariant) & 在对象的生命周期中始终为真的条件。良好的类设计应确保构造后建立不变式，并在所有公开操作中维护不变式。\\[6pt]
前置条件 (Precondition) & 调用函数之前必须满足的条件。由调用者负责保证。\\[6pt]
后置条件 (Postcondition) & 函数执行完毕后保证成立的条件。由被调用函数负责保证。\\[6pt]
未定义行为 (UB) & Undefined Behavior。C\texttt{++} 标准未规定其行为的操作。一旦发生 UB，程序的任何结果都是不可预测的。\\[6pt]
SL 系列规则 & Standard-library 系列规则，涉及标准库的使用建议。\\[6pt]
ES 系列规则 & Expressions and Statements 系列规则，涉及表达式和语句的编写建议。\\[6pt]
F 系列规则 & Functions 系列规则，涉及函数设计的建议。\\[6pt]
C 系列规则 & Classes and Class Hierarchies 系列规则，涉及类与类层次设计的建议。\\[6pt]
Con 系列规则 & Concurrency 系列规则，涉及并发编程的建议。\\[6pt]
R 系列规则 & Resource Management 系列规则，涉及资源管理的建议。\\[6pt]
E 系列规则 & Error Handling 系列规则，涉及错误处理的建议。\\[6pt]
P 系列规则 & Philosophical rules，涉及编程哲学的高层原则。\\[6pt]
\end{longtable}

% ------------------------------------------------------------
% 附录 B: 推荐阅读
% ------------------------------------------------------------
\chapter{推荐阅读}

以下书籍与资源是深入理解 C\texttt{++} Core Guidelines 及其背后设计理念的重要参考资料。

\section*{核心参考}

\begin{enumerate}[label=\textbf{[\arabic*]}]
  \item Bjarne Stroustrup. \textit{The C\texttt{++} Programming Language, 4th Edition}. Addison-Wesley, 2013.\\
  C\texttt{++} 之父撰写的权威参考，全面覆盖 C\texttt{++}11 标准。

  \item Bjarne Stroustrup. \textit{A Tour of C\texttt{++}, 3rd Edition}. Addison-Wesley, 2022.\\
  简明扼要的现代 C\texttt{++} 导览，适合快速了解 C\texttt{++}17/20 的核心特性。

  \item Scott Meyers. \textit{Effective Modern C\texttt{++}}. O'Reilly Media, 2014.\\
  针对 C\texttt{++}11/14 的 42 条具体建议，与核心指南高度互补。

  \item Herb Sutter, Andrei Alexandrescu. \textit{C\texttt{++} Coding Standards}. Addison-Wesley, 2004.\\
  核心指南的前身之一，101 条编码规范。
\end{enumerate}

\section*{进阶阅读}

\begin{enumerate}[label=\textbf{[\arabic*]}, start=5]
  \item Scott Meyers. \textit{Effective C\texttt{++}, 3rd Edition}. Addison-Wesley, 2005.\\
  经典 C\texttt{++} 最佳实践，虽然基于 C\texttt{++}03，但核心理念仍然适用。

  \item Nicolai M. Josuttis. \textit{The C\texttt{++} Standard Library, 2nd Edition}. Addison-Wesley, 2012.\\
  标准库的百科全书式参考，涵盖 C\texttt{++}11 标准库。

  \item Anthony Williams. \textit{C\texttt{++} Concurrency in Action, 2nd Edition}. Manning, 2019.\\
  并发编程的权威指南，覆盖 C\texttt{++}17 内存模型与线程库。

  \item Marc Gregoire. \textit{Professional C\texttt{++}, 5th Edition}. Wrox, 2021.\\
  全面覆盖 C\texttt{++}20 特性的专业参考书。
\end{enumerate}

\section*{在线资源}

\begin{itemize}
  \item \href{https://isocpp.github.io/CppCoreGuidelines/}{C\texttt{++} Core Guidelines 官方文档}
  \item \href{https://en.cppreference.com/}{cppreference.com --- C\texttt{++} 标准参考}
  \item \href{https://godbolt.org/}{Compiler Explorer (Godbolt) --- 在线编译器}
  \item \href{https://clang.llvm.org/extra/clang-tidy/}{clang-tidy --- 静态分析工具（支持部分核心指南检查）}
\end{itemize}

% ------------------------------------------------------------
% 附录 C: 规则索引
% ------------------------------------------------------------
\chapter{规则索引}

本附录按照规则编号列出所有在本书中讲解的核心指南规则，并标注其所在章节，方便读者快速查找。

\section*{P 系列：哲学规则}

\begin{longtable}{>{\ttfamily}p{2.5cm} p{7cm} p{4cm}}
\toprule
\textbf{规则} & \textbf{概要} & \textbf{章节} \\
\midrule
\endfirsthead

\toprule
\textbf{规则} & \textbf{概要} & \textbf{章节} \\
\midrule
\endhead

\bottomrule
\endfoot

P.1 & 直接表达意图 & 第 2 章 \\
P.2 & 尽量用 ISO C\texttt{++} 编写 & 第 2 章 \\
P.3 & 表达意图 & 第 2 章 \\
P.4 & 理想情况下，程序应该是静态类型安全的 & 第 2 章 \\
P.5 & 编译时能检查的东西不要留到运行时 & 第 2 章 \\
P.6 & 无法在编译时检查的内容应在运行时检查 & 第 2 章 \\
P.7 & 尽早捕获运行时错误 & 第 2 章 \\
P.8 & 不泄漏任何资源 & 第 2 章 \\
P.9 & 不要浪费资源 & 第 2 章 \\
P.10 & 优先使用默认值而非引用传递 & 第 2 章 \\
P.11 & 不在程序中隐藏指针 & 第 2 章 \\
\end{longtable}

\section*{ES 系列：表达式与语句}

\begin{longtable}{>{\ttfamily}p{2.5cm} p{7cm} p{4cm}}
\toprule
\textbf{规则} & \textbf{概要} & \textbf{章节} \\
\midrule
\endfirsthead

\toprule
\textbf{规则} & \textbf{概要} & \textbf{章节} \\
\midrule
\endhead

\bottomrule
\endfoot

ES.1 & 适当时使用抽象，避免使用原始语言特性 & 第 3 章 \\
ES.2 & 适当时使用抽象，避免使用原始语言特性 & 第 3 章 \\
ES.3 & 不要重复自己，避免冗余 & 第 3 章 \\
ES.5 & 避免使用宏 & 第 3 章 \\
ES.10 & 不要在同一表达式中初始化多个变量 & 第 3 章 \\
ES.11 & 用 \texttt{auto} 避免冗余的类型重复 & 第 3 章 \\
ES.20 & 始终初始化对象 & 第 3 章 \\
ES.21 & 不要在声明或初始化之前使用变量 & 第 3 章 \\
ES.22 & 不要在拥有未初始化的变量时声明函数 & 第 3 章 \\
ES.23 & 保留 \texttt{\{\}} 初始化列表语法 & 第 3 章 \\
ES.24 & 使用 \texttt{unique\_ptr} 或 \texttt{shared\_ptr} 管理堆上对象 & 第 3 章 \\
ES.25 & 不要在函数中声明同名变量来隐藏外层变量 & 第 3 章 \\
ES.26 & 不要用一个变量做两件不相关的事 & 第 3 章 \\
ES.27 & 使用 \texttt{std::array} 或 \texttt{stack\_string} 代替 C 风格数组 & 第 3 章 \\
ES.28 & 使用 \texttt{std::array} 或 \texttt{stack\_string} 代替 C 风格数组 & 第 3 章 \\
ES.30 & 不要用宏做程序文本处理 & 第 3 章 \\
ES.31 & 不要用宏做常量或"函数" & 第 3 章 \\
\end{longtable}

\section*{F 系列：函数}

\begin{longtable}{>{\ttfamily}p{2.5cm} p{7cm} p{4cm}}
\toprule
\textbf{规则} & \textbf{概要} & \textbf{章节} \\
\midrule
\endfirsthead

\toprule
\textbf{规则} & \textbf{概要} & \textbf{章节} \\
\midrule
\endhead

\bottomrule
\endfoot

F.1 & 将"打包"操作封装为函数 & 第 4 章 \\
F.2 & 函数应做且仅做一件事 & 第 4 章 \\
F.3 & 保持函数短小精悍 & 第 4 章 \\
F.4 & 如果函数可能需要在编译时求值，则声明为 \texttt{constexpr} & 第 4 章 \\
F.5 & 函数参数少 & 第 4 章 \\
F.6 & 使用 \texttt{noexcept} 标记不抛异常的函数 & 第 4 章 \\
F.7 & 对于通用引用，使用 \texttt{auto\&\&} 参数 & 第 4 章 \\
F.10 & 如果一个操作可以重用，应将其命名为函数 & 第 4 章 \\
F.15 & 遵循简单而惯用的参数传递方式 & 第 4 章 \\
F.16 & 对于"输入"参数，按值传递或按 \texttt{const} 引用传递 & 第 4 章 \\
F.17 & 对于"输入-输出"参数，按非 \texttt{const} 引用传递 & 第 4 章 \\
F.18 & 对于"将被移动"的参数，按值传递 & 第 4 章 \\
F.19 & 对于"转发"参数，使用 \texttt{auto\&\&} 并 \texttt{std::forward} & 第 4 章 \\
F.20 & 对于"输出"输出值，优先使用返回值 & 第 4 章 \\
F.21 & 需要返回多个值时，返回结构体 & 第 4 章 \\
F.36 & 使用 \texttt{[[nodiscard]]} 标记不应忽略返回值的函数 & 第 4 章 \\
F.38 & 使用返回值参数代替函数返回 void & 第 4 章 \\
F.40 & 使用 \texttt{final} 阻止覆写 & 第 4 章 \\
F.41 & 返回元组以返回多个不相关的值 & 第 4 章 \\
\end{longtable}

\section*{C 系列：类与类层次}

\begin{longtable}{>{\ttfamily}p{2.5cm} p{7cm} p{4cm}}
\toprule
\textbf{规则} & \textbf{概要} & \textbf{章节} \\
\midrule
\endfirsthead

\toprule
\textbf{规则} & \textbf{概要} & \textbf{章节} \\
\midrule
\endhead

\bottomrule
\endfoot

C.1 & 将相关的变量组织为结构体 & 第 5 章 \\
C.2 & 如果类有不变式，则使用构造函数建立不变式 & 第 5 章 \\
C.3 & 表示默认操作 & 第 5 章 \\
C.4 & 使默认操作在语义上正确 & 第 5 章 \\
C.5 & 将辅助函数放在与其辅助的类相同的命名空间中 & 第 5 章 \\
C.7 & 不要定义类或函数；声明它们 & 第 5 章 \\
C.8 & 用 \texttt{class} 而非 \texttt{struct} 表示有私有成员的类型 & 第 5 章 \\
C.9 & 最小化暴露的成员 & 第 5 章 \\
C.10 & 优先使用具体类型而非类层次 & 第 5 章 \\
C.11 & 使具体类型 \texttt{constexpr} 可初始化 & 第 5 章 \\
C.12 & 不要在拷贝赋值运算符中使用 \texttt{const} 参数 & 第 5 章 \\
C.20 & 如果能避免定义默认操作，就避免 & 第 5 章 \\
C.21 & 如果定义或禁用默认操作，对所有都这样做 & 第 5 章 \\
C.22 & 使默认操作一致 & 第 5 章 \\
C.30 & 如果析构函数需要执行显式操作，定义析构函数 & 第 5 章 \\
C.31 & 类分配的所有资源必须由析构函数释放 & 第 5 章 \\
C.32 & 如果类拥有原始指针或指针，考虑是否应该拥有 & 第 5 章 \\
C.33 & 如果类拥有带所有权的指针，定义拷贝操作 & 第 5 章 \\
C.35 & 基类需要有虚析构函数 & 第 5 章 \\
C.36 & 基类不应有数据成员 & 第 5 章 \\
C.37 & 使基类具有 \texttt{noexcept} 的析构函数 & 第 5 章 \\
C.40 & 如果类有虚函数，定义析构函数 & 第 5 章 \\
C.41 & 构造函数应创建完全初始化的对象 & 第 5 章 \\
C.42 & 构造函数不应抛出异常 & 第 5 章 \\
C.43 & 确保可复制的类有默认构造函数 & 第 5 章 \\
C.44 & 使默认构造函数简单且非抛出 & 第 5 章 \\
C.45 & 不要定义仅初始化数据成员的构造函数 & 第 5 章 \\
C.46 & 默认构造函数至少初始化所有成员 & 第 5 章 \\
C.47 & 按照成员声明的顺序定义和初始化成员变量 & 第 5 章 \\
C.48 & 对常量默认值使用就地初始化 & 第 5 章 \\
C.49 & 优先用初始化而非赋值来初始化成员 & 第 5 章 \\
C.50 & 使用"工厂函数"处理需要多态或复杂初始化的情形 & 第 5 章 \\
C.65 & 使移动操作 \texttt{noexcept} & 第 5 章 \\
C.67 & 多态基类应删除或保护拷贝操作 & 第 5 章 \\
C.80 & 用 \texttt{delete} 禁止默认行为 & 第 5 章 \\
C.81 & 用 \texttt{=default} 使用编译器生成的默认行为 & 第 5 章 \\
C.90 & 使拷贝赋值和拷贝构造 \texttt{noexcept} & 第 5 章 \\
C.120 & 使用类层次结构表示抽象概念（运行时多态） & 第 6 章 \\
C.122 & 用抽象类做接口 & 第 6 章 \\
C.126 & 用虚函数调用分派 & 第 6 章 \\
C.127 & 有虚函数的类应有虚析构函数 & 第 6 章 \\
C.128 & 虚函数应指定 \texttt{default}、\texttt{delete}、\texttt{override} 或 \texttt{final} & 第 6 章 \\
C.129 & 设计类层次时考虑抽象与实现 & 第 6 章 \\
C.130 & 对大型多态类使用虚函数 & 第 6 章 \\
C.131 & 避免琐碎的 getter 和 setter & 第 6 章 \\
C.133 & 避免 \texttt{protected} 数据 & 第 6 章 \\
C.134 & 用 \texttt{public} 成员保护非虚基类数据 & 第 6 章 \\
C.135 & 用多重继承表示独立的方面 & 第 6 章 \\
C.136 & 用 \texttt{dynamic\_cast} 安全地导航类层次 & 第 6 章 \\
C.137 & 用 \texttt{dynamic\_cast} 安全地导航类层次 & 第 6 章 \\
C.138 & 用 \texttt{dynamic\_cast} 代替类型转换 & 第 6 章 \\
C.139 & 使用 \texttt{final} 于虚函数 & 第 6 章 \\
C.140 & 不要为类层次使用不相关的转换运算符 & 第 6 章 \\
C.145 & 多态地访问对象 & 第 6 章 \\
C.146 & 使用 \texttt{dynamic\_cast} 访问类层次中的派生类对象 & 第 6 章 \\
C.147 & 对已初始化的引用使用 \texttt{dynamic\_cast} & 第 6 章 \\
C.149 & 用 \texttt{unique\_ptr} 或 \texttt{shared\_ptr} 避免忘记 \texttt{delete} & 第 6 章 \\
C.150 & 用 \texttt{unique\_ptr} 或 \texttt{shared\_ptr} 避免忘记 \texttt{delete} & 第 6 章 \\
C.151 & 用 \texttt{unique\_ptr} 或 \texttt{shared\_ptr} 避免忘记 \texttt{delete} & 第 6 章 \\
C.152 & 禁止赋值给基类子对象 & 第 6 章 \\
\end{longtable}

\section*{Con 系列：并发}

\begin{longtable}{>{\ttfamily}p{2.5cm} p{7cm} p{4cm}}
\toprule
\textbf{规则} & \textbf{概要} & \textbf{章节} \\
\midrule
\endfirsthead

\toprule
\textbf{规则} & \textbf{概要} & \textbf{章节} \\
\midrule
\endhead

\bottomrule
\endfoot

Con.1 & 默认使用线程而非进程 & 第 7 章 \\
Con.2 & 默认使用标准库的并发工具 & 第 7 章 \\
Con.3 & 不要对共享数据加锁，除非必须 & 第 7 章 \\
Con.4 & 用 \texttt{mutex} 或 \texttt{atomic} 保护共享数据 & 第 7 章 \\
Con.5 & 形成层次结构来避免死锁 & 第 7 章 \\
Con.6 & 用 \texttt{lock\_guard} 或 \texttt{scoped\_lock} 管理互斥量 & 第 7 章 \\
Con.7 & 避免使用裸 \texttt{atomic} 变量 & 第 7 章 \\
Con.8 & 不要用锁来保护个别变量 & 第 7 章 \\
Con.9 & 限制使用 \texttt{yield} & 第 7 章 \\
Con.10 & 优先使用标准库的并发设施 & 第 7 章 \\
Con.11 & 不要使用无界队列 & 第 7 章 \\
Con.12 & 使用原子操作时确保内存顺序正确 & 第 7 章 \\
Con.32 & 让共享的可访问对象用 \texttt{shared\_ptr} 引用 & 第 7 章 \\
Con.34 & 等待线程完成 & 第 7 章 \\
Con.35 & 使用 \texttt{jthread} 代替 \texttt{thread} & 第 7 章 \\
Con.36 & 绑定不可分离的线程 & 第 7 章 \\
Con.40 & 确保锁的作用域尽可能小 & 第 7 章 \\
Con.41 & 不要用 \texttt{volatile} 做线程间通信 & 第 7 章 \\
Con.42 & 不要 \texttt{sleep} 等待条件变化 & 第 7 章 \\
Con.43 & 不要对共享状态使用无锁编程 & 第 7 章 \\
Con.44 & 使用 \texttt{std::async} 而非裸线程 & 第 7 章 \\
Con.45 & 不要生成和吞噬异常 & 第 7 章 \\
Con.46 & 不要使用信号量做线程间通信 & 第 7 章 \\
\end{longtable}

\section*{R 系列：资源管理}

\begin{longtable}{>{\ttfamily}p{2.5cm} p{7cm} p{4cm}}
\toprule
\textbf{规则} & \textbf{概要} & \textbf{章节} \\
\midrule
\endfirsthead

\toprule
\textbf{规则} & \textbf{概要} & \textbf{章节} \\
\midrule
\endhead

\bottomrule
\endfoot

R.1 & 资源获取即初始化 (RAII) & 第 9 章 \\
R.2 & 在非全局函数中尽量少用裸 \texttt{new} & 第 9 章 \\
R.3 & 裸 \texttt{new} 表达式是代码异味 & 第 9 章 \\
R.4 & 裸 \texttt{delete} 表达式是代码异味 & 第 9 章 \\
R.5 & 优先使用作用域对象 & 第 9 章 \\
R.6 & 避免全局变量 & 第 9 章 \\
R.7 & 避免非 const 全局引用 & 第 9 章 \\
R.8 & 避免全局指针 & 第 9 章 \\
R.10 & 避免 \texttt{malloc()} 和 \texttt{free()} & 第 9 章 \\
R.11 & 避免显式调用 \texttt{new} 和 \texttt{delete} & 第 9 章 \\
R.12 & 避免多个裸的拥有指针 & 第 9 章 \\
R.13 & 立即执行至少一个拥有操作 & 第 9 章 \\
R.14 & 避免 \texttt{new[]} 和 \texttt{delete[]} 参数 & 第 9 章 \\
R.15 & 始终重载 \texttt{new} 和 \texttt{delete} & 第 9 章 \\
R.20 & 使用 \texttt{unique\_ptr} 或 \texttt{shared\_ptr} 避免忘记 \texttt{delete} & 第 9 章 \\
R.21 & 优先使用 \texttt{unique\_ptr} 而非 \texttt{shared\_ptr} & 第 9 章 \\
R.22 & 用 \texttt{make\_unique} 创建 \texttt{unique\_ptr} & 第 9 章 \\
R.23 & 用 \texttt{make\_shared} 创建 \texttt{shared\_ptr} & 第 9 章 \\
R.24 & 用 \texttt{std::atomic\_shared\_ptr<T>} 共享 \texttt{shared\_ptr} & 第 9 章 \\
R.30 & 用 \texttt{unique\_ptr} 或 \texttt{shared\_ptr} 传递所有权 & 第 9 章 \\
R.31 & 如果 \texttt{unique\_ptr} 模式不够，使用 \texttt{shared\_ptr} & 第 9 章 \\
R.32 & 取 \texttt{unique\_ptr} 参数表示获取所有权 & 第 9 章 \\
R.33 & 取 \texttt{const unique\_ptr\&} 参数表示保留所有权 & 第 9 章 \\
R.34 & 取 \texttt{const shared\_ptr\&} 参数表示不改变所有权 & 第 9 章 \\
R.35 & 取 \texttt{shared\_ptr\&} 参数表示重新绑定 & 第 9 章 \\
R.36 & 取 \texttt{shared\_ptr\&\&} 参数表示获取所有权 & 第 9 章 \\
R.37 & 不要将 \texttt{shared\_ptr} 传递给非拥有者 & 第 9 章 \\
\end{longtable}

\section*{E 系列：错误处理}

\begin{longtable}{>{\ttfamily}p{2.5cm} p{7cm} p{4cm}}
\toprule
\textbf{规则} & \textbf{概要} & \textbf{章节} \\
\midrule
\endfirsthead

\toprule
\textbf{规则} & \textbf{概要} & \textbf{章节} \\
\midrule
\endhead

\bottomrule
\endfoot

E.1 & 在开发阶段尽早系统地找出错误 & 第 10 章 \\
E.2 & 用异常表示无法完成的任务 & 第 10 章 \\
E.3 & 仅用异常来报告无法在当前作用域处理的错误 & 第 10 章 \\
E.4 & 围绕每个可能失败的操作设计错误处理策略 & 第 10 章 \\
E.5 & 让构造函数建立类不变式；若无法建立则抛出异常 & 第 10 章 \\
E.6 & 使用 RAII 防止资源泄漏 & 第 10 章 \\
E.7 & 注意前置条件 & 第 10 章 \\
E.12 & 不要在 \texttt{noexcept} 函数中抛出异常 & 第 10 章 \\
E.13 & 不要在直接拥有对象的函数中抛出异常 & 第 10 章 \\
E.14 & 使用自定义异常类型而非 \texttt{int} 或 \texttt{string} & 第 10 章 \\
E.15 & 通过合适的异常类型抛出异常 & 第 10 章 \\
E.16 & 用 \texttt{catch} 处理异常 & 第 10 章 \\
E.17 & 不要在每个函数中捕获所有异常 & 第 10 章 \\
E.18 & 最小化显式 \texttt{catch} 的使用 & 第 10 章 \\
E.19 & 用 \texttt{final\_action} 在函数退出时清理 & 第 10 章 \\
E.25 & 如果不能抛出异常，模拟 RAII 进行清理 & 第 10 章 \\
E.26 & 不要在无法 \texttt{abort} 的情况下 \texttt{throw} & 第 10 章 \\
E.27 & 使用 \texttt{std::terminate} 处理不可恢复的错误 & 第 10 章 \\
E.28 & 避免基于错误码的错误处理 & 第 10 章 \\
E.30 & 不要用异常传播做控制流 & 第 10 章 \\
E.31 & 正确区分异常和错误码 & 第 10 章 \\
\end{longtable}

\section*{SL 系列：标准库}

\begin{longtable}{>{\ttfamily}p{2.5cm} p{7cm} p{4cm}}
\toprule
\textbf{规则} & \textbf{概要} & \textbf{章节} \\
\midrule
\endfirsthead

\toprule
\textbf{规则} & \textbf{概要} & \textbf{章节} \\
\midrule
\endhead

\bottomrule
\endfoot

SL.1 & 使用标准库 & 第 8 章 \\
SL.2 & 优先使用标准库而非其他库 & 第 8 章 \\
SL.3 & 不要给标准库添加实体 & 第 8 章 \\
SL.4 & 安全地使用标准库 & 第 8 章 \\
SL.5 & 避免类型不安全 & 第 8 章 \\
SL.6 & 使用 ISO C\texttt{++} 标准库特性 & 第 8 章 \\
SL.7 & 不要将 \texttt{std::string} 与 \texttt{nullptr} 混用 & 第 8 章 \\
SL.8 & 避免 \texttt{std::string} 的 \texttt{c\_str()} & 第 8 章 \\
SL.9 & 避免使用窄化转换 & 第 8 章 \\
SL.10 & 使用 \texttt{std::string} 或 \texttt{string\_view} 代替 C 风格字符串 & 第 8 章 \\
SL.11 & 使用 \texttt{gsl::span} 代替指针加长度 & 第 8 章 \\
SL.12 & 使用 \texttt{gsl::narrow\_cast} 或 \texttt{narrow} 做窄化转换 & 第 8 章 \\
SL.13 & 使用 \texttt{string\_view} 或 \texttt{span} 避免越界访问 & 第 8 章 \\
SL.14 & 使用 \texttt{vector} 代替 C 风格数组 & 第 8 章 \\
SL.15 & 避免使用裸指针 & 第 8 章 \\
SL.20 & 使用 \texttt{string\_view} 或 \texttt{span} 代替指针加长度 & 第 8 章 \\
SL.21 & 不要用裸指针做参数 & 第 8 章 \\
SL.22 & 不要用裸指针做变量 & 第 8 章 \\
SL.23 & 避免抛出异常的临时变量 & 第 8 章 \\
SL.24 & 用 \texttt{unique\_ptr} 代替裸 \texttt{new} & 第 8 章 \\
SL.25 & 用 \texttt{make\_unique} 或 \texttt{make\_shared} & 第 8 章 \\
SL.26 & 用 \texttt{unique\_ptr} 或 \texttt{shared\_ptr} 代替裸 \texttt{new} & 第 8 章 \\
SL.27 & 用 \texttt{vector} 或 \texttt{string\_view} 代替 C 风格数组 & 第 8 章 \\
\end{longtable}

% ------------------------------------------------------------
% 后记
% ------------------------------------------------------------
\backmatter

\chapter{后记}

C\texttt{++} Core Guidelines 是一项持续演进的社区工程。本书所解读的规则和建议，也会随着语言标准的更新和社区的实践而不断修订。我们鼓励读者定期访问指南的官方仓库，关注最新的规则变更和新增内容。

编程规范的价值不在于被记住多少条，而在于能否在日常开发中自觉地运用。希望本书能帮助读者建立起对现代 C\texttt{++} 编程范式的深刻理解，在实践中写出更安全、更高效、更优雅的代码。

\printindex

\end{document}
